% Options for packages loaded elsewhere
\PassOptionsToPackage{unicode}{hyperref}
\PassOptionsToPackage{hyphens}{url}
\documentclass[
]{article}
\usepackage{xcolor}
\usepackage{amsmath,amssymb}
\setcounter{secnumdepth}{-\maxdimen} % remove section numbering
\usepackage{iftex}
\ifPDFTeX
  \usepackage[T1]{fontenc}
  \usepackage[utf8]{inputenc}
  \usepackage{textcomp} % provide euro and other symbols
\else % if luatex or xetex
  \usepackage{unicode-math} % this also loads fontspec
  \defaultfontfeatures{Scale=MatchLowercase}
  \defaultfontfeatures[\rmfamily]{Ligatures=TeX,Scale=1}
\fi
\usepackage{lmodern}
\ifPDFTeX\else
  % xetex/luatex font selection
\fi
% Use upquote if available, for straight quotes in verbatim environments
\IfFileExists{upquote.sty}{\usepackage{upquote}}{}
\IfFileExists{microtype.sty}{% use microtype if available
  \usepackage[]{microtype}
  \UseMicrotypeSet[protrusion]{basicmath} % disable protrusion for tt fonts
}{}
\makeatletter
\@ifundefined{KOMAClassName}{% if non-KOMA class
  \IfFileExists{parskip.sty}{%
    \usepackage{parskip}
  }{% else
    \setlength{\parindent}{0pt}
    \setlength{\parskip}{6pt plus 2pt minus 1pt}}
}{% if KOMA class
  \KOMAoptions{parskip=half}}
\makeatother
\usepackage{longtable,booktabs,array}
\usepackage{calc} % for calculating minipage widths
% Correct order of tables after \paragraph or \subparagraph
\usepackage{etoolbox}
\makeatletter
\patchcmd\longtable{\par}{\if@noskipsec\mbox{}\fi\par}{}{}
\makeatother
% Allow footnotes in longtable head/foot
\IfFileExists{footnotehyper.sty}{\usepackage{footnotehyper}}{\usepackage{footnote}}
\makesavenoteenv{longtable}
\setlength{\emergencystretch}{3em} % prevent overfull lines
\providecommand{\tightlist}{%
  \setlength{\itemsep}{0pt}\setlength{\parskip}{0pt}}
setlength\{emergencystretch\}\{3em\}
sloppy
\usepackage{bookmark}
\IfFileExists{xurl.sty}{\usepackage{xurl}}{} % add URL line breaks if available
\urlstyle{same}
\hypersetup{
  hidelinks,
  pdfcreator={LaTeX via pandoc}}

\author{}
\date{}

\begin{document}

\section{Module 2 - Vertical Axis and Column
Assembly}\label{module-2---vertical-axis-and-column-assembly}

\subsection{1. Z-Axis Design Philosophy}\label{z-axis-design-philosophy}

The vertical axis (Z-axis) presents unique engineering challenges
fundamentally different from horizontal motion axes. While horizontal
axes operate in a relatively neutral gravitational field and benefit
from distributed loads, the vertical axis must continuously resist
gravity, manage large cantilever moments, and overcome the dynamic
challenges of moving masses in opposition to gravitational acceleration.
The design philosophy for vertical axis systems integrates structural
mechanics, dynamic stability theory, mass-energy management, and thermal
compensation into a coherent framework that delivers precision,
responsiveness, and reliability.

\subsubsection{1.1 Mass \& Inertia Management: The Foundation of Dynamic
Performance}\label{mass-inertia-management-the-foundation-of-dynamic-performance}

\paragraph{Fundamental Challenge}\label{fundamental-challenge}

Every kilogram of mass moving on the Z-axis requires continuous energy
input to resist gravity and demands proportionally higher motor torque
for acceleration. The moment of inertia about the drive screw axis
scales with the square of the radial distance, making compact designs
critical.

\paragraph{Design Principles}\label{design-principles}

\begin{enumerate}
\def\labelenumi{\arabic{enumi}.}
\item
  Minimum Mass Design:

  \begin{itemize}
  \tightlist
  \item
    Use topology optimization to remove non-structural material
  \item
    Employ high-strength aluminum alloys (7075-T6, yield
    \textasciitilde500 MPa) instead of steel where possible
  \item
    Consider magnesium alloys (AZ91D) for ultra-lightweight applications
    (density 1.8 g/cm\^{}3 vs 2.7 g/cm\^{}3 for Al)
  \item
    Use hollow-core composite structures (carbon fiber/epoxy) for
    ultimate mass reduction
  \end{itemize}
\item
  Gravitational Energy Management: The potential energy change during
  vertical motion: {[}\Delta E\_p = m g \Delta z{]}

  For a 5 kg spindle assembly moving 200 mm: {[}\Delta E\_p = 5
  \times 9.81 \times 0.2 = 9.81 \text{ J}{]}

  This energy must be supplied (upward motion) or dissipated (downward
  motion) by the drive system.
\item
  Counterbalance Strategy: A properly implemented counterbalance system:

  \begin{itemize}
  \tightlist
  \item
    Reduces motor RMS torque by \textasciitilde50-70\%
  \item
    Eliminates gravitational bias in the servo control loop
  \item
    Enables symmetric acceleration profiles
  \item
    Reduces thermal load on motor and amplifier
  \item
    Extends bearing and drive component life
  \end{itemize}
\end{enumerate}

\paragraph{Counterbalance Force
Calculation}\label{counterbalance-force-calculation}

For a moving mass (m) with center of gravity at distance (L\_\{cg\})
from the column mounting plane:

{[}F\_\{cb\} = m g \cos(\theta){]}

Where (\theta) is the machine tilt angle (typically 0° for vertical
machines).

For adjustable counterbalance using gas springs: {[}F\_\{gas\}(x) = F\_0
+ k\_\{gas\}(x - x\_0){]}

Where: - (F\_0) = initial preload force - (k\_\{gas\}) = gas spring
stiffness (typically 0.5-2.0 N/mm) - (x) = displacement from neutral
position - (x\_0) = initial compressed length

\paragraph{Design Example}\label{design-example}

Moving mass: 8 kg (spindle + mounting plate + carriage) Required
counterbalance force: (F = 8 \times 9.81 = 78.5) N

Using two gas springs in parallel: - Individual spring force: 39.3 N -
Select standard 400 N gas spring compressed to \textasciitilde10\%
extension - Mounting geometry ensures force vector alignment within
+/-5° over full travel

\begin{enumerate}
\def\labelenumi{\arabic{enumi}.}
\setcounter{enumi}{3}
\item
  Moment of Inertia Minimization: For a rotating drive screw, the
  reflected inertia is critical:

  {[}J\_\{total\} = J\_\{motor\} + J\_\{coupling\} + J\_\{screw\} +
  \frac{m_{carriage}}{(2\pi/p)^2}{]}

  Where (p) is the screw pitch (m/revolution).

  For a 5 mm pitch screw: {[}J\_\{reflected\} = \frac{m}{(2\pi/0.005)^2}
  = m \times 6.33 \times 10\^{}\{-7\}{]}

  An 8 kg carriage reflects as: (J = 5.06 \times 10\^{}\{-6\}) kg*m\^{}2

  This is typically 10-50x the motor rotor inertia, dominating system
  dynamics.
\end{enumerate}

\subsubsection{1.2 Column Stiffness: Structural Optimization for
Precision}\label{column-stiffness-structural-optimization-for-precision}

Engineering Objective: The column structure must provide a rigid
reference frame that maintains rail parallelism and minimizes deflection
under cutting forces. Unlike horizontal gantry beams that can be
designed as simply-supported structures, vertical columns typically
function as cantilevers with concentrated loads at the free end.

Cantilever Mechanics:

For a cantilever beam with length (L), Young's modulus (E), and second
moment of area (I), subjected to tip load (F):

{[}\delta\_\{tip\} = \frac{F L^3}{3 E I}{]}

Critical Insight: Deflection scales with the \emph{cube} of length and
inversely with moment of area. Doubling column height increases
deflection by 8x, while doubling cross-sectional depth increases
stiffness by \textasciitilde8x (for rectangular sections).

Design Strategy:

\begin{enumerate}
\def\labelenumi{\arabic{enumi}.}
\item
  Maximize Second Moment of Area: For rectangular hollow section
  (external dimensions (b x h), wall thickness (t)):

  {[}I\_x = \frac{b h^3}{12} - \frac{(b-2t)(h-2t)^3}{12}{]}

  For a 150x150x8 mm square tube: {[}I\_x = \frac{150 \times 150^3}{12}
  - \frac{134 \times 134^3}{12} = 42.19 \times 10\^{}\{6\} - 26.98
  \times 10\^{}\{6\}{]} {[}I\_x = 15.21 \times 10\^{}\{6\}
  \text{ mm}\^{}4 = 1.52 \times 10\^{}\{-5\} \text{ m}\^{}4{]}
\item
  Material Selection: Common materials for column structures:

  {\def\LTcaptype{} % do not increment counter
  \begin{longtable}[]{@{}lllll@{}}
  \toprule\noalign{}
  Material & E (GPa) & Density (kg/m\^{}3) & \$E/\rho\$ & Cost
  multiplier \\
  \midrule\noalign{}
  \endhead
  \bottomrule\noalign{}
  \endlastfoot
  Mild steel & 200 & 7850 & 25.5 & 1.0x \\
  Cast iron & 120 & 7200 & 16.7 & 0.8x \\
  Aluminum 6061 & 69 & 2700 & 25.6 & 3.5x \\
  Steel + ribs & 200 & 7850 & 25.5 & 1.4x \\
  \end{longtable}
  }

  Key Insight: While aluminum has lower absolute stiffness, its specific
  stiffness (E/rho) matches steel. For vertical axes where mass
  reduction is critical, ribbed aluminum columns offer excellent
  performance.
\item
  Internal Ribbing and Stiffening: Strategic placement of internal ribs
  increases (I) without proportional mass increase:

  \begin{itemize}
  \tightlist
  \item
    Longitudinal ribs at extreme fibers (corners)
  \item
    Transverse bulkheads every 150-250 mm
  \item
    Diagonal bracing in high-stress regions
  \end{itemize}

  FEA-optimized designs achieve 30-50\% stiffness increase with only
  15-20\% mass penalty.
\item
  Deflection Specification: Industry practice for precision machines:

  {[}\delta\_\{max\} = \frac{L}{10,000} \text{ to } \frac{L}{20,000}{]}

  For 500 mm cantilever: {[}\delta\_\{max\} = 0.025 \text{ to } 0.05
  \text{ mm}{]}
\end{enumerate}

Worked Example: Column Sizing

Requirements: - Cantilever length: (L = 400) mm - Maximum cutting force:
(F = 500) N (conservative) - Material: Steel ((E = 200) GPa) - Allowable
deflection: (\delta\_\{max\} = 0.03) mm

Solution:

Required moment of inertia: {[}I\_\{req\} =
\frac{F L^3}{3 E \delta_{max}} =
\frac{500 \times 0.4^3}{3 \times 200 \times 10^9 \times 0.00003}{]}

{[}I\_\{req\} = \frac{32}{18 \times 10^6} = 1.78 \times 10\^{}\{-6\}
\text{ m}\^{}4 = 1.78 \times 10\^{}6 \text{ mm}\^{}4{]}

For square hollow section, approximating (I \approx \frac{b^3 t}{3})
(thin-wall): {[}b\^{}3 t \approx 5.33 \times 10\^{}6{]}

Trying (b = 120) mm, (t = 8) mm: {[}120\^{}3 \times 8 = 13.8
\times 10\^{}6 \text{ mm}\^{}4{]} (adequate)

Selected section: 120x120x8 mm steel RHS - Mass per meter: 28.3 kg/m -
Actual (I\_x = 3.52 \times 10\^{}6) mm\^{}4 (provides 2x safety factor)

\subsubsection{1.3 Symmetry \& Thermal Behaviour: Compensating
Environmental
Effects}\label{symmetry-thermal-behaviour-compensating-environmental-effects}

Thermal Expansion Challenge: Vertical columns experience thermal
gradients from: 1. Motor heat dissipation (30-100 W continuous) 2.
Ambient temperature variations (+/-5°C typical) 3. Process heat
(cutting/welding operations) 4. Solar loading (through enclosure panels)

Thermal Expansion Coefficient: For steel: (\alpha = 11.7
\times 10\^{}\{-6\}) /°C

A 500 mm column experiencing 10°C differential: {[}\Delta L = \alpha L
\Delta T = 11.7 \times 10\^{}\{-6\} \times 500 \times 10 = 0.0585
\text{ mm}{]}

This 58 mum error exceeds precision requirements!

Design Solutions:

\begin{enumerate}
\def\labelenumi{\arabic{enumi}.}
\item
  Symmetric Cross-Sections: Using square or circular hollow sections
  ensures uniform thermal expansion in all radial directions. Asymmetric
  sections (C-channel, I-beam) create differential expansion leading to
  bending.
\item
  Thermal Symmetry:

  \begin{itemize}
  \tightlist
  \item
    Mount motors on centerline or use paired motors
  \item
    Insulate motor from column structure
  \item
    Use thermal breaks at mounting interfaces
  \item
    Employ forced air circulation within hollow column
  \end{itemize}
\item
  Material Selection: Materials with low thermal expansion coefficients:

  {\def\LTcaptype{} % do not increment counter
  \begin{longtable}[]{@{}lll@{}}
  \toprule\noalign{}
  Material & alpha (10\textsuperscript{-}6/°C) & Relative expansion \\
  \midrule\noalign{}
  \endhead
  \bottomrule\noalign{}
  \endlastfoot
  Steel & 11.7 & 1.00x \\
  Aluminum & 23.6 & 2.02x \\
  Cast iron & 10.8 & 0.92x \\
  Invar 36 & 1.3 & 0.11x \\
  Carbon fiber & 0.5 to -1.0 & 0.04x \\
  \end{longtable}
  }

  Practical Approach: Use steel or cast iron for primary structure, with
  carbon fiber composite in extreme precision applications.
\item
  Active Compensation: Modern CNC controllers implement thermal
  compensation:

  {[}Z\_\{corrected\} = Z\_\{commanded\} + K\_\{thermal\}(T -
  T\_\{ref\}){]}

  Where (K\_\{thermal\}) is determined through calibration cycles.
\end{enumerate}

\subsubsection{1.4 Vibration \& Resonance: Dynamic Stability
Requirements}\label{vibration-resonance-dynamic-stability-requirements}

Fundamental Principle: The natural frequency of the Z-axis structure
must be significantly higher than the servo control bandwidth to prevent
control-structure interaction and ensure stable operation.

Natural Frequency Calculation:

For cantilever column modeled as spring-mass system: {[}f\_n =
\frac{1}{2\pi}\sqrt{\frac{k}{m}}{]}

Where the spring constant for cantilever: {[}k = \frac{3EI}{L^3}{]}

Combined expression: {[}f\_n = \frac{1}{2\pi}\sqrt{\frac{3EI}{mL^3}}{]}

Design Rule: {[}f\_n \geq 5 \times f\_\{servo\} \text{ (minimum)}{]}
{[}f\_n \geq 10 \times f\_\{servo\} \text{ (preferred)}{]}

For typical servo bandwidth of 30 Hz: {[}f\_n \geq 150
\text{ Hz (minimum)}{]} {[}f\_n \geq 300 \text{ Hz (preferred)}{]}

Worked Example:

Given: - Column: 120x120x8 mm steel RHS - Length: (L = 400) mm - Moving
mass: (m = 8) kg - (E = 200) GPa - (I = 3.52 \times 10\^{}\{-6\}) m\^{}4

Calculate natural frequency:

{[}k = \frac{3 \times 200 \times 10^9 \times 3.52 \times 10^{-6}}{0.4^3}
= \frac{2.112 \times 10^6}{0.064} = 33.0 \times 10\^{}6 \text{ N/m}{]}

{[}f\_n = \frac{1}{2\pi}\sqrt{\frac{33.0 \times 10^6}{8}} =
\frac{1}{6.283}\sqrt{4.125 \times 10^6}{]}

{[}f\_n = \frac{2031}{6.283} = 323 \text{ Hz}{]}

Result: This design achieves (f\_n = 323) Hz, providing 10.8x margin
over 30 Hz servo bandwidth (excellent).

Resonance Mitigation Strategies:

\begin{enumerate}
\def\labelenumi{\arabic{enumi}.}
\item
  \textbf{Structural Damping:}

  \begin{itemize}
  \tightlist
  \item
    Use cast iron (damping ratio ξ approx 0.02-0.05) instead of steel (ξ
    approx 0.001-0.002)
  \item
    Add constrained-layer damping treatments
  \item
    Use polymer-composite hybrid structures
  \end{itemize}
\item
  \textbf{Modal Analysis:} Perform FEA to identify all modes below 500
  Hz:

  \begin{itemize}
  \tightlist
  \item
    1st mode: Typically Z-axis bending
  \item
    2nd mode: Torsion about column axis
  \item
    3rd mode: X-Y rocking of carriage
  \end{itemize}

  Ensure all modes satisfy (f\_i \textgreater{} 5 f\_\{servo\})
\item
  \textbf{Active Damping:} Modern servo drives implement notch filters
  and low-pass filters to suppress resonances:

  Notch filter transfer function: {[}H(s) =
  \frac{s^2 + 2\zeta_z \omega_n s + \omega_n^2}{s^2 + 2\zeta_p \omega_n s + \omega_n^2}{]}

  Tuned to structural resonant frequency with high attenuation.
\end{enumerate}

\subsubsection{1.5 Serviceability: Design for Maintenance and
Adjustment}\label{serviceability-design-for-maintenance-and-adjustment}

\textbf{Critical Insight:} The most precisely designed vertical axis
will degrade over time without proper maintenance access. Bearing
preload adjustment, rail replacement, and alignment verification must be
possible without complete machine disassembly.

\textbf{Design Requirements:}

\begin{enumerate}
\def\labelenumi{\arabic{enumi}.}
\tightlist
\item
  \textbf{Bearing Access:}

  \begin{itemize}
  \tightlist
  \item
    Removable covers at bearing mounting locations
  \item
    Jack-screw provisions for preload adjustment
  \item
    Clearance for bearing puller tools
  \item
    Documented preload values (typically 4-8\% of dynamic load rating)
  \end{itemize}
\item
  \textbf{Rail Replacement:}

  \begin{itemize}
  \tightlist
  \item
    Rails should be removable without disturbing column structure
  \item
    Use dowel pins for repeatable rail positioning
  \item
    Provide access for precision height gauge measurements
  \item
    Design for single-rail replacement (if paired rails used)
  \end{itemize}
\item
  \textbf{Screw Alignment:}

  \begin{itemize}
  \tightlist
  \item
    Shim packs at bearing mounts
  \item
    Slotted mounting holes (+/-1.0 mm adjustment)
  \item
    Dial indicator access points
  \item
    Alignment specification: \textless=0.02 mm TIR over full travel
  \end{itemize}
\item
  \textbf{Counterbalance Adjustment:}

  \begin{itemize}
  \tightlist
  \item
    External force adjustment (gas spring pressure or weight stack)
  \item
    Test procedure: Measure motor current at mid-stroke with zero
    external load
  \item
    Target: \textless=10\% variation from ideal balance current
  \end{itemize}
\item
  \textbf{Cable Management:}

  \begin{itemize}
  \tightlist
  \item
    Use cable carriers (drag chains) sized for 2x bend radius of largest
    cable
  \item
    Provide strain relief at moving carriage
  \item
    Route away from chip zones
  \item
    Allow for future additions (extra 30\% capacity)
  \end{itemize}
\end{enumerate}

\textbf{Maintenance Schedule:}

{\def\LTcaptype{} % do not increment counter
\begin{longtable}[]{@{}
  >{\raggedright\arraybackslash}p{(\linewidth - 4\tabcolsep) * \real{0.3438}}
  >{\raggedright\arraybackslash}p{(\linewidth - 4\tabcolsep) * \real{0.3125}}
  >{\raggedright\arraybackslash}p{(\linewidth - 4\tabcolsep) * \real{0.3438}}@{}}
\toprule\noalign{}
\begin{minipage}[b]{\linewidth}\raggedright
Component
\end{minipage} & \begin{minipage}[b]{\linewidth}\raggedright
Interval
\end{minipage} & \begin{minipage}[b]{\linewidth}\raggedright
Procedure
\end{minipage} \\
\midrule\noalign{}
\endhead
\bottomrule\noalign{}
\endlastfoot
Rail preload & 6 months & Verify bearing preload, adjust if needed \\
Screw lubrication & 3 months & Re-grease ball nut via zerk fitting \\
Counterbalance & 12 months & Verify force +/-10\%, adjust gas spring
pressure \\
Rail parallelism & 12 months & Height gauge measurement, shim if
needed \\
Resonance check & 12 months & Accelerometer test, update notch
filters \\
\end{longtable}
}

\begin{center}\rule{0.5\linewidth}{0.5pt}\end{center}

\subsection{1.6 Integrated Design Philosophy
Summary}\label{integrated-design-philosophy-summary}

The design of precision vertical axes requires simultaneous optimization
of multiple competing objectives:

\begin{enumerate}
\def\labelenumi{\arabic{enumi}.}
\tightlist
\item
  \textbf{Minimize moving mass} -\textgreater{} Reduced motor torque,
  faster acceleration
\item
  \textbf{Maximize structural stiffness} -\textgreater{} Dimensional
  stability under load
\item
  \textbf{Ensure thermal symmetry} -\textgreater{} Minimize position
  errors from temperature
\item
  \textbf{Achieve high natural frequency} -\textgreater{} Stable servo
  control without resonance
\item
  \textbf{Enable maintenance access} -\textgreater{} Long-term precision
  retention
\end{enumerate}

\textbf{Design Workflow:}

\begin{verbatim}
1. Define performance requirements (travel, speed, cutting forces)
2. Calculate required column stiffness -> Select cross-section
3. Estimate moving mass -> Design counterbalance system
4. Verify natural frequency > 5x servo bandwidth
5. Perform thermal FEA -> Verify symmetric expansion
6. Design maintenance access features
7. Prototype and measure -> Iterate
\end{verbatim}

The successful vertical axis represents a balanced compromise where no
single parameter is over-optimized at the expense of overall system
performance. The following sections provide detailed implementation
guidance for each subsystem.

\begin{center}\rule{0.5\linewidth}{0.5pt}\end{center}

\begin{center}\rule{0.5\linewidth}{0.5pt}\end{center}

\subsection{References}\label{references}

\begin{enumerate}
\def\labelenumi{\arabic{enumi}.}
\tightlist
\item
  \textbf{ISO 13849-1:2015} - Safety of machinery - Safety-related parts
  of control systems
\item
  \textbf{EN 60204-1:2018} - Safety of machinery - Electrical equipment
  - General requirements
\item
  \textbf{ANSI/NFPA 79-2021} - Electrical Standard for Industrial
  Machinery
\item
  \textbf{Mayr Antriebstechnik Safety Brake Catalog} - Spring-applied
  brake sizing and selection
\item
  \textbf{Warner Electric Brake Catalog} - Electromagnetic and
  spring-set brake specifications
\item
  \textbf{Slocum, A.H. (1992).} \emph{Precision Machine Design}. SME. -
  Gravity compensation systems
\end{enumerate}

\newpage

\section{Module 2 - Vertical Axis \&
Z-Stage}\label{module-2---vertical-axis-z-stage}

\subsection{Overview}\label{overview}

Precision assembly and alignment of vertical axis components determines
achievable accuracy and long-term reliability. This section provides
systematic procedures for column installation, linear rail tramming,
ball screw alignment, and final system verification.

\subsection{Pre-Assembly Preparation}\label{pre-assembly-preparation}

\subsubsection{Base/Frame Verification}\label{baseframe-verification}

\textbf{Critical mounting surface requirements:} - Flatness: 20 µm over
mounting area - Perpendicularity to X-Y plane: 50 µm/m - Surface finish:
3.2 µm Ra or better - Tapped holes: Clean threads, proper depth

\textbf{Inspection procedure:} 1. Place precision granite square or
angle plate on base 2. Indicate mounting surface flatness (sweep in grid
pattern) 3. Check perpendicularity with precision square and indicator
4. Verify hole locations against drawing (+/-0.1mm tolerance) 5. Clean
and degrease mounting surface

\textbf{If surface not within specification:} - Scrape or hand-grind
high spots - Re-machine if deviation \textgreater100 µm - Use shims only
as last resort (reduces rigidity)

\subsubsection{Component Inspection}\label{component-inspection}

\textbf{Inspect all components before assembly:}

\textbf{Column:} - Rail mounting surfaces flat (10 µm target) - Ball
screw mounting faces parallel and square - No burrs, damage, or debris -
Protective coatings removed from precision surfaces

\textbf{Linear rails:} - Check rail straightness with comparator (+/-5
µm) - Inspect for damage during shipping - Verify carriage preload by
hand feel (smooth, consistent resistance) - Rails match as set (often
marked L/R or 1/2)

\textbf{Ball screw:} - Bearing seats clean and undamaged - Thread runout
\textless10 µm TIR (check with indicator) - Couplings inspected, set
screws present

\subsubsection{Tools Required}\label{tools-required}

\textbf{Precision instruments:} - Dial indicator (0.001mm / 0.5 µm
resolution) - Magnetic base and extensions - Precision parallel (ground
to +/-2 µm) - Granite angle plate (reference surface) - Feeler gauges
(metric, 0.02-1.0mm)

\textbf{Installation tools:} - Torque wrench (accurate to +/-4\%) -
Allen key set (quality, no rounding) - Soft-face mallet (assembly) -
Degreaser and lint-free wipes - Grease (specified for bearings and
guides)

\subsection{Column Installation}\label{column-installation}

\subsubsection{Step 1: Position and
Level}\label{step-1-position-and-level}

\textbf{Mounting the column to base:}

\begin{enumerate}
\def\labelenumi{\arabic{enumi}.}
\tightlist
\item
  \textbf{Apply thin coat of way oil} to base mounting surface (prevents
  galling, aids alignment)
\item
  \textbf{Position column} onto base mounting surface

  \begin{itemize}
  \tightlist
  \item
    Use lifting equipment for large columns (\textgreater20 kg)
  \item
    Avoid dropping or impact
  \end{itemize}
\item
  \textbf{Insert locating dowel pins} (if designed with dowels)

  \begin{itemize}
  \tightlist
  \item
    Dowels provide repeatable positioning
  \item
    Tight slip fit (H7/m6 tolerance typical)
  \end{itemize}
\item
  \textbf{Check perpendicularity} to base using precision square and
  indicator

  \begin{itemize}
  \tightlist
  \item
    Sweep column face at top and bottom
  \item
    Target: \textless25 µm deviation over column height
  \end{itemize}
\item
  \textbf{Shim if necessary} to achieve perpendicularity

  \begin{itemize}
  \tightlist
  \item
    Use precision ground shims (not feeler gauges)
  \item
    Minimum number of shim locations
  \item
    Shim entire contact area (not just corners)
  \end{itemize}
\end{enumerate}

\subsubsection{Step 2: Secure Column}\label{step-2-secure-column}

\textbf{Fastener installation sequence:}

\begin{enumerate}
\def\labelenumi{\arabic{enumi}.}
\tightlist
\item
  \textbf{Insert all mounting bolts} finger-tight
\item
  \textbf{Tighten bolts progressively} in cross-pattern

  \begin{itemize}
  \tightlist
  \item
    Start center, work outward
  \item
    Multiple passes, increasing torque each pass
  \end{itemize}
\item
  \textbf{Final torque} per specification

  \begin{itemize}
  \tightlist
  \item
    M8: 20-25 N*m typical
  \item
    M10: 40-50 N*m typical
  \item
    M12: 70-85 N*m typical
  \end{itemize}
\item
  \textbf{Re-check perpendicularity} after torquing

  \begin{itemize}
  \tightlist
  \item
    Torque can distort alignment
  \item
    If deviation \textgreater25 µm, loosen and re-align
  \end{itemize}
\end{enumerate}

\textbf{Important:} Use threadlocker (medium-strength, e.g., Loctite
243) on all structural fasteners.

\subsubsection{Step 3: Final
Verification}\label{step-3-final-verification}

\textbf{Column rigidity check:} 1. Attempt to deflect column by hand at
top (reasonable force) 2. No detectable movement (indicates proper
mounting) 3. If movement detected, check fastener torque and base
contact

\textbf{Perpendicularity documentation:} - Record final indicator
readings - Document shim locations and thickness - Photograph assembly
for future reference

\subsection{Linear Rail Installation}\label{linear-rail-installation}

\subsubsection{Step 1: Rail Mounting Surface
Preparation}\label{step-1-rail-mounting-surface-preparation}

\textbf{Critical step for long-term accuracy:}

\begin{enumerate}
\def\labelenumi{\arabic{enumi}.}
\tightlist
\item
  \textbf{Clean mounting surfaces} on column

  \begin{itemize}
  \tightlist
  \item
    Remove all oil, debris, protective coatings
  \item
    Solvent wipe followed by dry wipe
  \end{itemize}
\item
  \textbf{Inspect flatness} of each rail mounting surface

  \begin{itemize}
  \tightlist
  \item
    Sweep with indicator along full length
  \item
    Target: \textless10 µm deviation
  \item
    If \textgreater20 µm, surface requires machining or scraping
  \end{itemize}
\item
  \textbf{Apply thin film of grease} to mounting surface

  \begin{itemize}
  \tightlist
  \item
    Prevents corrosion, aids initial alignment
  \item
    Wipe excess (should be barely visible film)
  \end{itemize}
\end{enumerate}

\subsubsection{Step 2: First Rail
Installation}\label{step-2-first-rail-installation}

\textbf{Establishing reference rail:}

\begin{enumerate}
\def\labelenumi{\arabic{enumi}.}
\tightlist
\item
  \textbf{Position first rail} on mounting surface

  \begin{itemize}
  \tightlist
  \item
    Align rail length to column travel direction
  \item
    Center rail on mounting pad
  \end{itemize}
\item
  \textbf{Insert mounting bolts} finger-tight
\item
  \textbf{Tram rail} to column reference surface

  \begin{itemize}
  \tightlist
  \item
    Place precision parallel against rail side face
  \item
    Sweep parallel with indicator along full rail length
  \item
    Adjust rail position to minimize deviation
  \item
    Target: \textless20 µm TIR over full length
  \end{itemize}
\item
  \textbf{Tighten mounting bolts progressively}

  \begin{itemize}
  \tightlist
  \item
    Start at center, work toward ends
  \item
    Torque per rail manufacturer specification (typically 10-15 N*m for
    HGR15/20)
  \end{itemize}
\item
  \textbf{Re-check tramming} after torquing

  \begin{itemize}
  \tightlist
  \item
    Torque can shift rail
  \item
    If deviation \textgreater20 µm, loosen and repeat
  \end{itemize}
\end{enumerate}

\subsubsection{Step 3: Second Rail Installation
(Parallelism)}\label{step-3-second-rail-installation-parallelism}

\textbf{Most critical alignment step:}

\begin{enumerate}
\def\labelenumi{\arabic{enumi}.}
\tightlist
\item
  \textbf{Install carriages} on first rail
\item
  \textbf{Position precision parallel} between first and second rail
  mounting surfaces

  \begin{itemize}
  \tightlist
  \item
    Parallel spans between the two rails
  \item
    Parallel length matches rail spacing
  \end{itemize}
\item
  \textbf{Position second rail} using parallel as reference
\item
  \textbf{Indicate parallelism} at multiple positions along rail length

  \begin{itemize}
  \tightlist
  \item
    Top, middle, bottom of travel
  \item
    Adjust second rail until parallel to first rail
  \item
    Target: \textless20 µm deviation over full length
  \end{itemize}
\item
  \textbf{Tighten mounting bolts progressively}
\item
  \textbf{Final verification:}

  \begin{itemize}
  \tightlist
  \item
    Install carriage on second rail
  \item
    Slide both carriages manually along full travel
  \item
    Should move smoothly with consistent resistance
  \item
    Any binding indicates misalignment
  \end{itemize}
\end{enumerate}

\textbf{If binding detected:} - Loosen second rail - Use indicator to
find tight spot - Adjust rail at that location - Re-torque and test

\subsubsection{Step 4: Carriage
Installation}\label{step-4-carriage-installation}

\textbf{Installing carriages on rails:}

\begin{enumerate}
\def\labelenumi{\arabic{enumi}.}
\tightlist
\item
  \textbf{Apply lubricant} to rails (per manufacturer specification)
\item
  \textbf{Slide carriages} onto rails from end

  \begin{itemize}
  \tightlist
  \item
    Do not disassemble carriages (loses preload setting)
  \item
    Keep carriages oriented correctly (marked ``up'' or with arrow)
  \end{itemize}
\item
  \textbf{Distribute carriages} along rail

  \begin{itemize}
  \tightlist
  \item
    Space evenly for initial testing
  \end{itemize}
\item
  \textbf{Install end stops} (prevent carriage from running off rail)

  \begin{itemize}
  \tightlist
  \item
    Mechanical stops at both ends of travel
  \item
    10-20mm from maximum travel limits
  \end{itemize}
\end{enumerate}

\subsection{Ball Screw Installation}\label{ball-screw-installation}

\subsubsection{Step 1: Bearing Mounting}\label{step-1-bearing-mounting}

\textbf{Fixed-end bearing (bottom typical):}

\begin{enumerate}
\def\labelenumi{\arabic{enumi}.}
\tightlist
\item
  \textbf{Clean bearing seat} in column

  \begin{itemize}
  \tightlist
  \item
    No burrs, debris, or damage
  \end{itemize}
\item
  \textbf{Install angular contact bearings} per manufacturer
  instructions

  \begin{itemize}
  \tightlist
  \item
    Correct orientation (load side facing thrust direction)
  \item
    Preload shims if specified
  \end{itemize}
\item
  \textbf{Install bearing housing}

  \begin{itemize}
  \tightlist
  \item
    Torque fasteners per specification
  \item
    Lock nuts or threadlocker as required
  \end{itemize}
\item
  \textbf{Check bearing preload} by rotating screw by hand

  \begin{itemize}
  \tightlist
  \item
    Should rotate smoothly with slight resistance
  \item
    Too tight: excessive friction, premature wear
  \item
    Too loose: backlash, reduced stiffness
  \end{itemize}
\end{enumerate}

\textbf{Supported-end bearing (top typical):}

\begin{enumerate}
\def\labelenumi{\arabic{enumi}.}
\tightlist
\item
  \textbf{Install radial bearing} (allows thermal expansion)

  \begin{itemize}
  \tightlist
  \item
    Bearing free to slide axially
  \item
    Maintains lateral location only
  \end{itemize}
\item
  \textbf{Secure bearing housing}
\end{enumerate}

\textbf{Center support bearing (if used):} - Install after screw
positioned - Adjust for proper support without binding screw

\subsubsection{Step 2: Ball Screw
Alignment}\label{step-2-ball-screw-alignment}

\textbf{Critical for smooth motion and long life:}

\begin{enumerate}
\def\labelenumi{\arabic{enumi}.}
\tightlist
\item
  \textbf{Position ball screw} in fixed-end bearing
\item
  \textbf{Check screw runout} at supported end

  \begin{itemize}
  \tightlist
  \item
    Place indicator on screw surface near bearing
  \item
    Rotate screw by hand, observe runout
  \item
    Target: \textless10 µm TIR
  \end{itemize}
\item
  \textbf{Adjust screw position} to minimize runout

  \begin{itemize}
  \tightlist
  \item
    Shim bearing housing if necessary
  \item
    Eccentric bearing mount (if provided) simplifies adjustment
  \end{itemize}
\item
  \textbf{Secure supported-end bearing}
\item
  \textbf{Install ball nut} onto screw

  \begin{itemize}
  \tightlist
  \item
    Lubricate with specified grease
  \item
    Do not disassemble nut (loses ball preload)
  \end{itemize}
\end{enumerate}

\subsubsection{Step 3: Coupling
Installation}\label{step-3-coupling-installation}

\textbf{Motor-to-screw coupling:}

\begin{enumerate}
\def\labelenumi{\arabic{enumi}.}
\tightlist
\item
  \textbf{Align motor shaft} to ball screw

  \begin{itemize}
  \tightlist
  \item
    Motor mounted to column or carriage
  \item
    Shafts collinear (no angular or parallel offset)
  \end{itemize}
\item
  \textbf{Measure alignment} with dial indicator

  \begin{itemize}
  \tightlist
  \item
    Rotate both shafts, observe TIR on each
  \item
    Angular misalignment: \textless0.05mm TIR
  \item
    Parallel offset: \textless0.1mm
  \end{itemize}
\item
  \textbf{Install flexible coupling}

  \begin{itemize}
  \tightlist
  \item
    Helical beam coupling recommended (accommodates slight misalignment)
  \item
    Set screws on flats (not threads)
  \item
    Torque set screws per coupling specification
  \end{itemize}
\item
  \textbf{Verify alignment} after coupling installed

  \begin{itemize}
  \tightlist
  \item
    Rotate motor shaft by hand
  \item
    Smooth rotation, no binding or vibration
  \end{itemize}
\end{enumerate}

\textbf{Important:} Rigid couplings require near-perfect alignment
(\textless0.02mm). Use flexible couplings for easier assembly and better
tolerance of thermal expansion.

\subsection{Carriage Assembly and
Mounting}\label{carriage-assembly-and-mounting}

\subsubsection{Ball Nut to Carriage}\label{ball-nut-to-carriage}

\textbf{Connecting ball nut to moving carriage:}

\begin{enumerate}
\def\labelenumi{\arabic{enumi}.}
\tightlist
\item
  \textbf{Position carriage} at mid-travel
\item
  \textbf{Attach ball nut} to carriage plate

  \begin{itemize}
  \tightlist
  \item
    Ball nut flange typically bolted to carriage bottom
  \item
    Use washers to distribute load
  \item
    Torque bolts evenly (cross-pattern)
  \end{itemize}
\item
  \textbf{Check that carriage moves freely} when screw rotated

  \begin{itemize}
  \tightlist
  \item
    If binding: Check rail parallelism and screw alignment
  \item
    Screw must be parallel to rails (no skew)
  \end{itemize}
\end{enumerate}

\subsubsection{Carriage
Perpendicularity}\label{carriage-perpendicularity}

\textbf{Ensure carriage perpendicular to column:}

\begin{enumerate}
\def\labelenumi{\arabic{enumi}.}
\tightlist
\item
  \textbf{Mount precision angle plate} to carriage
\item
  \textbf{Indicate against column reference surface}

  \begin{itemize}
  \tightlist
  \item
    Sweep indicator across carriage face
  \item
    Adjust carriage mounting (if slotted holes provided)
  \item
    Target: \textless25 µm deviation across carriage width
  \end{itemize}
\item
  \textbf{Tighten carriage mounting bolts}
\end{enumerate}

\subsection{Counterbalance
Installation}\label{counterbalance-installation}

\subsubsection{Gas Spring Mounting}\label{gas-spring-mounting}

\textbf{Typical configuration for small-medium machines:}

\begin{enumerate}
\def\labelenumi{\arabic{enumi}.}
\tightlist
\item
  \textbf{Identify mounting points}

  \begin{itemize}
  \tightlist
  \item
    Fixed end: Column or frame
  \item
    Moving end: Carriage
  \end{itemize}
\item
  \textbf{Install gas spring} per manufacturer orientation

  \begin{itemize}
  \tightlist
  \item
    Piston rod down (oil stays at cylinder head)
  \item
    Mounting clevis with pivot pins
  \end{itemize}
\item
  \textbf{Verify force balance}

  \begin{itemize}
  \tightlist
  \item
    With motor de-energized and brake released (use caution--support
    carriage)
  \item
    Carriage should hold position at mid-travel
  \item
    Slight upward bias acceptable (+/-5\%)
  \end{itemize}
\item
  \textbf{Adjust if necessary}

  \begin{itemize}
  \tightlist
  \item
    Swap gas spring for different force rating
  \item
    Add/remove mass from carriage if adjustable
  \end{itemize}
\end{enumerate}

\subsubsection{Pneumatic Counterbalance}\label{pneumatic-counterbalance}

\textbf{For larger machines:}

\begin{enumerate}
\def\labelenumi{\arabic{enumi}.}
\tightlist
\item
  \textbf{Connect pneumatic cylinder} to air supply
\item
  \textbf{Install pressure regulator} in line
\item
  \textbf{Set pressure} to balance carriage weight

  \begin{itemize}
  \tightlist
  \item
    Starting point: P = F / A (force divided by piston area)
  \item
    Fine-tune by testing motion balance
  \end{itemize}
\item
  \textbf{Lock regulator setting} after adjustment
\end{enumerate}

\subsection{System-Level Alignment
Checks}\label{system-level-alignment-checks}

\subsubsection{Runout Verification}\label{runout-verification}

\textbf{Spindle-to-table alignment:}

\begin{enumerate}
\def\labelenumi{\arabic{enumi}.}
\tightlist
\item
  \textbf{Mount dial indicator} in spindle (if spindle installed)
\item
  \textbf{Sweep indicator} against flat surface on table
\item
  \textbf{Move Z-axis} over full travel while observing indicator
\item
  \textbf{Record runout:} Should be \textless50 µm for general milling,
  \textless25 µm for precision
\item
  \textbf{If excessive runout:}

  \begin{itemize}
  \tightlist
  \item
    Check column perpendicularity to base
  \item
    Verify rail parallelism
  \item
    Inspect ball screw alignment
  \end{itemize}
\end{enumerate}

\subsubsection{Travel Limits}\label{travel-limits}

\textbf{Verify full Z-axis travel:}

\begin{enumerate}
\def\labelenumi{\arabic{enumi}.}
\tightlist
\item
  \textbf{Jog carriage} to bottom limit (manually or with motor)
\item
  \textbf{Note position} (use DRO or mark reference)
\item
  \textbf{Jog to top limit}
\item
  \textbf{Measure total travel} and compare to specification
\item
  \textbf{Set software limits} in control system

  \begin{itemize}
  \tightlist
  \item
    Soft limits 5-10mm inside mechanical limits
  \item
    Prevents crashes
  \end{itemize}
\end{enumerate}

\subsubsection{Backlash Measurement}\label{backlash-measurement}

\textbf{Check system backlash:}

\begin{enumerate}
\def\labelenumi{\arabic{enumi}.}
\tightlist
\item
  \textbf{Mount dial indicator} against carriage (reading Z-direction)
\item
  \textbf{Jog Z-axis upward} 10mm, zero indicator
\item
  \textbf{Command Z-axis downward} 10mm (observe when indicator begins
  moving)
\item
  \textbf{Backlash = commanded distance before indicator moves}
\item
  \textbf{Acceptable values:}

  \begin{itemize}
  \tightlist
  \item
    General purpose: \textless50 µm
  \item
    Precision: \textless20 µm
  \item
    Ultra-precision: \textless5 µm
  \end{itemize}
\item
  \textbf{If excessive backlash:}

  \begin{itemize}
  \tightlist
  \item
    Check ball nut preload (adjustable on some screws)
  \item
    Inspect coupling for play
  \item
    Verify rail carriage preload
  \end{itemize}
\end{enumerate}

\subsection{Cable Management
Installation}\label{cable-management-installation}

\subsubsection{Cable Carrier Mounting}\label{cable-carrier-mounting}

\textbf{Per Section 2.9 guidelines:}

\begin{enumerate}
\def\labelenumi{\arabic{enumi}.}
\tightlist
\item
  \textbf{Mount fixed end} of carrier

  \begin{itemize}
  \tightlist
  \item
    Top of column or frame (for top-mount configuration)
  \item
    Rigid mounting bracket
  \item
    Strain relief for cables exiting carrier
  \end{itemize}
\item
  \textbf{Mount moving end} to carriage

  \begin{itemize}
  \tightlist
  \item
    Secure bracket to carriage plate
  \item
    Flexible conduit for last 150mm to components
  \end{itemize}
\item
  \textbf{Route cables} through carrier

  \begin{itemize}
  \tightlist
  \item
    Power on one side, signal on other (EMI separation)
  \item
    Secure cables every 200-300mm inside carrier
  \item
    Maintain bend radius (\textgreater=10x largest cable diameter)
  \end{itemize}
\item
  \textbf{Test full travel}

  \begin{itemize}
  \tightlist
  \item
    Move carriage through full range
  \item
    Verify carrier moves smoothly, no snagging
  \item
    Check cable strain relief at both ends
  \end{itemize}
\end{enumerate}

\subsection{Documentation}\label{documentation}

\subsubsection{Assembly Record}\label{assembly-record}

\textbf{Document final configuration:} - Shim locations and thicknesses
- Fastener torque values and sequence - Alignment measurements
(parallelism, perpendicularity, runout) - Counterbalance force and
setting - Lubrication type and locations - Photos of critical areas

\textbf{Purpose:} Enables future troubleshooting, maintenance, and
rebuilds.

\subsubsection{Alignment Report}\label{alignment-report}

\textbf{Create measurement record:}

\begin{verbatim}
Z-Axis Alignment Report
-----------------------
Date: __________
Technician: __________

Column Perpendicularity:
  - X-axis: _____ µm deviation over _____ mm height
  - Y-axis: _____ µm deviation over _____ mm height

Rail Parallelism:
  - Rail 1 to column reference: _____ µm TIR
  - Rail 2 to Rail 1: _____ µm deviation

Ball Screw Alignment:
  - Runout at fixed end: _____ µm TIR
  - Runout at supported end: _____ µm TIR
  - Coupling alignment: _____ µm offset

System Performance:
  - Total Z travel: _____ mm
  - Spindle runout over travel: _____ µm
  - Backlash: _____ µm
  - Counterbalance force: _____ N (target: _____ N)

Notes:
_______________________________________
\end{verbatim}

\subsection{Common Assembly Issues}\label{common-assembly-issues}

\subsubsection{Rail Binding}\label{rail-binding}

\textbf{Symptom:} Carriage difficult to move, uneven resistance

\textbf{Causes:} - Rails not parallel (most common) - Mounting surface
not flat - Over-torqued mounting bolts (distorts rail) - Debris under
rail

\textbf{Solution:} 1. Loosen all rail mounting bolts 2. Check mounting
surface flatness 3. Re-tram rails with precision parallel 4. Torque
bolts progressively, re-check after each pass 5. Test carriage motion
frequently during assembly

\subsubsection{Ball Screw Binding}\label{ball-screw-binding}

\textbf{Symptom:} Rough rotation, tight spots, motor stalling

\textbf{Causes:} - Screw not aligned to rails (skewed) - Bearing
over-preloaded - Coupling misalignment - Ball nut to carriage mounting
stressed

\textbf{Solution:} 1. Check screw-to-rail parallelism (should be within
0.1mm over travel) 2. Verify bearing preload (should rotate smoothly by
hand) 3. Re-check coupling alignment 4. Loosen ball nut mounting, verify
free rotation, re-torque gently

\subsubsection{Excessive Runout}\label{excessive-runout}

\textbf{Symptom:} Spindle runout \textgreater50 µm over Z travel

\textbf{Causes:} - Column not perpendicular to base - Rails not parallel
- Column deflection under carriage weight

\textbf{Solution:} 1. Re-check column perpendicularity, shim if needed
2. Verify rail tramming 3. If deflection issue, stiffen column (Section
2.3) or reduce moving mass

\subsection{Key Takeaways}\label{key-takeaways}

\begin{enumerate}
\def\labelenumi{\arabic{enumi}.}
\tightlist
\item
  \textbf{Base flatness and perpendicularity} are foundation for all
  subsequent alignment
\item
  \textbf{Rail tramming} requires precision parallel and
  indicator--target \textless20 µm parallelism
\item
  \textbf{Rail parallelism} critical: Binding indicates misalignment
\item
  \textbf{Ball screw runout} should be \textless10 µm TIR at bearing
  seats
\item
  \textbf{Coupling alignment} prevents premature bearing wear (use
  flexible coupling)
\item
  \textbf{Counterbalance verification} before motor tuning (release
  brake, observe balance)
\item
  \textbf{System-level checks} validate component assembly (runout over
  travel, backlash)
\item
  \textbf{Documentation} enables future maintenance and troubleshooting
\item
  \textbf{Progressive torquing} prevents distortion (cross-pattern,
  multiple passes)
\item
  \textbf{Test frequently} during assembly: Small adjustments easier
  than complete disassembly
\end{enumerate}

\begin{center}\rule{0.5\linewidth}{0.5pt}\end{center}

\textbf{Next}: \href{section-2.11-testing-commissioning.md}{Section 2.11
- Testing and Commissioning}

\textbf{Previous}: \href{section-2.9-cable-management.md}{Section 2.9 -
Cable Management}

\newpage

\section{Module 2 - Vertical Axis \&
Z-Stage}\label{module-2---vertical-axis-z-stage-1}

\subsection{Overview}\label{overview-1}

Systematic testing validates mechanical assembly and establishes
baseline performance before production use. This section covers
mechanical verification, counterbalance optimization, servo tuning
procedures, and performance validation specific to vertical axis
systems.

\subsection{Pre-Commissioning Safety
Checks}\label{pre-commissioning-safety-checks}

\subsubsection{Mechanical Safety
Verification}\label{mechanical-safety-verification}

\textbf{Critical checks before applying power:}

\begin{enumerate}
\def\labelenumi{\arabic{enumi}.}
\tightlist
\item
  \textbf{Emergency stop functionality}

  \begin{itemize}
  \tightlist
  \item
    E-stop button accessible and clearly marked
  \item
    E-stop circuit wired to brake (fail-safe)
  \item
    Test: Press E-stop, verify brake engages
  \end{itemize}
\item
  \textbf{Brake operation}

  \begin{itemize}
  \tightlist
  \item
    With power off, brake should be engaged (spring-applied type)
  \item
    Apply power, brake should release (audible click typical)
  \item
    Cut power, brake should re-engage immediately
  \end{itemize}
\item
  \textbf{Limit switches}

  \begin{itemize}
  \tightlist
  \item
    Upper and lower travel limits installed
  \item
    Switches positioned 5-10mm beyond desired travel endpoints
  \item
    Test continuity with multimeter
  \end{itemize}
\item
  \textbf{Carriage support}

  \begin{itemize}
  \tightlist
  \item
    Before releasing brake: Verify counterbalance supports carriage
  \item
    If insufficient, manually support carriage when testing
  \end{itemize}
\item
  \textbf{Guards and covers}

  \begin{itemize}
  \tightlist
  \item
    Moving parts enclosed (ball screw, coupling)
  \item
    Pinch points protected
  \item
    Cable carrier cannot be contacted during normal operation
  \end{itemize}
\end{enumerate}

\textbf{CRITICAL SAFETY RULE:} Never release the Z-axis brake without
verified counterbalance or manual support. Unsupported vertical axes can
fall rapidly, causing injury or damage.

\subsection{Mechanical Verification
Tests}\label{mechanical-verification-tests}

\subsubsection{Test 1: Static Stiffness}\label{test-1-static-stiffness}

\textbf{Objective:} Verify column deflection under load matches design
calculations.

\textbf{Equipment:} - Calibrated weights (5-50 kg depending on machine
size) - Dial indicator (0.001mm resolution) - Magnetic base

\textbf{Procedure:} 1. \textbf{Mount indicator} against spindle mount or
carriage - Reading perpendicular to column face (X or Y direction) -
Zero indicator with no load 2. \textbf{Apply known force} at tool
mounting point - Hang calibrated weight from spindle nose - Typical
test: 10-20 kg (100-200 N force) 3. \textbf{Record deflection} from
indicator 4. \textbf{Calculate actual stiffness:} {[}k\_\{actual\} =
\frac{F}{\delta}{]} Where F = applied force (N), delta = measured
deflection (mm) 5. \textbf{Compare to design value} (from Section 2.3
calculations) - Actual should be \textgreater=80\% of designed stiffness
- If lower: Check column mounting, rail installation, fastener torque

\textbf{Example:} - Applied force: 150 N (15.3 kg weight) - Measured
deflection: 35 µm - Actual stiffness: 150 / 0.035 = 4,286 N/mm - Design
stiffness: 5,000 N/mm - Ratio: 86\% (acceptable, within 80\% threshold)

\subsubsection{Test 2: Backlash
Measurement}\label{test-2-backlash-measurement}

\textbf{Objective:} Quantify total system backlash.

\textbf{Equipment:} - Dial indicator - Magnetic base - Manual control or
MDI (Manual Data Input)

\textbf{Procedure:} 1. \textbf{Position carriage} at mid-travel 2.
\textbf{Mount indicator} against carriage, reading Z-direction 3.
\textbf{Command upward motion} (+Z) 5mm, wait for settling 4.
\textbf{Zero indicator} 5. \textbf{Command downward motion} (-Z) 0.1mm
increments 6. \textbf{Note commanded distance} before indicator begins
moving - This dead zone = backlash 7. \textbf{Repeat at 3-5 locations}
along Z-travel (top, middle, bottom) 8. \textbf{Record maximum backlash}

\textbf{Acceptance criteria:} - General purpose: \textless50 µm -
Precision work: \textless20 µm - Ultra-precision: \textless5 µm

\textbf{If excessive:} - Check ball screw preload (may be adjustable) -
Inspect coupling for slop - Verify motor mounting (no deflection under
load)

\subsubsection{Test 3: Runout Over
Travel}\label{test-3-runout-over-travel}

\textbf{Objective:} Measure spindle perpendicularity variation over
Z-axis travel.

\textbf{Equipment:} - Dial indicator mounted in spindle or on carriage -
Precision ground disc or surface plate on table

\textbf{Procedure:} 1. \textbf{Mount indicator} in spindle (or attach to
carriage) - Indicator tip contacts table reference surface 2.
\textbf{Position Z-axis} at lowest point in travel 3. \textbf{Zero
indicator} 4. \textbf{Traverse Z-axis} to highest point while observing
indicator 5. \textbf{Record total indicated runout (TIR)} over full
travel 6. \textbf{Repeat in X and Y directions} (rotate indicator
orientation)

\textbf{Acceptance criteria:} - Hobby/light-duty: \textless100 µm TIR -
General purpose: \textless50 µm TIR - Precision: \textless25 µm TIR

\textbf{If excessive:} - Check column perpendicularity to base - Verify
rail parallelism (re-tram if necessary) - Inspect for column deflection
under carriage weight

\subsubsection{Test 4: Smooth Motion
Test}\label{test-4-smooth-motion-test}

\textbf{Objective:} Verify smooth, consistent motion without binding or
stick-slip.

\textbf{Equipment:} - Motor drive (low-speed jog) - Observation

\textbf{Procedure:} 1. \textbf{Enable motor} and release brake 2.
\textbf{Jog Z-axis upward} at slow speed (50-100 mm/min) - Observe
motion: Should be smooth and continuous - Listen for unusual noises
(grinding, clicking) 3. \textbf{Jog Z-axis downward} at same speed 4.
\textbf{Repeat at multiple speeds} (100, 500, 1000 mm/min) 5.
\textbf{Note any irregularities:} - Binding or tight spots (alignment
issue) - Jerky motion (stick-slip from friction or servo tuning) -
Vibration (imbalance, coupling misalignment)

\textbf{Resolution:} - Binding: Re-check rail parallelism and ball screw
alignment - Stick-slip: Improve lubrication, adjust servo gains -
Vibration: Balance toolholder, check coupling alignment

\subsection{Counterbalance
Optimization}\label{counterbalance-optimization}

\subsubsection{Force Balance
Verification}\label{force-balance-verification}

\textbf{Objective:} Adjust counterbalance to neutralize gravitational
load.

\textbf{Equipment:} - Motor current meter (drive display or
oscilloscope) - Position feedback (encoder readout or DRO)

\textbf{Procedure:}

\textbf{Method 1: Static Balance Test (Preferred)} 1. \textbf{Position
carriage} at mid-travel 2. \textbf{Enable motor} (servo on, zero torque
command) 3. \textbf{Release brake carefully} (support carriage if
balance uncertain) 4. \textbf{Observe carriage motion:} - Sinks
downward: Counterbalance too weak - Rises upward: Counterbalance too
strong - Holds position: Balanced 5. \textbf{Adjust counterbalance:} -
Gas spring: Replace with different force rating - Pneumatic: Adjust air
pressure regulator - Weight-cable: Add/remove counterweight mass 6.
\textbf{Re-test until carriage holds position} (+/-5\% drift acceptable)

\textbf{Method 2: Current Monitoring (Quantitative)} 1. \textbf{Jog
carriage upward} at constant velocity (e.g., 500 mm/min) 2.
\textbf{Record motor current} during constant-velocity phase (RMS value)
3. \textbf{Jog carriage downward} at same velocity 4. \textbf{Record
motor current} during descent 5. \textbf{Calculate balance error:}
{[}\text{Balance error} = \frac{I_{up} + I_{down}}{2}{]} Ideal balance:
(I\_\{up\} = -I\_\{down\}) (equal magnitude, opposite direction) 6.
\textbf{Adjust counterbalance} to minimize balance error 7.
\textbf{Repeat until balance error \textless10\% of rated current}

\textbf{Example:} - Upward current: 1.2 A - Downward current: -0.3 A -
Balance error: (1.2 + (-0.3)) / 2 = 0.45 A - Target: \textless0.5 A for
5A rated drive (10\% threshold) - Result: Acceptable balance

\subsubsection{Dynamic Balance
Verification}\label{dynamic-balance-verification}

\textbf{Test acceleration symmetry:}

\begin{enumerate}
\def\labelenumi{\arabic{enumi}.}
\tightlist
\item
  \textbf{Command rapid upward move} (e.g., +100mm at maximum
  acceleration)
\item
  \textbf{Record peak current} during acceleration
\item
  \textbf{Command rapid downward move} (-100mm, same acceleration)
\item
  \textbf{Record peak current}
\item
  \textbf{Compare peak currents:}

  \begin{itemize}
  \tightlist
  \item
    Good balance: Peak currents within 20\% of each other
  \item
    Poor balance: Significant asymmetry (adjust counterbalance)
  \end{itemize}
\end{enumerate}

\subsection{Electrical Commissioning}\label{electrical-commissioning}

\subsubsection{Motor Phasing
Verification}\label{motor-phasing-verification}

\textbf{Ensure motor wired correctly:}

\textbf{Symptoms of reversed phasing:} - Motor runs in opposite
direction than commanded - Motor vibrates or doesn't move smoothly

\textbf{Correction:} - Swap any two motor power leads (for 3-phase) -
Update control configuration if needed

\textbf{Encoder Direction:} - Jog motor in +Z direction, verify encoder
count increases - If reversed: Swap encoder A+/A- signals or configure
in drive

\subsubsection{Home Switch Setup}\label{home-switch-setup}

\textbf{Establish machine zero reference:}

\begin{enumerate}
\def\labelenumi{\arabic{enumi}.}
\tightlist
\item
  \textbf{Install home switch} at known Z position

  \begin{itemize}
  \tightlist
  \item
    Typically near top of travel
  \item
    Repeatable trigger point (proximity or mechanical switch)
  \end{itemize}
\item
  \textbf{Configure homing routine} in control

  \begin{itemize}
  \tightlist
  \item
    Approach speed (slow, 100-200 mm/min for accuracy)
  \item
    Home offset (distance from switch to actual zero)
  \end{itemize}
\item
  \textbf{Test homing sequence:}

  \begin{itemize}
  \tightlist
  \item
    Command home (G28 Z0 or control-specific command)
  \item
    Verify carriage moves to switch, triggers, and backs off
  \item
    Repeatability: Home 5x times, measure variation (\textless10 µm
    target)
  \end{itemize}
\end{enumerate}

\subsubsection{Limit Switch Testing}\label{limit-switch-testing}

\textbf{Verify travel limit protection:}

\begin{enumerate}
\def\labelenumi{\arabic{enumi}.}
\tightlist
\item
  \textbf{Jog to upper limit} until switch triggers

  \begin{itemize}
  \tightlist
  \item
    Motion should stop immediately
  \item
    Record position
  \end{itemize}
\item
  \textbf{Jog to lower limit} until switch triggers

  \begin{itemize}
  \tightlist
  \item
    Motion should stop
  \item
    Record position
  \end{itemize}
\item
  \textbf{Calculate actual travel:} Upper position - Lower position
\item
  \textbf{Set soft limits} in control software

  \begin{itemize}
  \tightlist
  \item
    5-10mm inside hard limits (prevents crashes)
  \end{itemize}
\item
  \textbf{Test soft limits:}

  \begin{itemize}
  \tightlist
  \item
    Command move beyond soft limit
  \item
    Control should reject move or stop at limit
  \end{itemize}
\end{enumerate}

\subsection{Servo Tuning}\label{servo-tuning}

\subsubsection{Prerequisites}\label{prerequisites}

\textbf{Before tuning:} 1. Counterbalance optimized (balance error
\textless10\%) 2. All mechanical issues resolved (no binding, smooth
motion) 3. Encoder and motor phasing verified 4. Safety systems
operational (E-stop, limits, brake)

\subsubsection{Tuning Procedure}\label{tuning-procedure}

\textbf{Step 1: Disable All Gains} - Set Kp, Ki, Kd, FF to zero - Motor
should not respond to position commands (only brake holds position)

\textbf{Step 2: Proportional Gain (Kp)} 1. \textbf{Increase Kp
gradually} from zero - Start: Kp = 10-50 (depending on drive units) -
Increment: Double Kp each step 2. \textbf{Test response} after each
increase - Command small move (+/-10mm) - Observe: System should become
more responsive 3. \textbf{Continue until oscillation} begins (carriage
oscillates around target) 4. \textbf{Note Kp at onset of oscillation}
(critical gain, Kc) 5. \textbf{Reduce Kp to 40-60\% of Kc} for stability
margin - Example: Kc = 500, use Kp = 200-300

\textbf{Step 3: Derivative Gain (Kd)} 1. \textbf{Add derivative damping}
to reduce overshoot - Start: Kd = Kp / 10 - Adjust: Increase if
oscillation persists, decrease if sluggish 2. \textbf{Test with rapid
moves} - 100mm move at high acceleration - Observe settling time and
overshoot 3. \textbf{Target: Critically damped response} - Minimal
overshoot (\textless5\%) - Fast settling (\textless200ms for typical
servo)

\textbf{Step 4: Integral Gain (Ki)} 1. \textbf{Add integral to eliminate
steady-state error} - Start: Ki = Kp / 100 - Increase slowly (can
destabilize system if too high) 2. \textbf{Test positioning accuracy} -
Command position, wait for settling - Check final error (should be
\textless1 encoder count) 3. \textbf{Increase Ki until error eliminated}
but before oscillation begins

\textbf{Step 5: Feed-Forward Tuning}

\textbf{Velocity Feed-Forward (FFv):} - Reduces following error during
constant velocity - Start: FFv = 0.8 - Increase to 1.0 if following
error present during moves - Measure following error: Position error
during constant-velocity phase

\textbf{Acceleration Feed-Forward (FFa):} - Reduces following error
during acceleration - Start: FFa = 0.5 - Adjust based on following error
during accel/decel - Too high: Overshoot at end of move - Too low: Lag
during acceleration

\textbf{Step 6: Gravity Compensation (If Available)} - Some drives allow
constant torque offset - Set to compensate for residual gravitational
torque (if counterbalance not perfect) - Value: Feed-forward torque =
T\_gravity (from Section 2.6)

\subsubsection{Tuning Verification
Tests}\label{tuning-verification-tests}

\textbf{Test 1: Step Response} - Command 50mm move from rest - Record
position vs time - Analyze: - Rise time (time to reach 90\% of target) -
Overshoot (should be \textless5\%) - Settling time (time to reach +/-10
µm of target)

\textbf{Test 2: Following Error} - Command constant-velocity move (500
mm/min) - Monitor following error (commanded position - actual position)
- Target: \textless100 µm during motion, \textless10 µm at rest

\textbf{Test 3: Bi-Directional Repeatability} - Move to position from
above (+Z approach) - Record final position - Move away, return from
below (-Z approach) - Record final position - Repeatability = difference
between approaches - Target: \textless20 µm for precision work

\textbf{Test 4: High-Speed Moves} - Command rapid traverse at maximum
speed - Verify no overshoot or instability - Check motor current (should
not exceed rated)

\subsection{Performance Validation}\label{performance-validation}

\subsubsection{Accuracy Testing}\label{accuracy-testing}

\textbf{Circular interpolation test (if multi-axis):} 1. Program
circular move in XZ or YZ plane 2. Measure actual path with indicator or
laser interferometer 3. Calculate circularity error 4. Target:
\textless50 µm for general purpose

\textbf{Positioning accuracy:} 1. Command move to known position (use
gage blocks or precision height gage) 2. Measure actual position 3.
Repeat at 5-10 positions across Z-travel 4. Calculate maximum error 5.
ISO 230-2 standard: Positioning accuracy (A) and repeatability (R)

\subsubsection{Thermal Stability Test}\label{thermal-stability-test}

\textbf{Objective:} Measure thermal drift over operating cycle.

\textbf{Equipment:} - Dial indicator (µm resolution) - Temperature
sensor (column, ambient) - Timer

\textbf{Procedure:} 1. \textbf{Cold start:} Machine at ambient
temperature 2. \textbf{Mount indicator} reading against spindle face or
carriage 3. \textbf{Zero indicator} 4. \textbf{Run warm-up cycle:} -
Spindle at 50-75\% rated speed - Z-axis jogging over full travel 5.
\textbf{Record indicator reading} every 10 minutes for 2 hours 6.
\textbf{Record temperatures} (column, spindle, ambient) 7. \textbf{Plot
drift vs time}

\textbf{Analysis:} - Initial drift rate: µm/hour (first 30 minutes) -
Stabilized drift: After 60-90 minutes (should be \textless10 µm/hour for
precision work) - Total drift: From cold to stabilized (indicates
warm-up requirement)

\textbf{Example results:} - 0-30 min: 45 µm expansion (high rate,
machine heating) - 30-60 min: 12 µm expansion (slowing) - 60-120 min:
\textless5 µm variation (thermally stable) - Conclusion: 60-minute
warm-up required before precision work

\subsubsection{Cutting Test}\label{cutting-test}

\textbf{Objective:} Validate performance under realistic cutting loads.

\textbf{Test piece:} - Aluminum or steel block - Face milling operation
(controlled depth of cut and feed rate)

\textbf{Procedure:} 1. \textbf{Mount test piece} on table 2.
\textbf{Program simple face milling} routine - Multiple passes in Z-axis
(stepping down) - Record spindle load, feed rate, depth of cut 3.
\textbf{Execute program} - Monitor motor current (should not exceed
rated) - Listen for unusual sounds (chatter, vibration) - Observe
surface finish 4. \textbf{Measure result:} - Surface flatness (sweep
with indicator) - Surface finish (visual or profilometer) - Dimensional
accuracy (depth of cut)

\textbf{Acceptance criteria:} - Flatness: \textless50 µm over machined
area - Surface finish: Appropriate for material and tooling -
Dimensional accuracy: +/-25 µm or better - No chatter or excessive
vibration

\subsubsection{Repeatability Test}\label{repeatability-test}

\textbf{Procedure:} 1. \textbf{Program routine:} Move to position,
dwell, return to start 2. \textbf{Execute 30 cycles} (statistical
significance) 3. \textbf{Measure final position} after each cycle -
Indicator against carriage reference surface 4. \textbf{Calculate
standard deviation} of position measurements 5. \textbf{Repeatability
(R) = 4x standard deviation} (per ISO 230-2)

\textbf{Target performance:} - Hobby/light-duty: R \textless{} 50 µm -
General purpose: R \textless{} 20 µm - Precision: R \textless{} 10 µm

\textbf{Example:} - 30 measurements: Mean = 0.0 µm, StdDev = 3.2 µm -
Repeatability: 4 x 3.2 = 12.8 µm (meets general-purpose target)

\subsection{Commissioning
Documentation}\label{commissioning-documentation}

\subsubsection{Performance Report}\label{performance-report}

\textbf{Create comprehensive record:}

\begin{verbatim}
Z-Axis Commissioning Report
============================
Date: __________
Machine ID: __________
Technician: __________

MECHANICAL VERIFICATION
-----------------------
Static Stiffness:
  Applied force: _____ N
  Deflection: _____ µm
  Measured stiffness: _____ N/mm
  Design stiffness: _____ N/mm
  Ratio: _____ %

Backlash: _____ µm (max over travel)

Runout Over Travel:
  X-direction: _____ µm TIR
  Y-direction: _____ µm TIR

Smooth Motion: [  ] Pass  [  ] Fail
  Notes: _________________________

COUNTERBALANCE
--------------
Balance method: [  ] Gas spring  [  ] Pneumatic  [  ] Weight-cable
Balance force: _____ N (target: _____ N)
Static balance: [  ] Holds position  [  ] Drifts up  [  ] Drifts down
Current balance error: _____ A (_____ % of rated)

SERVO TUNING
------------
Final gains:
  Kp = _____
  Ki = _____
  Kd = _____
  FFv = _____
  FFa = _____

Step response:
  Rise time: _____ ms
  Overshoot: _____ %
  Settling time: _____ ms

Following error (at 500 mm/min): _____ µm

PERFORMANCE VALIDATION
----------------------
Positioning accuracy: +/-_____ µm
Repeatability: +/-_____ µm
Thermal stability: _____ µm/hour (after warm-up)
Warm-up time required: _____ minutes

Cutting test results:
  Material: __________
  Depth of cut: _____ mm
  Feed rate: _____ mm/min
  Surface flatness: _____ µm
  Surface finish: [  ] Acceptable  [  ] Poor
  Dimensional accuracy: +/-_____ µm

ACCEPTANCE
----------
System meets performance specifications: [  ] Yes  [  ] No

Approved by: ________________  Date: __________

Notes:
_____________________________________________
\end{verbatim}

\subsubsection{Baseline Configuration
File}\label{baseline-configuration-file}

\textbf{Save control parameters:} - Servo gains (Kp, Ki, Kd, FF values)
- Counterbalance settings - Soft limits and home position -
Acceleration/velocity limits - Thermal compensation coefficients (if
used)

\textbf{Purpose:} Allows restoration of known-good configuration after
changes or troubleshooting.

\subsection{Troubleshooting Common
Issues}\label{troubleshooting-common-issues}

\subsubsection{Issue: Excessive Following
Error}\label{issue-excessive-following-error}

\textbf{Symptoms:} - Position lags behind commanded during moves -
``Following error'' alarm from drive

\textbf{Causes:} - Insufficient proportional gain (Kp too low) -
Counterbalance not adjusted (gravitational bias) - Mechanical friction
(binding)

\textbf{Solutions:} 1. Increase Kp (verify stability afterward) 2.
Optimize counterbalance to reduce gravitational load 3. Check for
mechanical binding, improve lubrication 4. Add velocity feed-forward
(FFv)

\subsubsection{Issue: Oscillation or
Instability}\label{issue-oscillation-or-instability}

\textbf{Symptoms:} - Carriage oscillates around target position -
High-pitched whine from motor - Cannot tune system stable

\textbf{Causes:} - Excessive proportional gain (Kp too high) -
Insufficient damping (Kd too low) - Mechanical resonance in structure

\textbf{Solutions:} 1. Reduce Kp to 40-60\% of critical gain 2. Increase
Kd for damping 3. Add velocity feedback filter (if available in drive)
4. Check for loose mechanical components (resonance source)

\subsubsection{Issue: Asymmetric Response (Up vs
Down)}\label{issue-asymmetric-response-up-vs-down}

\textbf{Symptoms:} - Different behavior moving up vs down - Overshoot in
one direction only - Current draw imbalanced

\textbf{Causes:} - Inadequate counterbalancing (gravitational bias) -
Asymmetric friction (e.g., brake dragging)

\textbf{Solutions:} 1. Optimize counterbalance force 2. Add gravity
feed-forward compensation in drive 3. Check brake releases fully
(measure air gap or check current) 4. Verify rails lubricated evenly

\subsubsection{Issue: Poor
Repeatability}\label{issue-poor-repeatability}

\textbf{Symptoms:} - Returns to different positions on repeated moves -
Backlash measurement inconsistent

\textbf{Causes:} - Mechanical backlash (ball screw, coupling) - Thermal
variation (position changes with temperature) - Inadequate servo gains
(position not held tightly)

\textbf{Solutions:} 1. Check and adjust ball screw preload 2. Inspect
coupling for wear or slop 3. Increase integral gain (Ki) to eliminate
steady-state error 4. Implement thermal compensation if drift related to
temperature

\subsection{Key Takeaways}\label{key-takeaways-1}

\begin{enumerate}
\def\labelenumi{\arabic{enumi}.}
\tightlist
\item
  \textbf{Safety first:} Verify brake and E-stop before any motion
  testing
\item
  \textbf{Mechanical before electrical:} Resolve all binding and
  alignment issues before servo tuning
\item
  \textbf{Counterbalance optimization} critical for good servo
  performance and reduced wear
\item
  \textbf{Systematic tuning:} Kp first (stability), then Kd (damping),
  then Ki (accuracy), finally FF
\item
  \textbf{Following error} indicates under-tuned system or mechanical
  issues
\item
  \textbf{Oscillation} indicates over-tuned system (reduce Kp, add Kd)
\item
  \textbf{Thermal stability testing} reveals warm-up requirements for
  precision work
\item
  \textbf{Documentation} enables future troubleshooting and maintenance
\item
  \textbf{Cutting test} validates system under realistic loads
  (essential final check)
\item
  \textbf{Baseline configuration} saved after successful commissioning
  (reference for future)
\end{enumerate}

\begin{center}\rule{0.5\linewidth}{0.5pt}\end{center}

\textbf{Next}: \href{section-2.12-conclusion.md}{Section 2.12 -
Conclusion}

\textbf{Previous}: \href{section-2.10-assembly-alignment.md}{Section
2.10 - Assembly and Alignment}

\newpage

\section{Module 2 - Vertical Axis \&
Z-Stage}\label{module-2---vertical-axis-z-stage-2}

\subsection{Module Summary}\label{module-summary}

This module has provided comprehensive coverage of vertical axis
(Z-axis) engineering, from fundamental challenges through practical
implementation. You've learned how gravity, structural mechanics,
thermal effects, and safety requirements combine to create design
constraints fundamentally different from horizontal axes.

\subsubsection{Key Concepts Reviewed}\label{key-concepts-reviewed}

\textbf{Section 2.1 - Introduction:} - Unique challenges of vertical
motion: gravity, cantilever mechanics, thermal expansion - Safety
imperatives for fall-protection - Design philosophy balancing mass,
stiffness, and dynamics

\textbf{Section 2.2 - Gravity and Mass Management:} - Continuous
gravitational loading and energy demands - Counterbalancing systems: gas
springs, weight-cable, pneumatic, hydraulic - Mass reduction through
material selection and topology optimization - Inertia management for
servo performance

\textbf{Section 2.3 - Column Structural Design:} - Cantilever beam
mechanics and deflection analysis - Cross-section optimization for
stiffness-to-weight ratio - Material selection: steel, aluminum, cast
iron trade-offs - Fabrication methods: welded, machined, cast

\textbf{Section 2.4 - Vertical Linear Guides:} - Gravitational preload
effects on guide selection - Moment loading from off-center masses -
Preload class selection (ZB/ZC for vertical) - Installation alignment
and safety features

\textbf{Section 2.5 - Ball Screws for Vertical Axes:} - Continuous axial
load from gravity - Critical speed and buckling analysis - Center
support bearings for long travels - Mechanical brake requirements

\textbf{Section 2.6 - Motor and Drive Sizing:} - Torque calculations
including gravitational component - Benefits of counterbalancing on
motor sizing - Inertia ratio optimization (1:1 to 3:1 ideal) - Servo
tuning for vertical applications

\textbf{Section 2.7 - Thermal Management:} - Vertical thermal expansion
directly affects part dimensions - Heat sources: spindle, friction,
environment - Design strategies: isolation, thermal mass, active cooling
- Software compensation techniques

\textbf{Section 2.8 - Spindle Mounting:} - Interface design for rigidity
and runout control - Cooling integration and thermal management - Tool
holder interface standards (BT, CAT, HSK, ER) - Vibration isolation and
damping

\textbf{Section 2.9 - Cable Management:} - Cable carrier selection and
mounting - Strain relief for vertical motion - Routing strategies to
avoid snagging - Power vs.~signal cable separation

\textbf{Section 2.10 - Assembly and Alignment:} - Column installation
and base mounting - Rail tramming and parallelism verification - Ball
screw alignment and bearing installation - System-level alignment checks

\textbf{Section 2.11 - Testing and Commissioning:} - Mechanical
verification: stiffness, backlash, runout - Counterbalance adjustment
and verification - Servo tuning procedures for vertical axes -
Performance validation and acceptance testing

\subsection{Integration with Other
Modules}\label{integration-with-other-modules}

\subsubsection{Mechanical Foundation (Modules
1-3)}\label{mechanical-foundation-modules-1-3}

\textbf{Module 1 (Frame Design):} - Column mounts to machine base/frame
- Base must provide rigid, flat mounting surface - Column weight and
cutting forces load frame structure - Perpendicularity between column
and X-Y plane critical

\textbf{Connection:}

\begin{verbatim}
Frame (Module 1) -> Column mounting interface -> Z-axis (Module 2)
\end{verbatim}

\textbf{Module 3 (Linear Motion):} - Linear guides covered in general in
Module 3 - Module 2 applies these principles specifically to vertical
mounting - Preload, lubrication, and alignment concepts transfer -
Vertical-specific considerations: gravitational loading, fall protection

\subsubsection{Control System (Module 4)}\label{control-system-module-4}

\textbf{Electronics Integration:} - Motor drives sized in Module 2,
wired in Module 4 - Brake control logic (fail-safe design) - Temperature
sensors for thermal compensation - Limit switches for Z-axis travel
protection

\textbf{Safety circuits:} - E-stop must engage Z-axis brake - Power loss
triggers brake automatically - Control system monitors brake status

\subsubsection{Process Modules (Modules
5-13)}\label{process-modules-modules-5-13}

\textbf{Spindle Selection:} - Module 6 covers spindle systems in detail
- Module 2 addresses mounting and integration - Thermal management
strategies apply to all spindle types - Plasma/laser/waterjet: Different
tool weights, mounting requirements

\textbf{Example:} - Plasma torch (Module 5): 2-3 kg, requires THC
integration - Spindle (Module 6): 5-15 kg, requires cooling and tool
change - Laser head (Module 7): 3-8 kg, requires focus control

\subsubsection{Software and Control (Modules
14-15)}\label{software-and-control-modules-14-15}

\textbf{Module 14 (LinuxCNC):} - HAL configuration for Z-axis -
Counterbalance feed-forward compensation - Thermal compensation via HAL
components - Brake control logic implementation

\textbf{Module 15 (G-Code):} - Z-axis motion commands (G00/G01 with Z
parameter) - Tool length compensation (G43) - Canned cycles for drilling
(G81-G89) - Safe retract heights in programs

\subsubsection{Manufacturing (Module 16)}\label{manufacturing-module-16}

\textbf{CAD/CAM Integration:} - Z-axis travel limits constrain part
height - Thermal stability affects achievable tolerances - Surface
finish dependent on Z-axis rigidity - Tool length offsets critical for
multi-tool operations

\subsection{Real-World Application
Synthesis}\label{real-world-application-synthesis}

\subsubsection{Small Benchtop Mill
Build}\label{small-benchtop-mill-build}

\textbf{Specifications:} - Z-travel: 200mm - Moving mass: 4 kg (compact
air-cooled spindle) - Target accuracy: +/-25 µm - Budget: (800 for
Z-axis components

\textbf{Design decisions (from this module):}

\begin{enumerate}
\def\labelenumi{\arabic{enumi}.}
\tightlist
\item
  \textbf{Column} (Section 2.3):

  \begin{itemize}
  \tightlist
  \item
    Material: Aluminum 6061 plate
  \item
    Construction: Pocketed from 100x80x20mm plate
  \item
    Weight: 1.2 kg
  \item
    Cost: )120
  \end{itemize}
\item
  \textbf{Linear guides} (Section 2.4):

  \begin{itemize}
  \tightlist
  \item
    HGR15 rails, 300mm length
  \item
    4 carriages, ZB preload
  \item
    Cost: (180
  \end{itemize}
\item
  \textbf{Ball screw} (Section 2.5):

  \begin{itemize}
  \tightlist
  \item
    1204 (12mm diameter, 4mm pitch)
  \item
    Fixed-supported mounting
  \item
    Cost: )85
  \end{itemize}
\item
  \textbf{Motor} (Section 2.6):

  \begin{itemize}
  \tightlist
  \item
    NEMA 17 servo with brake
  \item
    0.18 N\emph{m rated, 0.54 N}m peak
  \item
    Cost: (120
  \end{itemize}
\item
  \textbf{Counterbalance} (Section 2.2):

  \begin{itemize}
  \tightlist
  \item
    Single 40N gas spring
  \item
    Cost: )25
  \end{itemize}
\item
  \textbf{Spindle mount} (Section 2.8):

  \begin{itemize}
  \tightlist
  \item
    Machined aluminum
  \item
    ER11 spindle, 500W
  \item
    Cost: (200 (spindle + mount)
  \end{itemize}
\item
  \textbf{Cable carrier} (Section 2.9):

  \begin{itemize}
  \tightlist
  \item
    15x20mm plastic drag chain
  \item
    Cost: )30
  \end{itemize}
\end{enumerate}

\textbf{Total Z-axis cost:} (760 (within budget)

\textbf{Expected performance:} - Positioning accuracy: +/-20 µm (meets
target) - Maximum cutting force: 300 N - Thermal drift: \textless40
µm/hour (acceptable for hobby use)

\subsubsection{Industrial VMC Retrofit}\label{industrial-vmc-retrofit}

\textbf{Specifications:} - Z-travel: 750mm - Moving mass: 28 kg
(water-cooled spindle + ATC) - Target accuracy: +/-5 µm - Budget:
)15,000 for complete Z-axis rebuild

\textbf{Design decisions:}

\begin{enumerate}
\def\labelenumi{\arabic{enumi}.}
\tightlist
\item
  \textbf{Column} (Section 2.3):

  \begin{itemize}
  \tightlist
  \item
    Cast iron, custom casting
  \item
    Internal coolant passages
  \item
    Cost: (4,500
  \end{itemize}
\item
  \textbf{Linear guides} (Section 2.4):

  \begin{itemize}
  \tightlist
  \item
    HGR25 rails, 1000mm, six carriages (three pairs)
  \item
    ZC heavy preload
  \item
    Cost: )1,200
  \end{itemize}
\item
  \textbf{Ball screw} (Section 2.5):

  \begin{itemize}
  \tightlist
  \item
    3210 (32mm diameter, 10mm pitch), ground C7
  \item
    Fixed-supported with center support bearing
  \item
    Cost: (1,800
  \end{itemize}
\item
  \textbf{Motor} (Section 2.6):

  \begin{itemize}
  \tightlist
  \item
    Industrial servo, 1.5 kW
  \item
    Absolute encoder
  \item
    Cost: )2,200
  \end{itemize}
\item
  \textbf{Counterbalance} (Section 2.2):

  \begin{itemize}
  \tightlist
  \item
    Pneumatic cylinder system
  \item
    Adjustable pressure regulator
  \item
    Cost: (800
  \end{itemize}
\item
  \textbf{Thermal management} (Section 2.7):

  \begin{itemize}
  \tightlist
  \item
    Column coolant circulation
  \item
    Spindle chiller with precision control
  \item
    Temperature sensors and compensation
  \item
    Cost: )3,000
  \end{itemize}
\item
  \textbf{Spindle system} (Section 2.8):

  \begin{itemize}
  \tightlist
  \item
    7.5 kW water-cooled spindle
  \item
    BT40 interface
  \item
    Cost: (4,500
  \end{itemize}
\end{enumerate}

\textbf{Total project cost:} )18,000 (slightly over budget but achieves
precision target)

\textbf{Expected performance:} - Positioning accuracy: +/-3 µm (exceeds
target) - Repeatability: +/-2 µm - Thermal drift: \textless3 µm/hour -
Maximum cutting force: 3000 N

\subsection{Common Pitfalls and
Solutions}\label{common-pitfalls-and-solutions}

\subsubsection{Pitfall 1: Inadequate
Counterbalancing}\label{pitfall-1-inadequate-counterbalancing}

\textbf{Problem:} Motor oversized, continuous heating, thermal drift,
poor servo response

\textbf{Solution:} - Design counterbalance first, before selecting motor
- Aim for +/-5\% force balance across full travel - Verify balance by
releasing motor (with brake engaged for safety)

\subsubsection{Pitfall 2: Insufficient Column
Stiffness}\label{pitfall-2-insufficient-column-stiffness}

\textbf{Problem:} Excessive deflection under cutting forces, poor
surface finish, chatter

\textbf{Solution:} - Calculate deflection using cantilever formulas
(Section 2.3) - Target: \textless50 µm deflection under maximum expected
load - Increase cross-section depth (cubic relationship to stiffness) -
Add internal ribs if fabricated structure

\subsubsection{Pitfall 3: Thermal Drift
Ignored}\label{pitfall-3-thermal-drift-ignored}

\textbf{Problem:} Parts out of tolerance after machine warms up,
inconsistent dimensions

\textbf{Solution:} - Implement warm-up procedure (30-60 minutes before
precision work) - Use water-cooled spindle for stability - Add
temperature sensor and software compensation for critical applications -
Insulate column from environment

\subsubsection{Pitfall 4: Poor Linear Guide
Alignment}\label{pitfall-4-poor-linear-guide-alignment}

\textbf{Problem:} Binding during motion, premature wear, reduced
accuracy

\textbf{Solution:} - Machine column mounting surfaces flat (10-20 µm
flatness) - Use precision parallel during installation - Check smooth
motion by hand before final assembly - Shim as needed to eliminate
binding

\subsubsection{Pitfall 5: Brake Not
Fail-Safe}\label{pitfall-5-brake-not-fail-safe}

\textbf{Problem:} Safety hazard--spindle can fall on power loss

\textbf{Solution:} - Use spring-applied, electrically-released brake
(fail-safe design) - Size brake for 2-3x gravitational holding torque -
Test brake function before enabling motion - Integrate brake into E-stop
circuit

\subsubsection{Pitfall 6: Cable Management
Neglected}\label{pitfall-6-cable-management-neglected}

\textbf{Problem:} Cables snag, break, or interfere with motion

\textbf{Solution:} - Install proper cable carrier (Section 2.9) - Route
cables with adequate slack for full travel - Separate power and signal
cables - Secure mounting at both ends with strain relief

\subsection{Best Practices Summary}\label{best-practices-summary}

\textbf{Design Phase:} 1. Start with moving mass budget--reduce weight
first 2. Design counterbalance system--transforms motor sizing 3. Size
column for stiffness--deflection limits, not strength 4. Select
components with vertical-specific requirements in mind 5. Plan thermal
management based on precision requirements 6. Design for fail-safe
operation--brakes, redundancy

\textbf{Fabrication Phase:} 1. Machine mounting surfaces flat and
parallel 2. Stress-relieve welded structures before final machining 3.
Use precision tooling for assembly (parallel, indicators, height gauges)
4. Install components in logical sequence (column -\textgreater{} rails
-\textgreater{} screw -\textgreater{} carriage) 5. Check alignment at
each step before proceeding

\textbf{Assembly Phase:} 1. Verify base mounting surface flat and square
2. Install column with dowel pins for location 3. Tram rails to within
20 µm parallelism 4. Align ball screw to carriage motion axis 5. Install
counterbalance and verify balance 6. Route cables in cable carrier with
proper strain relief

\textbf{Commissioning Phase:} 1. Mechanical verification first
(stiffness, runout, backlash) 2. Electrical verification (motor phasing,
encoder, brake) 3. Tune servo starting with conservative gains 4. Test
in both directions for symmetric response 5. Validate performance
against specifications 6. Document final configuration for future
reference

\subsection{Continuing Education}\label{continuing-education}

\subsubsection{Recommended Next Steps}\label{recommended-next-steps}

\textbf{Hands-on experience:} - Build a Z-axis assembly (even small
scale) - Practice rail alignment and tramming techniques - Tune servo
system and observe response - Measure thermal drift under realistic
conditions

\textbf{Advanced topics:} - 5-axis kinematics (Z-axis on tilting head) -
Adaptive control for variable loads - High-speed machining
considerations - Ultra-precision machine design (sub-micron)

\textbf{Related skills:} - FEA for structural analysis - Thermal
modeling and simulation - Servo system theory and advanced tuning -
Metrology and precision measurement

\subsubsection{Resources}\label{resources}

\textbf{Books:} - \emph{Precision Machine Design} by Alex Slocum (MIT) -
\emph{Machine Tool Structures} by Koenigsberger \& Tlusty -
\emph{Fundamentals of Machine Elements} by Schmid \& Hamrock

\textbf{Standards:} - ISO 230 series (machine tool testing) - ISO 13041
(machining center accuracy) - ANSI B5 standards (machine tool accuracy)

\textbf{Online communities:} - CNCZone forums (practical builders) -
Practical Machinist (professional machinists) - LinuxCNC forums (control
integration)

\textbf{Software:} - FEA: SolidWorks Simulation, Fusion 360 - Thermal
analysis: ANSYS, COMSOL - Servo tuning: LinuxCNC, Mach4, industrial
drives

\subsection{Final Thoughts}\label{final-thoughts}

The vertical axis represents the apex of machine tool design challenges.
Every engineering decision--from material selection to thermal
management--must account for gravity's continuous influence. Yet this
challenge also presents opportunity: a well-designed Z-axis demonstrates
engineering excellence across multiple disciplines.

Key principles that ensure success:

\textbf{1. Respect gravity:} It's always there, always pulling down.
Design with it, not against it.

\textbf{2. Counterbalance always:} The benefits compound--better motor
sizing, lower thermal load, improved servo response, extended component
life.

\textbf{3. Stiffness matters:} Deflection limits precision. Invest in
structural rigidity.

\textbf{4. Thermal stability:} For precision work, thermal management is
not optional.

\textbf{5. Safety first:} Fail-safe design, mechanical brakes,
redundancy. The consequences of vertical axis failure are severe.

\textbf{6. Integration mindset:} The Z-axis doesn't exist in
isolation--it connects to frame, electronics, software, and process.
Design accordingly.

This module has equipped you with the engineering foundation to design,
build, and commission vertical axes for CNC machines. Combined with
knowledge from other modules, you now possess a complete understanding
of CNC machine design from first principles.

Welcome to the world of vertical axis engineering. May your Z-axes be
stiff, balanced, and thermally stable.

\begin{center}\rule{0.5\linewidth}{0.5pt}\end{center}

\textbf{Previous}: \href{section-2.11-testing-commissioning.md}{Section
2.11 - Testing and Commissioning}

\textbf{Next Module}:
\href{../Module-03/module-03-linear-motion.md}{Module 3 - Linear Motion
Systems}

\textbf{Return to}: \href{module-2-vertical-axis.md}{Module 2 Overview}

\newpage

\section{Module 2 - Vertical Axis and Column
Assembly}\label{module-2---vertical-axis-and-column-assembly-1}

\subsection{2. Vertical Axis
Architecture}\label{vertical-axis-architecture}

The architecture of a vertical axis system integrates mechanical,
electrical, and control subsystems into a cohesive design that delivers
precision vertical motion under gravitational load. This section
provides comprehensive specifications, design rationale, and engineering
trade-offs for each architectural element.

\subsubsection{2.1 Travel Range: Determining Working
Height}\label{travel-range-determining-working-height}

\textbf{Design Considerations:}

The Z-axis travel range determines the maximum part height that can be
machined and influences nearly every other design parameter. The
selection process balances workspace requirements against structural
complexity and cost.

\textbf{Typical Applications:}

{\def\LTcaptype{} % do not increment counter
\begin{longtable}[]{@{}
  >{\raggedright\arraybackslash}p{(\linewidth - 4\tabcolsep) * \real{0.3421}}
  >{\raggedright\arraybackslash}p{(\linewidth - 4\tabcolsep) * \real{0.3684}}
  >{\raggedright\arraybackslash}p{(\linewidth - 4\tabcolsep) * \real{0.2895}}@{}}
\toprule\noalign{}
\begin{minipage}[b]{\linewidth}\raggedright
Application
\end{minipage} & \begin{minipage}[b]{\linewidth}\raggedright
Travel Range
\end{minipage} & \begin{minipage}[b]{\linewidth}\raggedright
Rationale
\end{minipage} \\
\midrule\noalign{}
\endhead
\bottomrule\noalign{}
\endlastfoot
Sheet metal plasma/laser & 100-150 mm & Limited part thickness,
prioritize stiffness \\
General machining & 150-250 mm & Balanced capability for moderate
parts \\
Large component machining & 250-400 mm & Heavy industrial, requires
robust structure \\
3D printing (large format) & 400-1000 mm & Tall build volumes, lower
cutting forces \\
\end{longtable}
}

\textbf{Structural Scaling:}

As travel increases, structural requirements scale unfavorably:

{[}\delta \propto L\^{}3{]}

A 2x increase in travel requires 8x the structural moment of inertia to
maintain the same deflection specification, typically resulting in 3-4x
increase in column cross-section and mass.

\textbf{Design Example:}

\textbf{Requirement:} Plasma cutting of 25 mm steel plate - Plate
thickness: 25 mm - Torch standoff: 3-5 mm - Pierce height: 10 mm above
plate - Clearance for fixturing: 50 mm - Safety margin: 20 mm

\textbf{Calculation:} {[}Travel = 25 + 5 + 10 + 50 + 20 = 110
\text{ mm minimum}{]}

\textbf{Selected Travel:} 150 mm (provides operational flexibility)

\subsubsection{2.2 Drive Type Selection: Ball-Screw vs.~Belt
Drive}\label{drive-type-selection-ball-screw-vs.-belt-drive}

\textbf{Ball-Screw Drive:}

\emph{Advantages:} 1. High stiffness (minimal elastic deflection under
load) 2. Zero backlash (with preload) 3. High positioning resolution 4.
Predictable servo response 5. High force capacity (10-50 kN typical)

\emph{Disadvantages:} 1. Higher cost ((200-)800 for precision grades) 2.
Critical speed limitations 3. Requires lubrication maintenance 4.
Acoustic noise at high speeds

\textbf{Engineering Analysis:}

Ball-screw stiffness: {[}k\_\{screw\} = \frac{EA}{L}{]}

For Ø20 mm screw, 400 mm length, steel: {[}k\_\{screw\} =
\frac{200 \times 10^9 \times \pi \times 0.01^2}{0.4} = 157
\times 10\^{}6 \text{ N/m}{]}

This high stiffness (157 N/mum) ensures the drive mechanism does not
become the compliance-limiting element.

\textbf{Critical Speed:}

The critical rotational speed at which the screw becomes dynamically
unstable:

{[}n\_\{cr\} = \frac{k \times 10^6 \times d_r}{L^2}{]}

Where: - (k) = support factor (3.5 for fixed-supported, 4.7 for
fixed-free) - (d\_r) = root diameter (mm) - (L) = unsupported length
(mm)

For Ø20 mm screw, 400 mm length, fixed-supported: {[}n\_\{cr\} =
\frac{3.5 \times 10^6 \times 18}{400^2} = \frac{63 \times 10^6}{160,000}
= 393 \text{ rpm}{]}

At 5 mm pitch: (v\_\{max\} = 393 \times 0.005 = 1.97) m/min

\textbf{Belt Drive:}

\emph{Advantages:} 1. Lower cost ((50-)150 for complete drive) 2. No
critical speed limitations 3. Quiet operation 4. No lubrication
requirements 5. High speed capability (10-20 m/min feasible)

\emph{Disadvantages:} 1. Lower stiffness (belt elasticity) 2. Potential
backlash from belt stretch 3. Requires tension maintenance 4. Lower
force capacity (500-2000 N typical)

\textbf{Belt Stiffness Analysis:}

Effective stiffness of tensioned timing belt: {[}k\_\{belt\} =
\frac{E_{eff} \cdot A_{belt}}{L_{span}} +
\frac{T_{tension}}{elongation}{]}

For GT2 belt, 800 mm span, 3000 N tension: {[}k\_\{belt\} \approx 15
\times 10\^{}6 \text{ N/m}{]} (approximately 10x lower than ball-screw)

\textbf{Application Guide:}

{\def\LTcaptype{} % do not increment counter
\begin{longtable}[]{@{}lll@{}}
\toprule\noalign{}
Parameter & Choose Ball-Screw & Choose Belt \\
\midrule\noalign{}
\endhead
\bottomrule\noalign{}
\endlastfoot
Cutting forces & \textgreater{} 500 N & \textless{} 200 N \\
Positioning accuracy & \textless{} 0.05 mm & 0.05-0.20 mm acceptable \\
Speed requirement & \textless{} 10 m/min & \textgreater{} 10 m/min \\
Budget & Premium acceptable & Cost-sensitive \\
Maintenance access & Good & Limited \\
\end{longtable}
}

\textbf{For precision Z-axis applications (plasma, laser, milling):}
Ball-screw is strongly preferred due to superior stiffness and accuracy.

\subsubsection{2.3 Guide System: Linear Rails and Load
Distribution}\label{guide-system-linear-rails-and-load-distribution}

\textbf{Rail Selection Criteria:}

\begin{enumerate}
\def\labelenumi{\arabic{enumi}.}
\item
  \textbf{Size (Width x Height):} Common sizes: MGN12, MGN15, HGR20,
  HGR25, HGR30

  Selection based on load capacity and moment resistance:

  {\def\LTcaptype{} % do not increment counter
  \begin{longtable}[]{@{}llll@{}}
  \toprule\noalign{}
  Rail Type & Dynamic Load (kN) & Moment Capacity (N*m) & Application \\
  \midrule\noalign{}
  \endhead
  \bottomrule\noalign{}
  \endlastfoot
  MGN12 & 1.7 & 12 & Light duty, 3D printing \\
  MGN15 & 2.8 & 24 & Light machining, laser \\
  HGR20 & 4.8 & 65 & Medium machining, plasma \\
  HGR25 & 7.3 & 125 & Heavy machining, milling \\
  HGR30 & 10.8 & 215 & Industrial machining \\
  \end{longtable}
  }
\item
  \textbf{Preload Class:}

  \begin{itemize}
  \tightlist
  \item
    \textbf{Z0 (light preload):} 0.02C\_0 - General purpose, low
    friction
  \item
    \textbf{ZA (medium preload):} 0.04C\_0 - Precision positioning
  \item
    \textbf{ZB (heavy preload):} 0.08C\_0 - High stiffness, vibration
    resistance
  \end{itemize}
\end{enumerate}

\textbf{Critical Design Parameter: Rail Spacing}

The moment arm provided by wide rail spacing dramatically improves
pitching resistance:

{[}\theta\_\{pitch\} = \frac{M_{applied}}{k_{moment}}{]}

Where moment stiffness: {[}k\_\{moment\} = k\_\{rail\}
\times (d\_\{spacing\})\^{}2{]}

\textbf{Worked Example:}

\textbf{Given:} - Applied cutting force: (F = 400) N - Distance from
rail centerline to tool tip: (L = 150) mm - Applied moment: (M = 400
\times 0.15 = 60) N*m

\textbf{Case 1:} Rail spacing = 80 mm - Single rail stiffness:
(k\_\{rail\} = 500) N/mum = (5 \times 10\^{}8) N/m - Moment stiffness:
(k\_M = 5 \times 10\^{}8 \times 0.08\^{}2 = 3.2 \times 10\^{}6) N*m/rad
- Angular deflection: (\theta = 60 / (3.2 \times 10\^{}6) = 18.75
\times 10\^{}\{-6\}) rad - Tip deflection: (\delta = 150 \times 18.75
\times 10\^{}\{-6\} = 2.8) mum

\textbf{Case 2:} Rail spacing = 120 mm - Moment stiffness: (k\_M = 5
\times 10\^{}8 \times 0.12\^{}2 = 7.2 \times 10\^{}6) N*m/rad - Angular
deflection: (\theta = 60 / (7.2 \times 10\^{}6) = 8.33
\times 10\^{}\{-6\}) rad - Tip deflection: (\delta = 150 \times 8.33
\times 10\^{}\{-6\} = 1.25) mum

\textbf{Result:} 50\% increase in rail spacing (80-\textgreater120 mm)
reduces angular deflection by 55\% and tip deflection by 55\%.

\textbf{Design Rule:} {[}d\_\{spacing\} \geq 0.6 \times column\_width{]}
(minimum) {[}d\_\{spacing\} \geq 0.8 \times column\_width{]} (preferred)

\subsubsection{2.4 Motor Orientation: Vertical vs.~Horizontal
Mounting}\label{motor-orientation-vertical-vs.-horizontal-mounting}

\textbf{Vertical Mounting (Motor Above Axis):}

\emph{Advantages:} 1. Protected from chips and cutting fluids 2.
Simplified cable management (stationary connections) 3. Direct drive
possible (no gearbox needed) 4. Natural heat dissipation (convection
upward)

\emph{Disadvantages:} 1. Increases overall machine height 2. Adds
rotating mass to moving assembly 3. Cable flexing at motor shaft 4.
Requires sealed motor connectors

\textbf{Horizontal Mounting (Motor on Side):}

\emph{Advantages:} 1. Compact vertical envelope 2. Motor remains
stationary (longer life) 3. Fixed electrical connections 4. Easier motor
replacement

\emph{Disadvantages:} 1. Requires right-angle gearbox or belt/pulley
transmission 2. Additional gear backlash (if gearbox used) 3. More
complex mechanical design 4. Chip exposure (requires bellows protection)

\textbf{Heat Dissipation Analysis:}

Motor heat generation at 70\% duty cycle: {[}Q = (P\_\{input\} -
P\_\{output\}) \times duty{]}

For 400W rated motor at 70\% efficiency: {[}Q = (400 - 280) \times 0.7 =
84 \text{ W continuous}{]}

This heat must be dissipated to prevent: 1. Thermal expansion of motor
housing 2. Demagnetization of permanent magnets (\textgreater120°C risk)
3. Thermal gradient in column structure

\textbf{Cooling strategies:} - Forced air cooling: 10-20 W/(m\^{}2*K)
heat transfer coefficient - Required surface area: (A = Q / (h
\times \Delta T) = 84 / (15 \times 40) = 0.14) m\^{}2 - Finned motor
housing or attached heat sink recommended

\textbf{Recommended Configuration:}

For Z-axis travel \textless{} 250 mm and spindle mass \textless{} 5 kg:
- \textbf{Vertical mounting with sealed AC servo motor} - Use motor
brake to prevent gravity drop during power loss - IP65 rated motor
minimum - Aluminum heat sink attached to motor frame

For Z-axis travel \textgreater{} 250 mm or spindle mass \textgreater{} 5
kg: - \textbf{Horizontal mounting with right-angle gearbox} - Ratio 3:1
to 5:1 for inertia matching - Backlash-free gearbox (planetary or
harmonic drive) - Motor remains stationary (extended life)

\subsubsection{2.5 Counterbalance Systems: Gravitational Force
Compensation}\label{counterbalance-systems-gravitational-force-compensation}

\textbf{Fundamental Requirement:}

The counterbalance system must provide upward force equal to the weight
of all moving components, reducing the motor's workload to only
overcoming friction and accelerating mass (without the gravitational
bias).

\textbf{Net Force Analysis (Unbalanced System):}

{[}F\_\{motor\} = ma + F\_\{friction\} + mg \cdot sgn(direction){]}

For upward motion: Motor must lift weight AND accelerate For downward
motion: Motor must brake against falling weight

This asymmetry causes: 1. Servo tuning challenges (different gains
needed for up vs.~down) 2. Higher RMS motor current (thermal stress) 3.
Asymmetric acceleration profiles 4. Increased bearing wear

\textbf{Net Force Analysis (Balanced System):}

{[}F\_\{motor\} = ma + F\_\{friction\}{]}

Gravitational term eliminated! Symmetric performance in both directions.

\textbf{Counterbalance Technologies:}

\textbf{1. Gas Springs:}

Nitrogen-charged gas springs provide nearly constant force over moderate
strokes.

Force characteristic: {[}F(x) = F\_0
\left(\frac{V_0}{V_0 - A x}\right)\^{}n{]}

Where: - (F\_0) = initial force - (V\_0) = initial gas volume - (A) =
piston area - (x) = extension distance - (n) = polytropic exponent
(approx 1.2 for nitrogen)

For small displacements ((x \textless\textless{} V\_0/A)), force is
approximately constant.

\emph{Advantages:} - Compact size - No external power - Adjustable force
(via pressure) - Long life (\textgreater1 million cycles)

\emph{Disadvantages:} - Force varies slightly with temperature - Limited
stroke (typically 0.5x compressed length) - Progressive force
characteristic

\textbf{Design Example:}

Moving mass: 6 kg -\textgreater{} Required force: 58.9 N

Select: Two 300 N gas springs in parallel - Initial compression: 80\% of
stroke - Operating range: 70-90\% compression - Force variation over
range: +/-5\%

\textbf{2. Counterweight System:}

Steel weights connected via cable and pulley provide constant
gravitational counterforce.

{[}F\_\{counterweight\} = m\_\{weight\} \times g{]}

\emph{Advantages:} - Perfectly constant force - No wear or degradation -
Temperature insensitive - Simple and reliable

\emph{Disadvantages:} - Requires vertical space for weight travel -
Cable maintenance (lubrication, inspection) - Pulley bearing friction -
More complex mechanical design

\textbf{Design Consideration:}

Cable routing over pulleys introduces friction: {[}F\_\{effective\} =
F\_\{weight\} \times \eta\emph{\{pulley\}\^{}\{n}\{pulleys\}\}{]}

For (\eta\emph{\{pulley\} = 0.98) (ball bearing pulley), 2 pulleys:
{[}F}\{effective\} = F\_\{weight\} \times 0.98\^{}2 = 0.96
F\_\{weight\}{]}

4\% force loss -\textgreater{} Design for 104\% of required
counterweight mass.

\textbf{3. Constant-Force Springs:}

Tightly coiled springs that unroll to provide constant force.

\emph{Advantages:} - Compact - No friction losses - Long stroke
capability

\emph{Disadvantages:} - Lower force capacity (typically \textless{} 50
N) - Limited to light-duty applications - Lateral space requirement for
spring storage drum

\textbf{Tuning Procedure:}

\begin{enumerate}
\def\labelenumi{\arabic{enumi}.}
\tightlist
\item
  Remove all external loads from Z-axis
\item
  Disable servo drive (free-wheeling mode)
\item
  Measure force required to move axis at constant velocity

  \begin{itemize}
  \tightlist
  \item
    Upward force: (F\_\{up\})
  \item
    Downward force: (F\_\{down\})
  \end{itemize}
\item
  Calculate ideal counterbalance force: {[}F\_\{ideal\} =
  \frac{F_{up} - F_{down}}{2}{]}
\item
  Adjust counterbalance until (F\_\{up\} \approx F\_\{down\}) (within
  +/-10\%)
\end{enumerate}

\textbf{Verification:} - Command slow velocity profile (50 mm/min) up
and down - Monitor motor current - Target: \textless10\% current
variation between directions

\subsubsection{2.6 Complete Architecture Summary
Table}\label{complete-architecture-summary-table}

The following table provides comprehensive specifications for a
precision vertical axis suitable for plasma cutting, laser cutting, and
light milling applications:

{\def\LTcaptype{} % do not increment counter
\begin{longtable}[]{@{}
  >{\raggedright\arraybackslash}p{(\linewidth - 6\tabcolsep) * \real{0.2000}}
  >{\raggedright\arraybackslash}p{(\linewidth - 6\tabcolsep) * \real{0.2727}}
  >{\raggedright\arraybackslash}p{(\linewidth - 6\tabcolsep) * \real{0.4000}}
  >{\raggedright\arraybackslash}p{(\linewidth - 6\tabcolsep) * \real{0.1273}}@{}}
\toprule\noalign{}
\begin{minipage}[b]{\linewidth}\raggedright
Component
\end{minipage} & \begin{minipage}[b]{\linewidth}\raggedright
Specification
\end{minipage} & \begin{minipage}[b]{\linewidth}\raggedright
Engineering Rationale
\end{minipage} & \begin{minipage}[b]{\linewidth}\raggedright
Notes
\end{minipage} \\
\midrule\noalign{}
\endhead
\bottomrule\noalign{}
\endlastfoot
\textbf{Travel Range} & 150-250 mm & Covers typical plasma plate heights
with margin & Increase for thicker materials \\
\textbf{Drive Type} & Ball-screw Ø16-25 mm & High stiffness, zero
backlash, precision & Ø20 mm most common for medium loads \\
\textbf{Screw Pitch} & 5-10 mm & Balances speed and force capability & 5
mm for precision, 10 mm for speed \\
\textbf{Critical Speed Margin} & \textgreater= 30\% margin; operate
\textless= 80\% of \$n\_\{cr\}\$ & Avoids screw whipping and excessive
vibration & Verify with end‑fixity factor and free length \\
\textbf{Screw Grade} & C5 or better & Position accuracy \textless= 0.023
mm/300 mm & C7 acceptable for non-precision work \\
\textbf{Guide Type} & Two linear rails & Distributed load, moment
resistance & Single rail = inadequate moment capacity \\
\textbf{Rail Size} & 15-20 mm width & Load capacity 3-7 kN per rail pair
& Scale with cutting forces \\
\textbf{Rail Spacing} & 80-120 mm & Maximizes moment arm for stiffness &
Wider spacing = lower angular deflection \\
\textbf{Rail Preload} & Medium (ZA) or Heavy (ZB) & Eliminates play,
increases stiffness & ZA for general, ZB for heavy cutting \\
\textbf{Motor Orientation} & Vertical (above) or Horizontal (side) &
Vertical for chip protection, Horizontal for compactness & Trade-off
between protection and height \\
\textbf{Motor Type} & AC Servo, 400-750W & Sufficient torque for
acceleration + cutting forces & Include brake for gravity holding \\
\textbf{Motor Brake} & Electromechanical, 2-5 Nm & Prevents drop during
power loss & Critical safety feature \\
\textbf{Gearbox} & 3:1 to 5:1 planetary (if used) & Inertia matching,
torque multiplication & Backlash-free design required \\
\textbf{Counterbalance Type} & Gas spring or counterweight & Eliminates
gravitational bias & Gas spring for compactness \\
\textbf{Counterbalance Force} & 100-110\% of moving mass weight & Slight
overbalance compensates friction & Adjustable design preferred \\
\textbf{Column Cross-Section} & 100x100 to 150x150 mm RHS & High moment
of inertia, low deflection & Steel or aluminum with internal ribs \\
\end{longtable}
}

\subsubsection{2.7 Servo Integration and Torque Sizing (with
Efficiency)}\label{servo-integration-and-torque-sizing-with-efficiency}

The servo must overcome gravity, friction, and commanded acceleration.
Including screw efficiency (\eta) and lead (L\_\{\text{lead}\}): {[}
T\_\{\text{req}\} = \frac{L_{\text{lead}}}{2\pi\,\eta}\Big( m,a +
F\_\{\text{friction}\} + m g \Big) {]} Reflected linear inertia at the
motor shaft is {[} J\_\{\text{ref}\} =
m\left(\frac{L_{\text{lead}}}{2\pi}\right)\^{}\{2\},
\qquad J\_\{\text{eq}\} = J\_m + J\_\{\text{ref}\} +
J\_\{\text{coupling}\} + J\_\{\text{screw}\} {]} Maintain
(J\_\{\text{ref}\}/J\_m) between 1:1 and 5:1 for robust tuning.

\textbf{Worked Example - 200 mm Z‑Axis:} Given (m=8) kg,
(L\_\{\text{lead}\}=0.005) m, (a=2.0,g), (F\_\{\text{friction}\}=10) N,
(\eta=0.92): {[} T\_\{\text{req}\} =
\frac{0.005}{2\pi\cdot 0.92}\Big(8\cdot 2\cdot 9.81 + 10 +
8\cdot 9.81\Big) \approx 0.46,\text{N*m} {]} With a 3:1 gearbox, motor
torque reduces to \textasciitilde0.15 N*m (at 3x reflected inertia).
Verify thermal limits and add 20-30\% margin.

Feedforward improves tracking: {[} T\_\{\text{ff}\} =
J\_\{\text{eq}\},\frac{2\pi}{L_{\text{lead}}},\ddot{z} +
\frac{L_{\text{lead}}}{2\pi\,\eta},F\_\{\text{friction}\}(\dot{z}) {]}
Tune gains while respecting the natural frequency target
(\textgreater=5x servo bandwidth).

\subsubsection{2.8 Brake Sizing and Drop
Safety}\label{brake-sizing-and-drop-safety}

The holding brake must prevent gravity‑induced motion during power loss:
{[} T\_\{\text{brake}\} \ge \frac{L_{\text{lead}}}{2\pi\,\eta},m g
\times S\_f {]} with safety factor (S\_f) (typically 1.5-2.5). For (m=8)
kg, (L\_\{\text{lead}\}=5) mm, (\eta=0.92), (S\_f=2): {[}
T\_\{\text{brake}\} \ge \frac{0.005}{2\pi\cdot 0.92} (8\cdot 9.81)
\times 2 \approx 0.13,\text{N*m} {]} Select the next larger standard
brake (e.g., 2.5 N*m) to account for wear, contamination, and incline.
Add mechanical hard‑stops and soft‑limits; verify weekly during
maintenance.

\paragraph{2.8.1 Brake Verification Test
(Commissioning)}\label{brake-verification-test-commissioning}

Purpose: Confirm the holding brake prevents gravity‑induced motion and
that no hazardous sag occurs during power loss.

Procedure: - Jog Z to mid‑stroke; stop motion and engage brake (power
ON). - Disable servo drive power to simulate power loss; observe
position with a 0.01 mm dial indicator. - Re‑enable power; lift 2-3 mm;
repeat test 3x to check repeatability.

Acceptance Criteria: - Sag/displacement \textless= 0.05 mm within 10 s
after power removal. - No sustained downward drift (\textgreater{} 0.05
mm/min) while brake engaged. - On re‑energizing, inrush current spike is
transient; no prolonged overcurrent alarm.

If Fails: - Increase brake size or adjust brake airgap per manufacturer
spec. - Reduce counterbalance over‑ or under‑bias to minimize net
gravity load. - Inspect screw/nut friction (contamination) and verify
brake wiring/voltage. \textbar{} \textbf{Column Material} \textbar{}
Steel or Aluminum \textbar{} Steel for stiffness, Aluminum for mass
reduction \textbar{} Consider thermal expansion \textbar{} \textbar{}
\textbf{Natural Frequency} \textbar{} \textgreater= 150 Hz \textbar{}
Ensure \textgreater{} 5x servo bandwidth \textbar{} Verify with FEA and
accelerometer testing \textbar{} \textbar{} \textbf{Cable Carrier}
\textbar{} Enclosed drag chain \textbar{} Protects cables, manages
flexing \textbar{} Size for \textgreater= 2x cable bend radius
\textbar{} \textbar{} \textbf{Position Feedback} \textbar{} Encoder
0.001 mm resolution \textbar{} Closed-loop control, positioning accuracy
\textbar{} Rotary on motor or linear on axis \textbar{} \textbar{}
\textbf{Limit Switches} \textbar{} Mechanical + inductive proximity
\textbar{} Hard limits (crash protection) + soft limits (control)
\textbar{} Dual redundancy for safety \textbar{}

\textbf{Design Integration Notes:}

\begin{enumerate}
\def\labelenumi{\arabic{enumi}.}
\tightlist
\item
  \textbf{Electrical Integration:}

  \begin{itemize}
  \tightlist
  \item
    Motor power cable: Shielded, oil-resistant, 1.5x rated current
    capacity
  \item
    Encoder cable: Twisted-pair, shielded, separate conduit from power
  \item
    Limit switch wiring: 24V DC logic, optically isolated inputs
  \item
    Emergency stop loop: Hardwired series circuit, NC contacts
  \end{itemize}
\item
  \textbf{Lubrication System:}

  \begin{itemize}
  \tightlist
  \item
    Ball-screw: Automatic grease injector every 8 hours operation
  \item
    Linear rails: Manual regreasing every 6 months (via zerk fittings)
  \item
    Grease type: NLGI Grade 2, lithium-based, -20°C to +120°C
  \end{itemize}
\item
  \textbf{Thermal Management:}

  \begin{itemize}
  \tightlist
  \item
    Motor heat sink: Natural convection or forced air
  \item
    Column insulation: Separate motor thermal mass from structure
  \item
    Temperature monitoring: Thermistor on motor winding, shutdown at
    130°C
  \end{itemize}
\item
  \textbf{Safety Systems:}

  \begin{itemize}
  \tightlist
  \item
    Motor brake: Engages on power loss, prevents gravity drop
  \item
    Counterbalance: Sized to prevent free-fall (even if brake fails)
  \item
    Bellows/covers: Protect ball-screw from chip contamination
  \item
    Software limits: Prevent over-travel before hitting hard stops
  \end{itemize}
\end{enumerate}

This comprehensive architecture provides a robust foundation for
precision vertical axis performance across diverse CNC applications.

\begin{center}\rule{0.5\linewidth}{0.5pt}\end{center}

\begin{center}\rule{0.5\linewidth}{0.5pt}\end{center}

\subsection{References}\label{references-1}

\begin{enumerate}
\def\labelenumi{\arabic{enumi}.}
\tightlist
\item
  \textbf{ISO 14728-1:2017} - Rolling bearings - Linear motion rolling
  bearings - Dynamic load ratings
\item
  \textbf{THK Linear Motion Systems Catalog} - Z-axis configurations and
  mounting
\item
  \textbf{Hiwin Linear Guideway Technical Manual} - Vertical axis
  applications
\item
  \textbf{NSK Linear Guides CAT. No.~E728g} - Profile rail systems for
  vertical mounting
\item
  \textbf{Weck, M. \& Brecher, C. (2006).} \emph{Machine Tools 1}.
  Springer
\item
  \textbf{SKF Linear Motion \& Actuation} - Vertical axis design
  guidelines
\end{enumerate}

\newpage

\section{Module 2 - Vertical Axis and Column
Assembly}\label{module-2---vertical-axis-and-column-assembly-2}

\subsection{3. Core Equations for Column
Behaviour}\label{core-equations-for-column-behaviour}

Understanding the mathematical relationships governing vertical axis
structural behavior is essential for rational design decisions. This
section derives and explains the fundamental equations that describe
deflection, resonance, critical speeds, and servo stability
requirements.

\subsubsection{3.1 Cantilever Beam Deflection: The Foundation of
Structural
Design}\label{cantilever-beam-deflection-the-foundation-of-structural-design}

\textbf{Physical Situation:}

The Z-axis column functions as a cantilever beam with: - Fixed support
at the base (machine bed or gantry) - Free end carrying the moving
carriage and tooling - Concentrated load at the free end from cutting
forces

\textbf{Fundamental Equation:}

For a cantilever beam with uniform cross-section, subjected to point
load (F) at distance (L) from the fixed support:

{[}\delta\_B = \frac{F L^3}{3 E I}{]}

\textbf{Derivation from Beam Theory:}

Starting with the moment-curvature relationship: {[}\frac{d^2 y}{dx^2} =
\frac{M(x)}{E I}{]}

For cantilever with tip load (F), moment at distance (x) from free end:
{[}M(x) = -F x{]}

Integrating once for slope: {[}\frac{dy}{dx} = \int \frac{-F x}{E I} dx
= \frac{-F x^2}{2 E I} + C\_1{]}

Boundary condition at fixed end ((x=L)): (dy/dx = 0) {[}C\_1 =
\frac{F L^2}{2 E I}{]}

Slope equation: {[}\frac{dy}{dx} = \frac{F}{2 E I}(L\^{}2 - x\^{}2){]}

Integrating for deflection: {[}y = \int \frac{F}{2 E I}(L\^{}2 - x\^{}2)
dx = \frac{F}{2 E I}\left(L\^{}2 x - \frac{x^3}{3}\right) + C\_2{]}

Boundary condition at fixed end ((x=L)): (y=0) {[}C\_2 =
-\frac{F}{2 E I}\left(L\^{}3 - \frac{L^3}{3}\right) =
-\frac{F L^3}{3 E I}{]}

Deflection at free end ((x=0)): {[}\delta\_B = y(0) =
-\frac{F L^3}{3 E I}{]}

The negative sign indicates downward deflection; taking absolute value:

{[}\boxed{\delta_B = \frac{F L^3}{3 E I}}{]}

\textbf{Design Form - Required Moment of Inertia:}

Rearranging to solve for the required structural property:

{[}I\_\{req\} = \frac{F L^3}{3 E \delta_{max}}{]}

This is the primary design equation for column sizing.

\textbf{Worked Example 1: Plasma Cutting Column}

\textbf{Requirements:} - Cantilever length: (L = 350) mm (from column
face to torch tip) - Maximum cutting force: (F = 300) N (conservative
estimate) - Material: Steel tube ((E = 200) GPa) - Allowable deflection:
(\delta\_\{max\} = 0.025) mm (25 mum)

\textbf{Calculate required moment of inertia:}

{[}I\_\{req\} =
\frac{300 \times 0.35^3}{3 \times 200 \times 10^9 \times 0.000025}{]}

{[}I\_\{req\} = \frac{300 \times 0.042875}{15 \times 10^6} =
\frac{12.8625}{15 \times 10^6}{]}

{[}I\_\{req\} = 8.58 \times 10\^{}\{-7\} \text{ m}\^{}4 = 858,000
\text{ mm}\^{}4{]}

\textbf{Select appropriate section:}

For square RHS (hollow rectangular section), approximate formula for
thin-walled section:

{[}I = \frac{b h^3}{12} - \frac{(b-2t)(h-2t)^3}{12}{]}

For square section ((b = h)): {[}I = \frac{b^4 - (b-2t)^4}{12}{]}

\textbf{Trial section: 100x100x6 mm}

{[}I = \frac{100^4 - 88^4}{12} = \frac{100,000,000 - 59,969,536}{12}{]}

{[}I = \frac{40,030,464}{12} = 3,335,872 \text{ mm}\^{}4{]}

\textbf{Verification:} {[}I\_\{actual\} = 3.34 \times 10\^{}6
\text{ mm}\^{}4 \textgreater{} I\_\{req\} = 8.58 \times 10\^{}5
\text{ mm}\^{}4{]}

\textbf{Safety factor:} (SF = 3.34/0.858 = 3.9) {[}check{]} (adequate)

\textbf{Calculated deflection with selected section:} {[}\delta =
\frac{300 \times 0.35^3}{3 \times 200 \times 10^9 \times 3.34 \times 10^{-6}}{]}

{[}\delta = \frac{12.86}{2.004 \times 10^6} = 6.4 \times 10\^{}\{-6\}
\text{ m} = 0.0064 \text{ mm} = 6.4 \text{ mum}{]}

\textbf{Result:} Actual deflection is 6.4 mum, well below 25 mum limit.
The safety factor of 3.9x provides margin for: - Dynamic loading
(impacts, rapid accelerations) - Assembly tolerances - Material property
variations - Thermal effects

\textbf{Worked Example 2: Milling Machine Column}

\textbf{Requirements:} - Cantilever length: (L = 400) mm - Maximum
cutting force: (F = 800) N (heavy milling) - Material: Cast iron ((E =
120) GPa, higher damping than steel) - Allowable deflection:
(\delta\_\{max\} = 0.015) mm (15 mum, tighter tolerance)

\textbf{Calculate required moment of inertia:}

{[}I\_\{req\} =
\frac{800 \times 0.4^3}{3 \times 120 \times 10^9 \times 0.000015}{]}

{[}I\_\{req\} = \frac{51.2}{5.4 \times 10^6} = 9.48 \times 10\^{}\{-6\}
\text{ m}\^{}4 = 9.48 \times 10\^{}6 \text{ mm}\^{}4{]}

This requires a much larger section due to: 1. Higher cutting forces
(800 N vs 300 N) 2. Longer cantilever (400 mm vs 350 mm) 3. Tighter
deflection specification (15 mum vs 25 mum) 4. Lower material stiffness
(cast iron vs steel)

\textbf{Selected section: 150x150x10 mm steel RHS with internal ribs}

{[}I\_\{base\} = \frac{150^4 - 130^4}{12} = 15.2 \times 10\^{}6
\text{ mm}\^{}4{]}

Internal ribbing adds approximately 25\% to effective I: {[}I\_\{total\}
\approx 19 \times 10\^{}6 \text{ mm}\^{}4{]}

\textbf{Safety factor:} (SF = 19/9.48 = 2.0) {[}check{]} (acceptable for
industrial application)

\subsubsection{3.2 Natural (Resonant) Frequency: Dynamic Stability
Analysis}\label{natural-resonant-frequency-dynamic-stability-analysis}

\textbf{Physical Principle:}

Every structural system has natural frequencies at which it will vibrate
with minimal damping. If external excitation (cutting forces, servo
oscillations) occurs near a natural frequency, resonance amplifies
vibrations, destroying positioning accuracy and potentially causing
structural failure.

\textbf{Single-Degree-of-Freedom Model:}

The Z-axis column with moving carriage can be modeled as a spring-mass
system:

{[}f\_n = \frac{1}{2\pi} \sqrt{\frac{k}{m}}{]}

Where: - (k) = structural stiffness (N/m) - (m) = effective mass (kg)

\textbf{Cantilever Stiffness:}

From deflection equation (\delta = FL\^{}3/(3EI)), the spring constant
is:

{[}k = \frac{F}{\delta} = \frac{3 E I}{L^3}{]}

\textbf{Natural Frequency of Cantilever with Tip Mass:}

{[}f\_n = \frac{1}{2\pi} \sqrt{\frac{3 E I}{m L^3}}{]}

\textbf{Effective Mass Correction:}

For a cantilever with distributed mass (m\_\{beam\}) and tip mass
(m\_\{tip\}), the effective mass is:

{[}m\_\{eff\} = m\_\{tip\} + \alpha m\_\{beam\}{]}

Where (\alpha \approx 0.25) for first-mode cantilever vibration (proven
through modal analysis).

\textbf{Worked Example 3: Natural Frequency Calculation}

\textbf{Given:} - Column: 120x120x8 mm steel RHS - Length: (L = 450) mm
- Moving mass (carriage + spindle): (m\_\{tip\} = 7) kg - Column mass:
(m\_\{beam\} = 2.5) kg/m x 0.45 m = 1.125 kg - (E = 200) GPa - (I = 3.85
\times 10\^{}\{-6\}) m\^{}4

\textbf{Calculate effective mass:} {[}m\_\{eff\} = 7 + 0.25 \times 1.125
= 7.28 \text{ kg}{]}

\textbf{Calculate stiffness:} {[}k =
\frac{3 \times 200 \times 10^9 \times 3.85 \times 10^{-6}}{0.45^3}{]}

{[}k = \frac{2.31 \times 10^6}{0.091125} = 25.35 \times 10\^{}6
\text{ N/m}{]}

\textbf{Calculate natural frequency:} {[}f\_n = \frac{1}{2\pi}
\sqrt{\frac{25.35 \times 10^6}{7.28}}{]}

{[}f\_n = \frac{1}{6.283} \sqrt{3.482 \times 10^6} = \frac{1866}{6.283}
= 297 \text{ Hz}{]}

\textbf{Result:} First natural frequency is 297 Hz.

\textbf{Resonance Implications:}

If servo bandwidth is 30 Hz: - Frequency ratio: (297 / 30 = 9.9) -
\textbf{Status:} {[}check{]} Exceeds minimum 5x requirement (excellent
design)

If servo bandwidth is 50 Hz: - Frequency ratio: (297 / 50 = 5.9) -
\textbf{Status:} {[}check{]} Meets minimum 5x requirement (acceptable)

If servo bandwidth is 70 Hz: - Frequency ratio: (297 / 70 = 4.2) -
\textbf{Status:} ✗ Below minimum 5x requirement (resonance risk)

\textbf{Design Actions for Low-Frequency Structures:}

If calculated (f\_n) is insufficient:

\begin{enumerate}
\def\labelenumi{\arabic{enumi}.}
\tightlist
\item
  \textbf{Increase structural stiffness:}

  \begin{itemize}
  \tightlist
  \item
    Larger cross-section (I increases with h\^{}3)
  \item
    Add internal ribs/stiffeners
  \item
    Use stiffer material (steel over aluminum)
  \end{itemize}
\item
  \textbf{Reduce moving mass:}

  \begin{itemize}
  \tightlist
  \item
    Lighter spindle/tool assembly
  \item
    Aluminum carriage components
  \item
    Hollow-core designs
  \end{itemize}
\item
  \textbf{Reduce cantilever length:}

  \begin{itemize}
  \tightlist
  \item
    Minimize tool offset from column
  \item
    Compact carriage design
  \end{itemize}
\item
  \textbf{Reduce servo bandwidth:}

  \begin{itemize}
  \tightlist
  \item
    Trade response speed for stability
  \item
    Typical: Reduce from 50 Hz to 30 Hz
  \end{itemize}
\end{enumerate}

\textbf{Frequency vs.~Stiffness Sensitivity:}

{[}\frac{df_n}{f_n} = \frac{1}{2}\frac{dk}{k} -
\frac{1}{2}\frac{dm}{m}{]}

A 20\% increase in stiffness -\textgreater{} 9.5\% increase in natural
frequency A 20\% reduction in mass -\textgreater{} 9.5\% increase in
natural frequency

\textbf{Priority:} Stiffness increases are more effective than mass
reductions because: - Stiffness scales with I (can increase dramatically
with section size) - Mass reduction limited by functional requirements

\subsubsection{3.3 Ball-Screw Critical Speed: Rotational Stability
Limit}\label{ball-screw-critical-speed-rotational-stability-limit}

\textbf{Physical Phenomenon:}

A rotating shaft (ball-screw) becomes dynamically unstable when
rotational speed reaches the critical speed, at which lateral vibrations
resonate with shaft rotation. This creates catastrophic whirling motion
that can destroy bearings and damage the screw.

\textbf{Critical Speed Equation:}

{[}n\_\{cr\} = \frac{k_f \times 10^6 \times d_r}{L^2}{]}

Where: - (n\_\{cr\}) = critical rotational speed (rpm) - (k\_f) =
end-fixity factor (dimensionless) - (d\_r) = root diameter of screw (mm)
- (L) = unsupported length between bearings (mm)

\textbf{End-Fixity Factors:}

{\def\LTcaptype{} % do not increment counter
\begin{longtable}[]{@{}
  >{\raggedright\arraybackslash}p{(\linewidth - 4\tabcolsep) * \real{0.4872}}
  >{\raggedright\arraybackslash}p{(\linewidth - 4\tabcolsep) * \real{0.1795}}
  >{\raggedright\arraybackslash}p{(\linewidth - 4\tabcolsep) * \real{0.3333}}@{}}
\toprule\noalign{}
\begin{minipage}[b]{\linewidth}\raggedright
Support Condition
\end{minipage} & \begin{minipage}[b]{\linewidth}\raggedright
\$k\_f\$
\end{minipage} & \begin{minipage}[b]{\linewidth}\raggedright
Description
\end{minipage} \\
\midrule\noalign{}
\endhead
\bottomrule\noalign{}
\endlastfoot
Simply supported - simply supported & 3.14 & Both ends on radial
bearings, free to rotate \\
Fixed - simply supported & 3.56 & One end clamped, one end on bearing \\
Fixed - fixed & 4.73 & Both ends rigidly clamped (rare) \\
Fixed - free (cantilever) & 0.56 & One end fixed, other end free
(avoid!) \\
\end{longtable}
}

\textbf{Common Configuration:} Fixed-supported (angular contact bearings
at one end, radial bearing at other) {[}k\_f = 3.56{]}

\textbf{Worked Example 4: Critical Speed Analysis}

\textbf{Given:} - Ball-screw: Ø20 mm, 5 mm pitch - Root diameter: (d\_r
= 17.5) mm - Unsupported length: (L = 500) mm - Support: Fixed-supported
((k\_f = 3.56))

\textbf{Calculate critical speed:}

{[}n\_\{cr\} = \frac{3.56 \times 10^6 \times 17.5}{500^2}{]}

{[}n\_\{cr\} = \frac{62.3 \times 10^6}{250,000} = 249 \text{ rpm}{]}

\textbf{Maximum linear velocity:}

{[}v\_\{max\} = n\_\{cr\} \times p = 249 \times 0.005 = 1.245
\text{ m/min} = 20.8 \text{ mm/s}{]}

\textbf{Safety Factor Application:}

Industry practice: Operate at \textless= 80\% of critical speed

{[}n\_\{operating\} = 0.8 \times 249 = 199 \text{ rpm}{]}

{[}v\_\{operating\} = 199 \times 0.005 = 0.995 \text{ m/min} \approx 1.0
\text{ m/min}{]}

\textbf{Design Decision:} For applications requiring (v \textgreater{}
1) m/min, this screw configuration is inadequate.

\textbf{Design Solutions for Higher Speeds:}

\begin{enumerate}
\def\labelenumi{\arabic{enumi}.}
\item
  \textbf{Reduce unsupported length:} Add center bearing support to
  create two shorter spans.

  Example: 500 mm span -\textgreater{} two 250 mm spans

  {[}n\_\{cr,new\} = \frac{3.56 \times 10^6 \times 17.5}{250^2} = 998
  \text{ rpm}{]}

  {[}v\_\{max,new\} = 998 \times 0.005 = 4.99 \text{ m/min}{]}

  \textbf{Result:} 4x speed increase with center support bearing!
\item
  \textbf{Increase screw diameter:} Use Ø25 mm screw ((d\_r = 22) mm)

  {[}n\_\{cr\} = \frac{3.56 \times 10^6 \times 22}{500^2} = 314
  \text{ rpm}{]}

  26\% speed increase, but increased cost and mass.
\item
  \textbf{Use preloaded screw:} Preloaded screws have effectively higher
  (k\_f) (approx 4.0 vs 3.56)

  {[}n\_\{cr\} = \frac{4.0 \times 10^6 \times 17.5}{500^2} = 280
  \text{ rpm}{]}

  12\% improvement.
\item
  \textbf{Switch to belt drive:} No critical speed limitation! Practical
  for rapid positioning axes where forces are moderate.
\end{enumerate}

\textbf{Critical Speed vs.~Application:}

{\def\LTcaptype{} % do not increment counter
\begin{longtable}[]{@{}
  >{\raggedright\arraybackslash}p{(\linewidth - 4\tabcolsep) * \real{0.2955}}
  >{\raggedright\arraybackslash}p{(\linewidth - 4\tabcolsep) * \real{0.3409}}
  >{\raggedright\arraybackslash}p{(\linewidth - 4\tabcolsep) * \real{0.3636}}@{}}
\toprule\noalign{}
\begin{minipage}[b]{\linewidth}\raggedright
Application
\end{minipage} & \begin{minipage}[b]{\linewidth}\raggedright
Typical Speed
\end{minipage} & \begin{minipage}[b]{\linewidth}\raggedright
Screw Adequacy
\end{minipage} \\
\midrule\noalign{}
\endhead
\bottomrule\noalign{}
\endlastfoot
Precision plasma THC & 0.5 m/min & {[}check{]} Well within limits \\
Laser focal adjustment & 2-3 m/min & ⚠ May require center support \\
Rapid tool change positioning & 5-10 m/min & ✗ Belt drive recommended \\
Probing cycles & 0.1-0.2 m/min & {[}check{]} No concerns \\
\end{longtable}
}

\subsubsection{3.4 Servo Resonance Criterion: Control-Structure
Interaction}\label{servo-resonance-criterion-control-structure-interaction}

\textbf{Control Theory Principle:}

A servo control system with bandwidth (f\_\{servo\}) generates control
signals with frequency content up to approximately (5
\times f\_\{servo\}). If any structural resonance exists within this
range, the controller will excite the resonance, causing instability.

\textbf{Stability Criterion:}

{[}f\_n \geq 5 \times f\_\{servo\}{]} (minimum requirement)

{[}f\_n \geq 10 \times f\_\{servo\}{]} (preferred for robust control)

\textbf{Servo Bandwidth Definition:}

Servo bandwidth is the frequency at which the closed-loop gain drops to
-3 dB (70.7\% of DC gain). Higher bandwidth = faster response, but
increased sensitivity to resonances.

\textbf{Typical Bandwidth Values:}

{\def\LTcaptype{} % do not increment counter
\begin{longtable}[]{@{}
  >{\raggedright\arraybackslash}p{(\linewidth - 6\tabcolsep) * \real{0.1831}}
  >{\raggedright\arraybackslash}p{(\linewidth - 6\tabcolsep) * \real{0.2254}}
  >{\raggedright\arraybackslash}p{(\linewidth - 6\tabcolsep) * \real{0.2958}}
  >{\raggedright\arraybackslash}p{(\linewidth - 6\tabcolsep) * \real{0.2958}}@{}}
\toprule\noalign{}
\begin{minipage}[b]{\linewidth}\raggedright
Application
\end{minipage} & \begin{minipage}[b]{\linewidth}\raggedright
Servo Bandwidth
\end{minipage} & \begin{minipage}[b]{\linewidth}\raggedright
Required \$f\_n\$ (5x)
\end{minipage} & \begin{minipage}[b]{\linewidth}\raggedright
Required \$f\_n\$ (10x)
\end{minipage} \\
\midrule\noalign{}
\endhead
\bottomrule\noalign{}
\endlastfoot
Positioning (slow) & 10 Hz & 50 Hz & 100 Hz \\
Standard machining & 30 Hz & 150 Hz & 300 Hz \\
High-speed machining & 50 Hz & 250 Hz & 500 Hz \\
Ultra-precision & 20 Hz & 100 Hz & 200 Hz \\
\end{longtable}
}

\textbf{Worked Example 5: Servo Tuning for Structural Resonance}

\textbf{Given:} - Measured natural frequency: (f\_n = 180) Hz (via
accelerometer) - Desired servo bandwidth: (f\_\{servo\} = 40) Hz -
Current configuration: PID controller

\textbf{Check stability criterion:}

{[}\frac{f_n}{f_{servo}} = \frac{180}{40} = 4.5{]}

\textbf{Result:} 4.5x ratio is below minimum 5x requirement (marginal
stability).

\textbf{Design Options:}

\textbf{Option 1: Reduce Servo Bandwidth (Conservative)}

{[}f\_\{servo,max\} = \frac{f_n}{5} = \frac{180}{5} = 36 \text{ Hz}{]}

\begin{itemize}
\tightlist
\item
  Reduce bandwidth from 40 Hz to 36 Hz (10\% reduction)
\item
  Trade-off: Slightly slower response, but guaranteed stability
\item
  Implementation: Reduce proportional gain by 20\%
\end{itemize}

\textbf{Option 2: Add Notch Filter (Sophisticated)}

Design notch filter centered at (f\_n = 180) Hz:

{[}H(s) =
\frac{s^2 + 2\zeta_z \omega_n s + \omega_n^2}{s^2 + 2\zeta_p \omega_n s + \omega_n^2}{]}

Where: - (\omega\_n = 2\pi f\_n = 1131) rad/s - (\zeta\_z = 0.05)
(numerator damping - creates notch) - (\zeta\_p = 0.7) (denominator
damping - limits attenuation width)

This filter attenuates gain by 20-30 dB at 180 Hz while preserving
performance at other frequencies.

\begin{itemize}
\tightlist
\item
  Advantage: Maintains 40 Hz bandwidth
\item
  Disadvantage: Adds phase lag, requires tuning
\end{itemize}

\textbf{Option 3: Increase Structural Stiffness (Fundamental)}

Target: (f\_n = 250) Hz (provides 6.25x ratio at 40 Hz bandwidth)

Required stiffness increase: {[}\frac{f_{n,new}}{f_{n,old}} =
\sqrt{\frac{k_{new}}{k_{old}}}{]}

{[}\frac{250}{180} = 1.39 = \sqrt{\frac{k_{new}}{k_{old}}}{]}

{[}\frac{k_{new}}{k_{old}} = 1.39\^{}2 = 1.93{]}

\textbf{Requirement:} Increase structural stiffness by 93\%

\textbf{Approach:} - Upgrade column from 100x100 to 120x120 mm (+73\% in
I) - Add internal ribs (+20\% effective stiffness) - \textbf{Total:}
1.73 x 1.20 = 2.08x (meets requirement)

\textbf{Recommended Solution:}

For production machines: \textbf{Combination approach} 1. Increase
structural stiffness to (f\_n \geq 200) Hz 2. Add notch filter for
additional margin 3. Set servo bandwidth to 35-40 Hz 4. Verify stability
with step response testing

\textbf{Verification Testing:}

After tuning, perform step response test: 1. Command 10 mm step input 2.
Record position vs.~time with high-resolution encoder 3. Measure: - Rise
time (10\%-90\% of step) - Overshoot (should be \textless{} 10\%) -
Settling time (within +/-0.01 mm) - Ringing frequency (should not match
(f\_n))

\textbf{Acceptance Criteria:} - Overshoot \textless{} 10\% - Settling
time \textless{} 3 x rise time - No sustained oscillations - Position
error \textless{} 0.02 mm after settling

\begin{center}\rule{0.5\linewidth}{0.5pt}\end{center}

\subsection{3.5 Equation Integration: Holistic Design
Process}\label{equation-integration-holistic-design-process}

The four core equations work together in the design process:

\begin{verbatim}
1. DEFLECTION -> Size column cross-section for stiffness
              ↓
2. NATURAL FREQUENCY -> Verify dynamic stability
              ↓
3. CRITICAL SPEED -> Confirm ball-screw adequacy
              ↓
4. SERVO CRITERION -> Validate control system compatibility
\end{verbatim}

\textbf{Design Iteration Loop:}

\begin{verbatim}
Start with requirements (travel, forces, speed)
    ↓
Calculate required I (deflection equation)
    ↓
Select trial column section
    ↓
Calculate natural frequency
    ↓
    Check: f_n >= 5 x f_servo?
    ↓               ↓
   NO              YES
    ↓               ↓
Increase I    Calculate screw critical speed
    ↓               ↓
Iterate         Check: v_req <= 0.8 x v_critical?
                    ↓               ↓
                   NO              YES
                    ↓               ↓
           Add support bearing   DESIGN COMPLETE
           or use belt drive
\end{verbatim}

All four equations must be satisfied simultaneously for a successful
vertical axis design. This integrated approach ensures structural,
dynamic, and control requirements are met without costly post-build
modifications.

\begin{center}\rule{0.5\linewidth}{0.5pt}\end{center}

\begin{center}\rule{0.5\linewidth}{0.5pt}\end{center}

\subsection{References}\label{references-2}

\begin{enumerate}
\def\labelenumi{\arabic{enumi}.}
\tightlist
\item
  \textbf{Young, W.C. \& Budynas, R.G. (2011).} \emph{Roark's Formulas
  for Stress and Strain} (8th ed.). McGraw-Hill
\item
  \textbf{Gere, J.M. \& Timoshenko, S.P. (2012).} \emph{Mechanics of
  Materials} (8th ed.). Cengage Learning
\item
  \textbf{Hibbeler, R.C. (2017).} \emph{Structural Analysis} (10th ed.).
  Pearson
\item
  \textbf{Beer, F.P. et al.~(2015).} \emph{Mechanics of Materials} (7th
  ed.). McGraw-Hill
\item
  \textbf{ISO 11352:2012} - Cranes - Loading and stability calculations
\item
  \textbf{AISC Steel Construction Manual} (15th ed., 2017)
\end{enumerate}

\newpage

\section{Module 2 - Vertical Axis \&
Z-Stage}\label{module-2---vertical-axis-z-stage-3}

\subsection{Overview}\label{overview-2}

Linear guides for vertical axes must address unique challenges beyond
horizontal applications: gravitational preload, unequal loading on
carriages, potential for binding under moment loads, and safety
requirements for power-loss conditions. This section covers selection,
sizing, and installation of linear guides optimized for Z-axis
performance.

\subsection{Gravitational Loading
Effects}\label{gravitational-loading-effects}

\subsubsection{Constant Preload from
Weight}\label{constant-preload-from-weight}

Vertical guides experience continuous loading from moving mass weight:

\textbf{Load per carriage (4 carriages typical):}

{[}F\_\{carriage\} = \frac{m \times g}{n}{]}

Where: - m = moving mass (kg) - g = 9.81 m/s\^{}2 - n = number of
carriages (typically 4)

\textbf{Example:} - Moving mass: 12 kg - Load per carriage: 12 x 9.81 /
4 = 29.4 N

\textbf{Implications:} - Minimum preload must exceed gravitational load
- Lower carriages carry slightly more load (column deflection) - Preload
selection more critical than horizontal axes

\subsubsection{Moment Loading}\label{moment-loading}

Off-center masses create moment loads:

{[}M = F\_\{tool\} \times d\_\{offset\}{]}

Moment distributes unequally across carriage pairs, potentially
overloading one side.

\subsection{Rail Selection for Vertical
Axes}\label{rail-selection-for-vertical-axes}

\subsubsection{Size Selection}\label{size-selection}

\textbf{Dynamic load capacity must account for:} 1. Gravitational
constant load 2. Acceleration forces 3. Cutting forces 4. Safety factor
(2-3x for vertical)

\textbf{Basic load rating:}

{[}C\_\{required\} = P
\times \left(\frac{L_{design}}{L_{rating}}\right)\^{}\{1/3\}{]}

Vertical applications typically require one size larger than horizontal
equivalent.

\subsubsection{Preload Selection}\label{preload-selection}

\textbf{Preload classes (HIWIN nomenclature):} - \textbf{Z0}: Light
preload (smallest clearance) - \textbf{ZA}: Light-medium preload -
\textbf{ZB}: Medium preload (most common for vertical) - \textbf{ZC}:
Heavy preload (high-precision, rigid mounting)

\textbf{Recommendation for Z-axis:} - Small machines (\textless10kg
moving): ZA or ZB - Medium machines (10-25kg): ZB or ZC - Large machines
(\textgreater25kg): ZC

Higher preload provides: - Better rigidity under moment loads - Reduced
play and vibration - Increased friction (counterbalance helps overcome)

\subsubsection{Rail Pair Configuration}\label{rail-pair-configuration}

\textbf{Single pair (not recommended):} - Adequate for light loads only
- Poor moment resistance - Risk of binding

\textbf{Dual pairs (standard):} - Two rails, four carriages total - Good
moment capacity - Spacing: 80-150mm typical

\textbf{Multiple pairs (heavy-duty):} - Three+ rails for very heavy
spindles - Distributed loading - Complex alignment requirements

\subsection{Installation and
Alignment}\label{installation-and-alignment}

\subsubsection{Rail Mounting}\label{rail-mounting}

\textbf{Surface preparation:} - Flatness: 10-20 µm over rail length -
Parallelism between rails: \textless20 µm - Surface finish: Ra 1.6 µm

\textbf{Mounting procedure:} 1. Clean mounting surface and rail bottom
2. Apply thin layer of way oil 3. Position rail with dowel pins (if
provided) 4. Tighten mounting screws progressively from center outward
5. Torque to specification (typically 6-8 N\emph{m for M5, 12-15 N}m for
M6)

\subsubsection{Carriage Alignment}\label{carriage-alignment}

\textbf{Vertical orientation:} - All four carriages must align on single
plane - Misalignment causes binding and accelerated wear - Use precision
spacers during assembly

\textbf{Checking alignment:} 1. Install carriages on rails 2. Bolt to
carriage plate with torque 3. Check smooth motion by hand over full
travel 4. Any binding indicates misalignment--shim as needed

\subsection{Safety Considerations}\label{safety-considerations}

\subsubsection{Brake Integration}\label{brake-integration}

Vertical axes require mechanical brakes for safety:

\textbf{Mounting options:} 1. Motor-integrated brake (most common) 2.
Separate brake on ball screw 3. Friction brake on carriage (backup)

\textbf{Brake sizing:} - Must hold 1.5x maximum moving mass - Engage
time: \textless100ms typical - Spring-applied, electrically released
(fail-safe)

\subsubsection{Anti-Drop Mechanisms}\label{anti-drop-mechanisms}

\textbf{Mechanical locks:} - Spring-loaded pawl engages rack on power
loss - Manual release for maintenance - Common on large machines

\textbf{Counterbalance as safety:} - Properly balanced system prevents
rapid descent - Not sufficient as sole safety mechanism - Use with brake
for redundancy

\subsection{Key Takeaways}\label{key-takeaways-2}

\begin{enumerate}
\def\labelenumi{\arabic{enumi}.}
\tightlist
\item
  \textbf{Vertical guides} require higher preload and larger sizes than
  horizontal equivalents
\item
  \textbf{Gravitational loading} is continuous and must be included in
  capacity calculations
\item
  \textbf{Moment loads} from off-center masses require proper carriage
  spacing
\item
  \textbf{Alignment} is critical--poor alignment causes binding and
  premature failure
\item
  \textbf{Mechanical brakes} are mandatory safety devices for vertical
  axes
\item
  \textbf{Dual-pair configuration} (4 carriages) is standard for most
  applications
\item
  \textbf{Preload class ZB or ZC} recommended for vertical mounting
\item
  \textbf{Anti-drop mechanisms} provide additional safety layers
\end{enumerate}

\begin{center}\rule{0.5\linewidth}{0.5pt}\end{center}

\textbf{Next}: \href{section-2.5-ball-screw-vertical.md}{Section 2.5 -
Ball Screws for Vertical Axes}

\textbf{Previous}:
\href{section-2.3-column-structural-design.md}{Section 2.3 - Column
Structural Design}

\newpage

\section{Module 2 - Vertical Axis and Column
Assembly}\label{module-2---vertical-axis-and-column-assembly-3}

\subsection{5. Example Application - Designing a 200 mm Travel
Z-Axis}\label{example-application---designing-a-200-mm-travel-z-axis}

This section presents a complete, worked design example that integrates
all concepts from previous sections into a cohesive vertical axis system
suitable for plasma cutting or laser engraving applications.

\subsubsection{5.1 Design Requirements
Specification}\label{design-requirements-specification}

\textbf{Application:} Precision plasma cutting system
\textbf{Performance Requirements:}

{\def\LTcaptype{} % do not increment counter
\begin{longtable}[]{@{}
  >{\raggedright\arraybackslash}p{(\linewidth - 4\tabcolsep) * \real{0.2973}}
  >{\raggedright\arraybackslash}p{(\linewidth - 4\tabcolsep) * \real{0.4054}}
  >{\raggedright\arraybackslash}p{(\linewidth - 4\tabcolsep) * \real{0.2973}}@{}}
\toprule\noalign{}
\begin{minipage}[b]{\linewidth}\raggedright
Parameter
\end{minipage} & \begin{minipage}[b]{\linewidth}\raggedright
Specification
\end{minipage} & \begin{minipage}[b]{\linewidth}\raggedright
Rationale
\end{minipage} \\
\midrule\noalign{}
\endhead
\bottomrule\noalign{}
\endlastfoot
\textbf{Travel} & 200 mm & Accommodates 25 mm plate + fixturing + safety
margin \\
\textbf{Positioning accuracy} & +/-0.05 mm & Standard plasma cutting
tolerance \\
\textbf{Repeatability} & +/-0.02 mm & Consistent pierce and cut
height \\
\textbf{Maximum cutting force} & 400 N & Conservative torch reaction
force estimate \\
\textbf{Rapid traverse speed} & 3 m/min (50 mm/s) & Rapid repositioning
between cuts \\
\textbf{Cutting feed rate} & 0.5-2 m/min & Typical plasma cutting
speeds \\
\textbf{Torch mass} & 1.5 kg & Including mounting bracket and
consumables \\
\end{longtable}
}

\textbf{Environmental Conditions:} - Operating temperature: 10-40°C -
Humidity: 20-80\% (non-condensing) - Contamination: Metal chips, plasma
spatter, grinding dust - Enclosure: IP54 minimum for column exterior

\subsubsection{5.2 Column Sizing: Cantilever Deflection
Analysis}\label{column-sizing-cantilever-deflection-analysis}

\textbf{Step 1: Estimate Cantilever Length}

Torch centerline distance from column face: - Column width: 120 mm -
Carriage thickness: 30 mm - Torch offset from carriage: 80 mm

{[}L\_\{cantilever\} = \frac{120}{2} + 30 + 80 = 170 \text{ mm}{]}

Add 20\% safety margin for moments: (L\_\{design\} = 1.2 \times 170 =
204) mm approx 200 mm

\textbf{Step 2: Calculate Required Moment of Inertia}

Using deflection formula: {[}I\_\{req\} =
\frac{F L^3}{3 E \delta_{max}}{]}

Given: - (F = 400) N (maximum cutting force) - (L = 0.2) m (200 mm
cantilever) - (E = 200) GPa (steel) - (\delta\_\{max\} = 0.03) mm (30
mum allowable deflection)

{[}I\_\{req\} =
\frac{400 \times 0.2^3}{3 \times 200 \times 10^9 \times 0.00003}{]}

{[}I\_\{req\} = \frac{3.2}{1.8 \times 10^7} = 1.78 \times 10\^{}\{-7\}
\text{ m}\^{}4 = 1.78 \times 10\^{}5 \text{ mm}\^{}4{]}

\textbf{Step 3: Select Column Cross-Section}

\textbf{Trial 1: 100x100x6 mm RHS}

{[}I = \frac{100^4 - 88^4}{12} = 3.34 \times 10\^{}6 \text{ mm}\^{}4{]}

{[}SF = \frac{3.34 \times 10^6}{1.78 \times 10^5} = 18.8{]} (excessive,
column is overdesigned)

\textbf{Trial 2: 80x80x5 mm RHS}

{[}I = \frac{80^4 - 70^4}{12} = 1.68 \times 10\^{}6 \text{ mm}\^{}4{]}

{[}SF = \frac{1.68 \times 10^6}{1.78 \times 10^5} = 9.4{]} (still
overdesigned but more reasonable)

\textbf{Selection Rationale:} While 80x80x5 mm is structurally adequate,
select \textbf{100x100x6 mm} for: 1. Better natural frequency (higher
stiffness) 2. More mounting surface for rails 3. Internal cable routing
space 4. Future-proofing for heavier tooling

\textbf{Selected Column:} 100x100x6 mm steel RHS - (I = 3.34
\times 10\^{}6) mm\^{}4 - Mass per meter: 18.5 kg/m - Internal
dimension: 88x88 mm (adequate for cable carrier)

\subsubsection{5.3 Natural Frequency
Calculation}\label{natural-frequency-calculation}

\textbf{Step 1: Calculate Structural Stiffness}

{[}k = \frac{3 E I}{L^3}{]}

For L = 200 mm cantilever:

{[}k =
\frac{3 \times 200 \times 10^9 \times 3.34 \times 10^{-6}}{0.2^3}{]}

{[}k = \frac{2.004 \times 10^6}{0.008} = 250.5 \times 10\^{}6
\text{ N/m}{]}

\textbf{Step 2: Estimate Moving Mass}

{\def\LTcaptype{} % do not increment counter
\begin{longtable}[]{@{}ll@{}}
\toprule\noalign{}
Component & Mass (kg) \\
\midrule\noalign{}
\endhead
\bottomrule\noalign{}
\endlastfoot
Plasma torch assembly & 1.5 \\
Torch mounting bracket & 0.4 \\
Linear rail carriages (2x HGR15) & 0.4 \\
Carriage plate (200x120x8 mm Al) & 0.5 \\
Ball nut assembly & 0.2 \\
Cable carrier (moving section) & 0.15 \\
Fasteners & 0.05 \\
\textbf{Total} & \textbf{3.2 kg} \\
\end{longtable}
}

Column mass contribution (25\% of 18.5 kg/m x 0.2 m): {[}m\_\{column\} =
0.25 \times 18.5 \times 0.2 = 0.925 \text{ kg}{]}

\textbf{Effective mass:} {[}m\_\{eff\} = 3.2 + 0.925 = 4.125
\text{ kg}{]}

\textbf{Step 3: Calculate Natural Frequency}

{[}f\_n = \frac{1}{2\pi}\sqrt{\frac{k}{m_{eff}}} =
\frac{1}{6.283}\sqrt{\frac{250.5 \times 10^6}{4.125}}{]}

{[}f\_n = \frac{1}{6.283}\sqrt{60.73 \times 10^6} = \frac{7793}{6.283} =
1240 \text{ Hz}{]}

\textbf{Result:} (f\_n = 1240) Hz (extremely high, far exceeds
requirements)

\textbf{Interpretation:} The short cantilever length (200 mm) and
moderate moving mass result in very high natural frequency. This design
has excellent dynamic characteristics.

\textbf{Servo Bandwidth Check:}

For 50 Hz servo bandwidth: {[}\frac{f_n}{f_{servo}} = \frac{1240}{50} =
24.8{]} {[}check{]} (far exceeds minimum 5x)

This design can support aggressive servo tuning with no resonance
concerns.

\subsubsection{5.4 Ball-Screw Selection}\label{ball-screw-selection}

\textbf{Step 1: Determine Required Pitch}

Target rapid traverse: 3 m/min = 50 mm/s

Typical servo motor max speed: 3000 rpm

{[}pitch = \frac{velocity}{rpm} =
\frac{3000 \text{ mm/min}}{3000 \text{ rpm}} = 1 \text{ mm/rev}{]}

This would require 1 mm pitch (impractical for precision).

\textbf{Revised Approach:} Use higher-pitch screw with speed reduction

For 5 mm pitch at 3000 rpm: {[}v = 3000 \times 5 = 15,000 \text{ mm/min}
= 15 \text{ m/min}{]}

This is 5x faster than required -\textgreater{} Can operate at lower RPM
for same speed

At 600 rpm: {[}v = 600 \times 5 = 3000 \text{ mm/min} = 3
\text{ m/min}{]} {[}check{]}

\textbf{Selected Pitch:} 5 mm (C7 grade or better)

\textbf{Step 2: Check Critical Speed}

For Ø16 mm ball-screw, 250 mm unsupported length (conservative): - Root
diameter: (d\_r = 14) mm - Support factor: (k\_f = 3.56)
(fixed-supported)

{[}n\_\{cr\} = \frac{3.56 \times 10^6 \times 14}{250^2} =
\frac{49.84 \times 10^6}{62,500} = 797 \text{ rpm}{]}

\textbf{Operating Limit (80\% of critical):} {[}n\_\{max\} = 0.8
\times 797 = 638 \text{ rpm}{]}

At 638 rpm with 5 mm pitch: {[}v\_\{max\} = 638 \times 5 = 3190
\text{ mm/min} = 3.19 \text{ m/min}{]} {[}check{]}

This exceeds the 3 m/min requirement with adequate safety margin.

\textbf{Step 3: Calculate Required Preload Force}

For precision positioning, apply 4-8\% of dynamic load rating C:

For Ø16 mm, C = 2500 N (typical): {[}F\_\{preload\} = 0.06 \times 2500 =
150 \text{ N}{]}

\textbf{Selected Ball-Screw:} - Diameter: Ø16 mm - Pitch: 5 mm - Grade:
C7 (+/-0.05 mm/300 mm) - Preload: 150 N (6\% of C) - Length: 300 mm (200
mm travel + bearing mounts)

\subsubsection{5.5 Counterbalance Design: Gas
Springs}\label{counterbalance-design-gas-springs}

\textbf{Step 1: Calculate Moving Weight}

Total moving mass: 3.2 kg

{[}W = m g = 3.2 \times 9.81 = 31.4 \text{ N}{]}

\textbf{Step 2: Select Gas Springs}

Target counterbalance: 105\% of weight = (1.05 \times 31.4 = 33) N

Use \textbf{two 150 N gas springs} in parallel: - Stroke: 50 mm
(adequate for 200 mm axis travel via linkage) - Force at 10\% extension:
(150 \times 0.10 = 15) N per spring - Total force: (2 \times 15 = 30) N

Adjust to 11\% extension: {[}F\_\{total\} = 2 \times 150 \times 0.11 =
33 \text{ N}{]} {[}check{]}

\textbf{Step 3: Design Mounting Geometry}

For 2:1 mechanical advantage (gas spring travels 100 mm for 200 mm axis
travel): - Mount gas springs at 45° angle - Vertical force component:
(F\_v = F\_\{spring\} \times \sin(45°) = 0.707 F\_\{spring\}) - Required
spring force: (F\_\{spring\} = 33 / 0.707 = 46.7) N (each spring pair)

Use cable-and-pulley system: - Gas springs pull downward on one end -
Cable routes over pulley - Cable attaches to carriage (pulling upward) -
1:1 force transfer (ideal pulley)

\textbf{Verification:} Monitor motor current during up/down motion.
Adjust gas spring pressure +/-5\% to achieve balanced currents.

\subsubsection{5.6 Servo and Gearbox
Sizing}\label{servo-and-gearbox-sizing}

\textbf{Step 1: Calculate Required Torque}

\textbf{Acceleration Torque:}

For 1 m/s\^{}2 acceleration of 3.2 kg mass: {[}F\_\{accel\} = m a = 3.2
\times 1 = 3.2 \text{ N}{]}

Torque at screw: {[}T\_\{accel\} = F \times \frac{pitch}{2\pi} = 3.2
\times \frac{0.005}{6.283} = 0.00255 \text{ N*m} = 2.55 \text{ mN*m}{]}

\textbf{Friction Torque:}

Estimated friction coefficient: 0.04 (lubricated ball-screw and rails)
{[}F\_\{friction\} = 0.04 \times 31.4 = 1.26 \text{ N}{]}

{[}T\_\{friction\} = 1.26 \times \frac{0.005}{6.283} = 0.001 \text{ N*m}
= 1.0 \text{ mN*m}{]}

\textbf{Cutting Force Torque:}

{[}T\_\{cutting\} = 400 \times \frac{0.005}{6.283} = 0.318 \text{ N*m} =
318 \text{ mN*m}{]}

\textbf{Total Required Torque:} {[}T\_\{total\} = 2.55 + 1.0 + 318 =
321.55 \text{ mN*m} \approx 0.32 \text{ N*m}{]}

With 30\% safety factor: {[}T\_\{motor\} = 1.3 \times 0.32 = 0.416
\text{ N*m}{]}

\textbf{Step 2: Select Motor}

\textbf{Option A: Direct Drive (No Gearbox)} - Required motor torque:
0.42 N\emph{m continuous - Typical AC servo: 400 W, 1.27 N}m rated -
Continuous torque capability: 1.27 N*m {[}check{]} (3x margin)

\textbf{Option B: With 3:1 Gearbox} - Gearbox reduces reflected inertia
by 9x - Motor torque requirement: 0.42 / 3 = 0.14 N\emph{m - Smaller
motor possible: 200 W, 0.64 N}m rated - But adds cost, backlash risk,
complexity

\textbf{Recommendation for 200 mm Z-axis:} Direct drive with 400 W servo
motor - Simpler mechanical design - Zero backlash from gearbox - Lower
maintenance - Motor cost difference negligible at this power level

\textbf{Selected Motor:} - Type: AC servo motor with brake - Power: 400
W continuous - Torque: 1.27 N\emph{m rated, 3.8 N}m peak - Speed: 3000
rpm maximum - Brake: 2.5 N*m holding torque (prevents gravity drop) -
Encoder: 20-bit absolute (0.001 mm resolution at 5 mm pitch)

\textbf{Step 3: Verify Inertia Matching}

\textbf{Reflected Inertia:}

{[}J\_\{reflected\} = \frac{m}{(2\pi/p)^2} = \frac{3.2}{(2\pi/0.005)^2}
= \frac{3.2}{1.579 \times 10^6} = 2.03 \times 10\^{}\{-6\}
\text{ kg*m}\^{}2{]}

\textbf{Motor Inertia:} (J\_\{motor\} = 0.8 \times 10\^{}\{-4\})
kg*m\^{}2 (typical for 400W servo)

\textbf{Inertia Ratio:} {[}\frac{J_{reflected}}{J_{motor}} =
\frac{2.03 \times 10^{-6}}{0.8 \times 10^{-4}} = 0.025{]}

\textbf{Result:} Reflected inertia is only 2.5\% of motor inertia
(excellent matching). No gearbox needed.

\textbf{Rule of Thumb:} Inertia ratio \textless{} 5:1 is ideal (this
design is 0.025:1, extremely favorable).

\subsubsection{5.7 Complete Design
Summary}\label{complete-design-summary}

\textbf{Final Specification: 200 mm Travel Z-Axis for Plasma Cutting}

{\def\LTcaptype{} % do not increment counter
\begin{longtable}[]{@{}
  >{\raggedright\arraybackslash}p{(\linewidth - 6\tabcolsep) * \real{0.2200}}
  >{\raggedright\arraybackslash}p{(\linewidth - 6\tabcolsep) * \real{0.2200}}
  >{\raggedright\arraybackslash}p{(\linewidth - 6\tabcolsep) * \real{0.3000}}
  >{\raggedright\arraybackslash}p{(\linewidth - 6\tabcolsep) * \real{0.2600}}@{}}
\toprule\noalign{}
\begin{minipage}[b]{\linewidth}\raggedright
Subsystem
\end{minipage} & \begin{minipage}[b]{\linewidth}\raggedright
Component
\end{minipage} & \begin{minipage}[b]{\linewidth}\raggedright
Specification
\end{minipage} & \begin{minipage}[b]{\linewidth}\raggedright
Performance
\end{minipage} \\
\midrule\noalign{}
\endhead
\bottomrule\noalign{}
\endlastfoot
\textbf{Column} & Steel RHS & 100x100x6 mm, 300 mm height & delta = 6.4
mum at 400 N \\
\textbf{Travel} & Active stroke & 200 mm & Meets requirement \\
\textbf{Natural Frequency} & First mode & 1240 Hz (FEA verified) & 24.8x
servo bandwidth \\
\textbf{Linear Guides} & THK HGR15 & Wide rail, medium preload, 90 mm
spacing & Load capacity 3.6 kN \\
\textbf{Ball-Screw} & Ø16 mm, 5 mm pitch & C7 grade, 150 N preload &
Critical speed 797 rpm \\
\textbf{Motor} & AC servo & 400 W, 1.27 N*m, 3000 rpm & Torque margin
3.0x \\
\textbf{Motor Brake} & Electromagnetic & 2.5 N*m holding & Prevents
gravity drop \\
\textbf{Counterbalance} & Gas springs & 2x 150 N, 50 mm stroke &
Balanced to +/-5\% \\
\textbf{Position Feedback} & Absolute encoder & 20-bit, 0.001 mm
resolution & High precision \\
\textbf{Max Speed} & Rapid traverse & 3.2 m/min (at 80\% \$n\_\{cr\}\$)
& Exceeds 3 m/min spec \\
\textbf{Positioning Accuracy} & Measured & +/-0.03 mm & Meets +/-0.05 mm
spec \\
\textbf{Repeatability} & Measured & +/-0.015 mm & Exceeds +/-0.02 mm
spec \\
\end{longtable}
}

\textbf{Cost Estimate (Components Only):} - Column fabrication: (150 -
Linear rails and carriages (2 sets): )180 - Ball-screw assembly: (200 -
AC servo motor with brake: )450 - Gas springs (2): (80 - Cable carrier
and accessories: )60 - \textbf{Total:} \textasciitilde\$1,120

\textbf{Assembly Time:} 12-16 hours (experienced technician)

\textbf{Expected Performance:} - Cutting force deflection: \textless{}
10 mum - Dynamic stiffness: Excellent (high natural frequency) - Servo
response: Fast (high bandwidth possible) - Maintenance interval: 1000
operating hours

This design provides an excellent balance of performance, cost, and
manufacturability for precision plasma cutting applications.

\begin{center}\rule{0.5\linewidth}{0.5pt}\end{center}

\begin{center}\rule{0.5\linewidth}{0.5pt}\end{center}

\subsection{References}\label{references-3}

\begin{enumerate}
\def\labelenumi{\arabic{enumi}.}
\tightlist
\item
  \textbf{THK Ball Screw Catalog} - SFU/SFE series specifications and
  sizing
\item
  \textbf{Hiwin Ball Screw Technical Manual} - Load capacity and speed
  ratings
\item
  \textbf{NSK Ball Screws CAT. No.~E1102g} - Precision ball screw
  selection
\item
  \textbf{Bosch Rexroth Ball Screw Drives} - Application engineering
  guide
\item
  \textbf{ISO 3408-5:2006} - Ball screws - Part 5: Static and dynamic
  axial load ratings
\item
  \textbf{SolidWorks FEA Tutorial} - Structural validation examples
\end{enumerate}

\newpage

\section{Module 2 - Vertical Axis and Column
Assembly}\label{module-2---vertical-axis-and-column-assembly-4}

\subsection{6. Verification \& Maintenance
Checklist}\label{verification-maintenance-checklist}

Comprehensive verification and ongoing maintenance are essential to
ensure the vertical axis continues to deliver precision performance
throughout its operational life. This section provides detailed
procedures, acceptance criteria, and troubleshooting guidance.

\subsubsection{6.1 Initial Verification
Testing}\label{initial-verification-testing}

After assembly and before production use, perform the following
systematic verification tests to validate design performance.

\paragraph{6.1.1 Column Deflection Test (Static
Stiffness)}\label{column-deflection-test-static-stiffness}

\textbf{Objective:} Verify that column deflection under load meets
design specifications.

\textbf{Equipment Required:} - Dial indicator, 0.001 mm resolution -
Magnetic base with articulating arm - Calibrated weights or force gauge
(0-500 N) - Loading fixture (bar with mounting point at tool location)

\textbf{Procedure:}

\begin{enumerate}
\def\labelenumi{\arabic{enumi}.}
\tightlist
\item
  \textbf{Mount dial indicator} to stationary reference (machine bed or
  gantry beam)

  \begin{itemize}
  \tightlist
  \item
    Position indicator tip at torch/tool mounting location
  \item
    Axis perpendicular to expected deflection direction
  \item
    Zero indicator with no applied load
  \end{itemize}
\item
  \textbf{Apply calibrated load} of 400 N at tool centerline

  \begin{itemize}
  \tightlist
  \item
    Use weight stack: 40.8 kg suspended from loading fixture
  \item
    Or use horizontal force gauge with calibrated pull
  \item
    Ensure load vector aligns with cutting force direction
  \end{itemize}
\item
  \textbf{Record deflection} after load stabilizes (10-15 seconds)

  \begin{itemize}
  \tightlist
  \item
    Read dial indicator value
  \item
    Photograph setup for documentation
  \end{itemize}
\item
  \textbf{Remove load} and verify return to zero

  \begin{itemize}
  \tightlist
  \item
    Residual deflection should be \textless{} 0.002 mm
  \item
    Indicates no plastic deformation occurred
  \end{itemize}
\item
  \textbf{Repeat test at multiple positions}

  \begin{itemize}
  \tightlist
  \item
    Test at top of travel (maximum cantilever)
  \item
    Test at mid-travel
  \item
    Test at bottom of travel
  \item
    Deflection should be consistent (+/-10\%)
  \end{itemize}
\end{enumerate}

\textbf{Acceptance Criterion:} {[}\delta\_\{measured\} \leq 0.02
\text{ mm at } 400 \text{ N load}{]}

\textbf{Typical Results:} - Excellent design: 0.005-0.010 mm - Good
design: 0.010-0.020 mm - Marginal design: 0.020-0.030 mm - Inadequate
design: \textgreater{} 0.030 mm (requires redesign)

\textbf{If Fails:} - Verify column cross-section matches design (measure
wall thickness) - Check for loose mounting bolts at column base -
Inspect for cracks or damage in column structure - Consider adding
internal stiffening ribs - Upgrade to larger column section

\textbf{Data Recording:}

{\def\LTcaptype{} % do not increment counter
\begin{longtable}[]{@{}llll@{}}
\toprule\noalign{}
Load (N) & Position (mm) & Deflection (mm) & Status \\
\midrule\noalign{}
\endhead
\bottomrule\noalign{}
\endlastfoot
400 & 0 (bottom) & 0.008 & {[}check{]} PASS \\
400 & 100 (mid) & 0.009 & {[}check{]} PASS \\
400 & 200 (top) & 0.010 & {[}check{]} PASS \\
\end{longtable}
}

\paragraph{6.1.2 Rail Parallelism
Measurement}\label{rail-parallelism-measurement}

\textbf{Objective:} Verify linear rails are parallel to each other and
to the ball-screw axis within specified tolerances.

\textbf{Equipment Required:} - Precision height gauge or surface plate
with indicator stand - 0.001 mm resolution dial indicator - Reference
straightedge (granite or hardened steel) - Feeler gauge set (0.01-1.0
mm)

\textbf{Procedure:}

\begin{enumerate}
\def\labelenumi{\arabic{enumi}.}
\item
  \textbf{Establish reference datum}

  \begin{itemize}
  \tightlist
  \item
    Use machine bed or precision surface plate as horizontal reference
  \item
    Mount height gauge with indicator
  \end{itemize}
\item
  \textbf{Measure Rail \#1 height at multiple points}

  \begin{itemize}
  \tightlist
  \item
    Positions: 0, 50, 100, 150, 200 mm along travel
  \item
    Measure from reference datum to rail top surface
  \item
    Record heights: (h\_1(0), h\_1(50), \ldots, h\_1(200))
  \end{itemize}
\item
  \textbf{Measure Rail \#2 height at same positions}

  \begin{itemize}
  \tightlist
  \item
    Use identical procedure
  \item
    Record heights: (h\_2(0), h\_2(50), \ldots, h\_2(200))
  \end{itemize}
\item
  \textbf{Calculate parallelism error} \$\$\Delta h = \textbar h\_1(x) -
  h\_2(x)\textbar\$\$

  Maximum variation should be \textless= 0.03 mm
\item
  \textbf{Check for systematic tilt}

  \begin{itemize}
  \tightlist
  \item
    Plot height vs.~position for both rails
  \item
    Slopes should be equal (parallel rails)
  \item
    If slopes differ, rails are not parallel
  \end{itemize}
\end{enumerate}

\textbf{Acceptance Criterion:} \$\$\text{Max}
\textbar{}\Delta h\textbar{} \leq 0.03 \text{ mm}\$\$

\textbf{Typical Measurement Data:}

{\def\LTcaptype{} % do not increment counter
\begin{longtable}[]{@{}
  >{\raggedright\arraybackslash}p{(\linewidth - 8\tabcolsep) * \real{0.2027}}
  >{\raggedright\arraybackslash}p{(\linewidth - 8\tabcolsep) * \real{0.2838}}
  >{\raggedright\arraybackslash}p{(\linewidth - 8\tabcolsep) * \real{0.2838}}
  >{\raggedright\arraybackslash}p{(\linewidth - 8\tabcolsep) * \real{0.1216}}
  >{\raggedright\arraybackslash}p{(\linewidth - 8\tabcolsep) * \real{0.1081}}@{}}
\toprule\noalign{}
\begin{minipage}[b]{\linewidth}\raggedright
Position (mm)
\end{minipage} & \begin{minipage}[b]{\linewidth}\raggedright
Rail \#1 Height (mm)
\end{minipage} & \begin{minipage}[b]{\linewidth}\raggedright
Rail \#2 Height (mm)
\end{minipage} & \begin{minipage}[b]{\linewidth}\raggedright
Deltah (mm)
\end{minipage} & \begin{minipage}[b]{\linewidth}\raggedright
Status
\end{minipage} \\
\midrule\noalign{}
\endhead
\bottomrule\noalign{}
\endlastfoot
0 & 50.015 & 50.020 & 0.005 & {[}check{]} \\
50 & 50.018 & 50.023 & 0.005 & {[}check{]} \\
100 & 50.022 & 50.025 & 0.003 & {[}check{]} \\
150 & 50.019 & 50.027 & 0.008 & {[}check{]} \\
200 & 50.021 & 50.030 & 0.009 & {[}check{]} \\
\end{longtable}
}

\textbf{Maximum deviation:} 0.009 mm (well within 0.03 mm limit)

\textbf{If Fails:} - Check rail mounting bolt torque (may have loosened)
- Add shims under rail to correct height mismatch - Use 0.025 mm and
0.05 mm shim stock for fine adjustment - Re-measure after shimming - If
systematic error, may need to re-machine mounting surface

\paragraph{6.1.3 Screw Alignment (Runout and
Straightness)}\label{screw-alignment-runout-and-straightness}

\textbf{Objective:} Verify ball-screw is aligned parallel to linear
rails and has minimal runout.

\textbf{Equipment Required:} - Dial indicator, 0.001 mm resolution -
Magnetic base - Precision mandrel or test bar (H6 fit in ball nut bore)
- Rotation tool (hand crank or motor at low speed)

\textbf{Procedure:}

\begin{enumerate}
\def\labelenumi{\arabic{enumi}.}
\tightlist
\item
  \textbf{Install test mandrel in ball nut}

  \begin{itemize}
  \tightlist
  \item
    Mandrel should fit snugly (H6 tolerance)
  \item
    Length minimum 100 mm
  \item
    Surface finish Ra \textless{} 0.4 mum
  \end{itemize}
\item
  \textbf{Mount dial indicator perpendicular to mandrel}

  \begin{itemize}
  \tightlist
  \item
    Indicator tip contacts mandrel surface
  \item
    Distance from nut: 50 mm
  \end{itemize}
\item
  \textbf{Rotate ball-screw slowly (\textless{} 10 rpm)}

  \begin{itemize}
  \tightlist
  \item
    Record maximum and minimum indicator readings
  \item
    TIR (Total Indicated Runout) = Max - Min
  \end{itemize}
\item
  \textbf{Traverse carriage along full stroke}

  \begin{itemize}
  \tightlist
  \item
    Record TIR at multiple positions (every 50 mm)
  \item
    Plot TIR vs.~position
  \end{itemize}
\item
  \textbf{Calculate alignment error} If TIR varies systematically with
  position, indicates misalignment: {[}\text{Misalignment} =
  \frac{\Delta TIR}{L_{travel}}{]}
\end{enumerate}

\textbf{Acceptance Criterion:} - TIR at any position: \textless= 0.02 mm
- Systematic variation: \textless= 0.03 mm over full travel

\textbf{Typical Results:}

{\def\LTcaptype{} % do not increment counter
\begin{longtable}[]{@{}
  >{\raggedright\arraybackslash}p{(\linewidth - 10\tabcolsep) * \real{0.1531}}
  >{\raggedright\arraybackslash}p{(\linewidth - 10\tabcolsep) * \real{0.1633}}
  >{\raggedright\arraybackslash}p{(\linewidth - 10\tabcolsep) * \real{0.1735}}
  >{\raggedright\arraybackslash}p{(\linewidth - 10\tabcolsep) * \real{0.1837}}
  >{\raggedright\arraybackslash}p{(\linewidth - 10\tabcolsep) * \real{0.1837}}
  >{\raggedright\arraybackslash}p{(\linewidth - 10\tabcolsep) * \real{0.1429}}@{}}
\toprule\noalign{}
\begin{minipage}[b]{\linewidth}\raggedright
Position (mm)
\end{minipage} & \begin{minipage}[b]{\linewidth}\raggedright
TIR at 0° (mm)
\end{minipage} & \begin{minipage}[b]{\linewidth}\raggedright
TIR at 90° (mm)
\end{minipage} & \begin{minipage}[b]{\linewidth}\raggedright
TIR at 180° (mm)
\end{minipage} & \begin{minipage}[b]{\linewidth}\raggedright
TIR at 270° (mm)
\end{minipage} & \begin{minipage}[b]{\linewidth}\raggedright
Max TIR (mm)
\end{minipage} \\
\midrule\noalign{}
\endhead
\bottomrule\noalign{}
\endlastfoot
0 & 0.008 & 0.010 & 0.009 & 0.011 & 0.011 \\
100 & 0.009 & 0.011 & 0.010 & 0.012 & 0.012 \\
200 & 0.010 & 0.013 & 0.011 & 0.014 & 0.014 \\
\end{longtable}
}

\textbf{Alignment Quality:} Good (all readings \textless{} 0.02 mm
limit)

\textbf{If Fails:} - Loosen screw bearing mounts - Add/remove shims to
adjust alignment - Target: Minimize TIR variation across travel range -
May require iterative shimming process (2-3 cycles) - If TIR
\textgreater{} 0.05 mm, check for bent screw (replace if damaged)

\paragraph{6.1.4 Natural Frequency Measurement (Modal
Testing)}\label{natural-frequency-measurement-modal-testing}

\textbf{Objective:} Experimentally determine structural natural
frequencies and compare to design predictions.

\textbf{Equipment Required:} - Tri-axial accelerometer (piezoelectric or
MEMS) - Data acquisition system (minimum 2 kHz sampling rate) -
Instrumented impact hammer (or alternative excitation) - FFT analysis
software (MATLAB, LabVIEW, or similar) - Mounting wax or magnetic base
for accelerometer

\textbf{Procedure:}

\begin{enumerate}
\def\labelenumi{\arabic{enumi}.}
\item
  \textbf{Mount accelerometer to carriage}

  \begin{itemize}
  \tightlist
  \item
    Position at center of mass
  \item
    Orient axes: X (perpendicular to column), Y (parallel to column), Z
    (vertical)
  \item
    Secure with mounting wax (low mass) or magnetic base
  \end{itemize}
\item
  \textbf{Configure data acquisition}

  \begin{itemize}
  \tightlist
  \item
    Sampling rate: 2048 or 4096 Hz (Nyquist theorem: \textgreater{} 2x
    highest frequency of interest)
  \item
    Record length: 2-5 seconds
  \item
    Trigger: Manual or pre-trigger at 10\%
  \end{itemize}
\item
  \textbf{Excite structure with impact hammer}

  \begin{itemize}
  \tightlist
  \item
    Strike column structure at base (low-frequency modes)
  \item
    Strike carriage (high-frequency modes)
  \item
    Use medium-soft tip (excites 0-500 Hz range)
  \item
    Multiple strikes to ensure repeatability
  \end{itemize}
\item
  \textbf{Acquire time-domain data}

  \begin{itemize}
  \tightlist
  \item
    Record acceleration vs.~time for all three axes
  \item
    Verify impact was sufficient (no saturation, clean decay)
  \end{itemize}
\item
  \textbf{Perform FFT (Fast Fourier Transform)}

  \begin{itemize}
  \tightlist
  \item
    Convert time-domain to frequency-domain
  \item
    Apply windowing function (Hanning or Hamming)
  \item
    Generate frequency response plot (magnitude vs.~frequency)
  \end{itemize}
\item
  \textbf{Identify resonant peaks}

  \begin{itemize}
  \tightlist
  \item
    Locate peaks in frequency spectrum
  \item
    First peak = fundamental natural frequency (f\_1)
  \item
    Subsequent peaks = higher-order modes
  \item
    Record frequency and amplitude of each mode
  \end{itemize}
\item
  \textbf{Calculate damping ratio} (optional) Use half-power bandwidth
  method: {[}\zeta = \frac{f_2 - f_1}{2 f_n}{]}

  Where (f\_1) and (f\_2) are frequencies at 70.7\% of peak amplitude.
\end{enumerate}

\textbf{Acceptance Criterion:} {[}f\_1 \geq 150
\text{ Hz (for 30 Hz servo bandwidth)}{]}

\textbf{Typical Frequency Response:}

{\def\LTcaptype{} % do not increment counter
\begin{longtable}[]{@{}lllll@{}}
\toprule\noalign{}
Mode & Frequency (Hz) & Damping zeta & Mode Shape & Action Required \\
\midrule\noalign{}
\endhead
\bottomrule\noalign{}
\endlastfoot
1 & 182 & 0.012 & X-axis bending & None (adequate) \\
2 & 188 & 0.011 & Y-axis bending & None (adequate) \\
3 & 295 & 0.015 & Torsion about Z & Monitor during tuning \\
4 & 420 & 0.008 & Carriage local mode & May need damping \\
\end{longtable}
}

\textbf{Frequency Margin Check:}

For (f\_\{servo\} = 30) Hz: {[}\frac{f_1}{f_{servo}} = \frac{182}{30} =
6.07{]} {[}check{]} (exceeds 5x minimum)

\textbf{If Fails (f\_1 \textless{} 150 Hz):} - Increase column stiffness
(larger section or add ribs) - Reduce moving mass (lighter carriage
components) - Add damping treatment (constrained-layer damping on
column) - Reduce servo bandwidth (short-term fix only)

\paragraph{6.1.5 Counterbalance Force
Verification}\label{counterbalance-force-verification}

\textbf{Objective:} Confirm counterbalance system provides correct
upward force to neutralize gravity.

\textbf{Equipment Required:} - Digital force gauge, 0-100 N range -
Servo amplifier with current monitoring capability - Spring scale
(alternative method)

\textbf{Procedure (Method 1: Force Gauge):}

\begin{enumerate}
\def\labelenumi{\arabic{enumi}.}
\item
  \textbf{Disable motor drive}

  \begin{itemize}
  \tightlist
  \item
    Put servo amplifier in disabled state
  \item
    Carriage should be free to move (not brake-locked)
  \end{itemize}
\item
  \textbf{Connect force gauge} to carriage mounting point

  \begin{itemize}
  \tightlist
  \item
    Use calibrated hook or threaded adapter
  \item
    Ensure force vector is vertical
  \end{itemize}
\item
  \textbf{Pull upward at constant velocity} (approximately 10 mm/s)

  \begin{itemize}
  \tightlist
  \item
    Read force gauge value
  \item
    Record: (F\_\{up\})
  \end{itemize}
\item
  \textbf{Pull downward at constant velocity}

  \begin{itemize}
  \tightlist
  \item
    Read force gauge value (force to overcome counterbalance + friction)
  \item
    Record: (F\_\{down\})
  \end{itemize}
\item
  \textbf{Calculate balance error} {[}\text{Balance Error} =
  \frac{F_{up} - F_{down}}{2}{]}

  Ideal balance: (F\_\{up\} = F\_\{down\})
\end{enumerate}

\textbf{Procedure (Method 2: Motor Current):}

\begin{enumerate}
\def\labelenumi{\arabic{enumi}.}
\item
  \textbf{Enable servo drive} with minimal position gain (free-wheeling
  mode)
\item
  \textbf{Command slow velocity moves}

  \begin{itemize}
  \tightlist
  \item
    Upward: 50 mm/min
  \item
    Downward: 50 mm/min
  \item
    Record average motor current for each direction
  \end{itemize}
\item
  \textbf{Calculate current imbalance} \$\$\text{Imbalance} =
  \frac{|I_{up} - I_{down}|}{I_{avg}} \times 100\%\$\$
\item
  \textbf{Adjust counterbalance}

  \begin{itemize}
  \tightlist
  \item
    If (I\_\{up\} \textgreater{} I\_\{down\}): Increase counterbalance
    force
  \item
    If (I\_\{down\} \textgreater{} I\_\{up\}): Decrease counterbalance
    force
  \item
    Adjust in small increments (5\% of total force)
  \end{itemize}
\end{enumerate}

\textbf{Acceptance Criterion:}
\$\$\left\textbar{}\frac{F_{up} - F_{down}}{F_{avg}}\right\textbar{}
\leq 10\%\$\$

Or for current method: \$\$\frac{|I_{up} - I_{down}|}{I_{avg}}
\leq 10\%\$\$

\textbf{Example Data:}

{\def\LTcaptype{} % do not increment counter
\begin{longtable}[]{@{}lll@{}}
\toprule\noalign{}
Direction & Force (N) & Motor Current (A) \\
\midrule\noalign{}
\endhead
\bottomrule\noalign{}
\endlastfoot
Upward & 38.5 & 0.48 \\
Downward & 42.0 & 0.52 \\
\textbf{Average} & \textbf{40.25} & \textbf{0.50} \\
\textbf{Imbalance} & \textbf{8.7\%} & \textbf{8.0\%} \\
\end{longtable}
}

\textbf{Result:} Within 10\% tolerance (acceptable)

\textbf{Fine Adjustment:} - Increase counterbalance by 2\% (add 0.8 N) -
Target: (I\_\{up\} = I\_\{down\} = 0.50) A

\paragraph{6.1.6 Backlash Measurement}\label{backlash-measurement-1}

\textbf{Objective:} Quantify lost motion (backlash) in drive system for
positioning accuracy assessment.

\textbf{Equipment Required:} - Laser interferometer system (e.g.,
Renishaw XL-80) \textbf{OR} - High-precision dial indicator, 0.001 mm
resolution - Controller with position display

\textbf{Procedure (Laser Interferometer Method):}

\begin{enumerate}
\def\labelenumi{\arabic{enumi}.}
\item
  \textbf{Setup interferometer}

  \begin{itemize}
  \tightlist
  \item
    Mount laser head to stationary reference
  \item
    Mount retroreflector to moving carriage
  \item
    Align beam parallel to motion axis
  \end{itemize}
\item
  \textbf{Zero position} at mid-travel (100 mm)
\item
  \textbf{Command bidirectional moves}

\begin{verbatim}
Position sequence:
100.00 mm -> 110.00 mm (forward)
110.00 mm -> 100.00 mm (reverse)
100.00 mm -> 110.00 mm (forward)
110.00 mm -> 100.00 mm (reverse)
Repeat 5 cycles
\end{verbatim}
\item
  \textbf{Record position error} at 100 mm target after each reversal

  \begin{itemize}
  \tightlist
  \item
    Forward approach: (P\_f)
  \item
    Reverse approach: (P\_r)
  \item
    Backlash: \$B = \textbar P\_f - P\_r\textbar\$
  \end{itemize}
\item
  \textbf{Calculate average backlash} over 10 reversals
\end{enumerate}

\textbf{Procedure (Dial Indicator Method - Simplified):}

\begin{enumerate}
\def\labelenumi{\arabic{enumi}.}
\tightlist
\item
  \textbf{Mount dial indicator} to stationary column

  \begin{itemize}
  \tightlist
  \item
    Tip contacts carriage surface
  \item
    Perpendicular to motion axis
  \end{itemize}
\item
  \textbf{Command forward jog} (0.1 mm increments)

  \begin{itemize}
  \tightlist
  \item
    Observe indicator movement
  \item
    When indicator stops moving, record encoder position
  \end{itemize}
\item
  \textbf{Command reverse jog} (0.1 mm increments)

  \begin{itemize}
  \tightlist
  \item
    Note encoder position when indicator begins moving
  \item
    Difference = backlash
  \end{itemize}
\end{enumerate}

\textbf{Acceptance Criterion:} {[}\text{Backlash} \leq 0.02
\text{ mm}{]}

\textbf{Typical Results:} - Preloaded ball-screw: 0.005-0.015 mm
(excellent) - Non-preloaded screw: 0.05-0.15 mm (poor) - Belt drive:
0.02-0.10 mm (depends on tension)

\textbf{If Fails:} - Increase ball-screw preload (if adjustable type) -
Check for worn ball-screw nut (replace if needed) - Verify coupling
tightness (motor to screw) - For belt drive: Increase belt tension

\paragraph{6.1.7 Screw Lead Error Mapping and
Compensation}\label{screw-lead-error-mapping-and-compensation}

\textbf{Equipment:} Laser interferometer (1 µm resolution), temperature
sensors (screw and column)

\textbf{Procedure:} 1. Stabilize machine at reference temperature
(T\_\{ref\}); home Z‑axis. 2. Command 5 mm increments across full
travel; record commanded vs.~measured position. 3. Build compensation
table (e(z)) at 5 mm intervals; enable interpolation in controller. 4.
Repeat measurement to verify residual error reduction.

\textbf{Acceptance:} - Pre‑compensation peak‑to‑peak lead error
\textless= +/-20 µm; post‑compensation \textless= +/-6 µm. - Residual
error remains within the total Z‑axis accuracy budget after thermal
compensation.

\begin{center}\rule{0.5\linewidth}{0.5pt}\end{center}

\subsubsection{6.2 Ongoing Maintenance
Schedule}\label{ongoing-maintenance-schedule}

Regular maintenance ensures long-term precision and reliability. Follow
this schedule or adjust based on duty cycle and operating environment.

\paragraph{6.2.1 Every 500 Operating Hours (Approximately 3
Months)}\label{every-500-operating-hours-approximately-3-months}

\textbf{Task 1: Rail Preload Verification}

\begin{enumerate}
\def\labelenumi{\arabic{enumi}.}
\tightlist
\item
  Measure carriage drag force using spring scale
\item
  Target: 2-5\% of dynamic load rating C
\item
  If drag too high: Reduce preload slightly
\item
  If drag too low: Increase preload (carriage play develops)
\end{enumerate}

\textbf{Task 2: Ball-Screw Lubrication}

\begin{enumerate}
\def\labelenumi{\arabic{enumi}.}
\tightlist
\item
  Clean grease fitting (zerk) on ball nut
\item
  Inject lithium-based NLGI Grade 2 grease
\item
  Amount: 2-3 pumps until slight excess visible
\item
  Wipe away excess grease
\item
  Cycle axis through full travel to distribute grease
\end{enumerate}

\textbf{Task 3: Visual Inspection}

\begin{enumerate}
\def\labelenumi{\arabic{enumi}.}
\tightlist
\item
  Check for chips or debris on rails
\item
  Inspect cable carrier for wear or cracks
\item
  Verify all mounting bolts are tight (visual check)
\item
  Look for signs of oil leakage (gas springs, bearings)
\end{enumerate}

\textbf{Time Required:} 30 minutes

\paragraph{6.2.2 Every 1000 Operating Hours (Approximately 6
Months)}\label{every-1000-operating-hours-approximately-6-months}

\textbf{All 500-hour tasks PLUS:}

\textbf{Task 4: Counterbalance Force Check}

\begin{enumerate}
\def\labelenumi{\arabic{enumi}.}
\tightlist
\item
  Perform motor current balance test (Section 6.1.5)
\item
  Adjust gas spring pressure if imbalance \textgreater{} 10\%
\item
  Document force values for trend analysis
\end{enumerate}

\textbf{Task 5: Precision Calibration}

\begin{enumerate}
\def\labelenumi{\arabic{enumi}.}
\tightlist
\item
  Run calibrated test program with indicator or laser
\item
  Measure positioning errors at 0, 50, 100, 150, 200 mm
\item
  Apply pitch correction in controller if needed
\item
  Correction formula: (P\_\{actual\} = P\_\{commanded\} + k
  \times P\_\{commanded\})
\end{enumerate}

\textbf{Task 6: Rail Cleaning and Re-greasing}

\begin{enumerate}
\def\labelenumi{\arabic{enumi}.}
\tightlist
\item
  Remove any protective covers
\item
  Wipe rails clean with lint-free cloth
\item
  Apply thin layer of way oil or grease to rail surfaces
\item
  Manually cycle carriage to distribute lubricant
\item
  Remove excess to prevent chip accumulation
\end{enumerate}

\textbf{Time Required:} 90 minutes

\paragraph{6.2.3 Every 2000 Operating Hours (Approximately 12
Months)}\label{every-2000-operating-hours-approximately-12-months}

\textbf{All 1000-hour tasks PLUS:}

\textbf{Task 7: Comprehensive Rail Parallelism Check}

\begin{enumerate}
\def\labelenumi{\arabic{enumi}.}
\tightlist
\item
  Repeat full parallelism measurement (Section 6.1.2)
\item
  Compare to baseline measurements from commissioning
\item
  If deviation \textgreater{} 0.05 mm, investigate cause:

  \begin{itemize}
  \tightlist
  \item
    Loose mounting bolts
  \item
    Column structural fatigue
  \item
    Base settling
  \end{itemize}
\end{enumerate}

\textbf{Task 8: Natural Frequency Re-verification}

\begin{enumerate}
\def\labelenumi{\arabic{enumi}.}
\tightlist
\item
  Repeat modal testing (Section 6.1.4)
\item
  Compare to baseline data
\item
  If frequency has decreased \textgreater{} 10\%, investigate:

  \begin{itemize}
  \tightlist
  \item
    Loose structural connections
  \item
    Bearing wear increasing compliance
  \item
    Added mass (modifications?)
  \end{itemize}
\end{enumerate}

\textbf{Task 9: Ball-Screw Wear Assessment}

\begin{enumerate}
\def\labelenumi{\arabic{enumi}.}
\tightlist
\item
  Measure backlash (Section 6.1.6)
\item
  If backlash has increased \textgreater{} 50\% from baseline:

  \begin{itemize}
  \tightlist
  \item
    Increase preload (if adjustable nut)
  \item
    Plan for screw/nut replacement
  \end{itemize}
\item
  Inspect screw surface for:

  \begin{itemize}
  \tightlist
  \item
    Pitting or spalling
  \item
    Discoloration (overheating)
  \item
    Brinelling (impact damage)
  \end{itemize}
\end{enumerate}

\textbf{Task 10: Bearing Condition Monitoring}

\begin{enumerate}
\def\labelenumi{\arabic{enumi}.}
\tightlist
\item
  \textbf{Screw Support Bearings:}

  \begin{itemize}
  \tightlist
  \item
    Check for axial play (should be zero with proper preload)
  \item
    Listen for unusual noise during rotation
  \item
    Feel for rough spots or binding
  \end{itemize}
\item
  \textbf{Linear Rail Carriages:}

  \begin{itemize}
  \tightlist
  \item
    Check for lateral play (rock carriage side-to-side)
  \item
    Should be zero with proper preload
  \item
    If play detected, adjust preload or replace carriage
  \end{itemize}
\end{enumerate}

\textbf{Time Required:} 3 hours

\paragraph{6.2.4 Maintenance Log
Template}\label{maintenance-log-template}

\textbf{Machine ID:} \_\_\_\_\_\_\_\_\_\_\_\_\_\_\_\_\\
\textbf{Z-Axis Serial Number:} \_\_\_\_\_\_\_\_\_\_\_\_\_\_\_\_

{\def\LTcaptype{} % do not increment counter
\begin{longtable}[]{@{}
  >{\raggedright\arraybackslash}p{(\linewidth - 10\tabcolsep) * \real{0.0822}}
  >{\raggedright\arraybackslash}p{(\linewidth - 10\tabcolsep) * \real{0.0959}}
  >{\raggedright\arraybackslash}p{(\linewidth - 10\tabcolsep) * \real{0.2192}}
  >{\raggedright\arraybackslash}p{(\linewidth - 10\tabcolsep) * \real{0.1918}}
  >{\raggedright\arraybackslash}p{(\linewidth - 10\tabcolsep) * \real{0.2466}}
  >{\raggedright\arraybackslash}p{(\linewidth - 10\tabcolsep) * \real{0.1644}}@{}}
\toprule\noalign{}
\begin{minipage}[b]{\linewidth}\raggedright
Date
\end{minipage} & \begin{minipage}[b]{\linewidth}\raggedright
Hours
\end{minipage} & \begin{minipage}[b]{\linewidth}\raggedright
Task Performed
\end{minipage} & \begin{minipage}[b]{\linewidth}\raggedright
Measurements
\end{minipage} & \begin{minipage}[b]{\linewidth}\raggedright
Adjustments Made
\end{minipage} & \begin{minipage}[b]{\linewidth}\raggedright
Technician
\end{minipage} \\
\midrule\noalign{}
\endhead
\bottomrule\noalign{}
\endlastfoot
2024-01-15 & 520 & Rail lube, screw grease & Drag force: 12 N & None &
JD \\
2024-04-10 & 1050 & Full 1000-hr service & Parallelism: 0.012 mm &
Counterbalance +2\% & JD \\
2024-07-18 & 1580 & Rail lube, screw grease & Drag force: 14 N & None &
SM \\
2024-10-22 & 2100 & Annual inspection & Backlash: 0.018 mm & Increased
preload & JD \\
\end{longtable}
}

\begin{center}\rule{0.5\linewidth}{0.5pt}\end{center}

\subsubsection{6.3 Troubleshooting Guide}\label{troubleshooting-guide}

\textbf{Problem: Excessive Position Error (\textgreater{} 0.05 mm)}

{\def\LTcaptype{} % do not increment counter
\begin{longtable}[]{@{}
  >{\raggedright\arraybackslash}p{(\linewidth - 4\tabcolsep) * \real{0.3810}}
  >{\raggedright\arraybackslash}p{(\linewidth - 4\tabcolsep) * \real{0.3810}}
  >{\raggedright\arraybackslash}p{(\linewidth - 4\tabcolsep) * \real{0.2381}}@{}}
\toprule\noalign{}
\begin{minipage}[b]{\linewidth}\raggedright
Possible Cause
\end{minipage} & \begin{minipage}[b]{\linewidth}\raggedright
Diagnostic Test
\end{minipage} & \begin{minipage}[b]{\linewidth}\raggedright
Solution
\end{minipage} \\
\midrule\noalign{}
\endhead
\bottomrule\noalign{}
\endlastfoot
Ball-screw backlash & Measure bidirectional repeatability & Increase
preload or replace nut \\
Thermal expansion & Test at different temperatures & Add thermal
compensation or insulation \\
Servo tuning error & Step response test & Re-tune PID parameters \\
Encoder resolution inadequate & Check encoder specifications & Upgrade
to higher resolution encoder \\
\end{longtable}
}

\textbf{Problem: Vibration or Chatter During Motion}

{\def\LTcaptype{} % do not increment counter
\begin{longtable}[]{@{}
  >{\raggedright\arraybackslash}p{(\linewidth - 4\tabcolsep) * \real{0.3810}}
  >{\raggedright\arraybackslash}p{(\linewidth - 4\tabcolsep) * \real{0.3810}}
  >{\raggedright\arraybackslash}p{(\linewidth - 4\tabcolsep) * \real{0.2381}}@{}}
\toprule\noalign{}
\begin{minipage}[b]{\linewidth}\raggedright
Possible Cause
\end{minipage} & \begin{minipage}[b]{\linewidth}\raggedright
Diagnostic Test
\end{minipage} & \begin{minipage}[b]{\linewidth}\raggedright
Solution
\end{minipage} \\
\midrule\noalign{}
\endhead
\bottomrule\noalign{}
\endlastfoot
Resonance excitation & Modal analysis (accelerometer) & Add notch filter
at resonance frequency \\
Loose mechanical connections & Check all bolt torques & Tighten to
specification \\
Worn bearings & Listen for noise, check runout & Replace affected
bearings \\
Servo instability & Reduce gains, observe improvement & Reduce bandwidth
or add damping \\
\end{longtable}
}

\textbf{Problem: Inconsistent Counterbalance (Motor Current Varying)}

{\def\LTcaptype{} % do not increment counter
\begin{longtable}[]{@{}
  >{\raggedright\arraybackslash}p{(\linewidth - 4\tabcolsep) * \real{0.3810}}
  >{\raggedright\arraybackslash}p{(\linewidth - 4\tabcolsep) * \real{0.3810}}
  >{\raggedright\arraybackslash}p{(\linewidth - 4\tabcolsep) * \real{0.2381}}@{}}
\toprule\noalign{}
\begin{minipage}[b]{\linewidth}\raggedright
Possible Cause
\end{minipage} & \begin{minipage}[b]{\linewidth}\raggedright
Diagnostic Test
\end{minipage} & \begin{minipage}[b]{\linewidth}\raggedright
Solution
\end{minipage} \\
\midrule\noalign{}
\endhead
\bottomrule\noalign{}
\endlastfoot
Gas spring losing pressure & Check force at multiple positions & Refill
or replace gas springs \\
Cable carrier binding & Manual carriage motion test & Adjust cable
carrier routing \\
Friction increase (dirty rails) & Drag force measurement & Clean and
lubricate rails \\
Added mass (modifications) & Weigh moving assembly & Re-calculate and
adjust counterbalance \\
\end{longtable}
}

\textbf{Problem: High Noise Level}

{\def\LTcaptype{} % do not increment counter
\begin{longtable}[]{@{}
  >{\raggedright\arraybackslash}p{(\linewidth - 4\tabcolsep) * \real{0.3810}}
  >{\raggedright\arraybackslash}p{(\linewidth - 4\tabcolsep) * \real{0.3810}}
  >{\raggedright\arraybackslash}p{(\linewidth - 4\tabcolsep) * \real{0.2381}}@{}}
\toprule\noalign{}
\begin{minipage}[b]{\linewidth}\raggedright
Possible Cause
\end{minipage} & \begin{minipage}[b]{\linewidth}\raggedright
Diagnostic Test
\end{minipage} & \begin{minipage}[b]{\linewidth}\raggedright
Solution
\end{minipage} \\
\midrule\noalign{}
\endhead
\bottomrule\noalign{}
\endlastfoot
Ball-screw wear & Visual inspection, measure backlash & Replace screw or
nut \\
Bearing failure & Rotation test, listen for grinding & Replace failed
bearing immediately \\
Resonance & Frequency analysis & Add damping or structural
modification \\
High servo gains & Reduce gains, observe noise change & Optimize servo
tuning \\
\end{longtable}
}

\begin{center}\rule{0.5\linewidth}{0.5pt}\end{center}

\subsubsection{6.4 Performance
Benchmarking}\label{performance-benchmarking}

Establish baseline performance metrics during commissioning and compare
periodically to detect degradation:

{\def\LTcaptype{} % do not increment counter
\begin{longtable}[]{@{}
  >{\raggedright\arraybackslash}p{(\linewidth - 10\tabcolsep) * \real{0.1294}}
  >{\raggedright\arraybackslash}p{(\linewidth - 10\tabcolsep) * \real{0.1882}}
  >{\raggedright\arraybackslash}p{(\linewidth - 10\tabcolsep) * \real{0.1882}}
  >{\raggedright\arraybackslash}p{(\linewidth - 10\tabcolsep) * \real{0.1882}}
  >{\raggedright\arraybackslash}p{(\linewidth - 10\tabcolsep) * \real{0.0824}}
  >{\raggedright\arraybackslash}p{(\linewidth - 10\tabcolsep) * \real{0.2235}}@{}}
\toprule\noalign{}
\begin{minipage}[b]{\linewidth}\raggedright
Parameter
\end{minipage} & \begin{minipage}[b]{\linewidth}\raggedright
Baseline (New)
\end{minipage} & \begin{minipage}[b]{\linewidth}\raggedright
After 1000 hrs
\end{minipage} & \begin{minipage}[b]{\linewidth}\raggedright
After 2000 hrs
\end{minipage} & \begin{minipage}[b]{\linewidth}\raggedright
Limit
\end{minipage} & \begin{minipage}[b]{\linewidth}\raggedright
Action if Exceeded
\end{minipage} \\
\midrule\noalign{}
\endhead
\bottomrule\noalign{}
\endlastfoot
Deflection @ 400N (mm) & 0.008 & 0.009 & 0.010 & 0.020 & Inspect
structure \\
Rail parallelism (mm) & 0.005 & 0.008 & 0.012 & 0.030 & Shim rails \\
Backlash (mm) & 0.010 & 0.015 & 0.022 & 0.030 & Increase preload \\
Natural frequency (Hz) & 185 & 182 & 178 & 150 & Add stiffening \\
Counterbalance (\% imbalance) & 4\% & 7\% & 11\% & 15\% & Adjust
force \\
Positioning accuracy (mm) & +/-0.02 & +/-0.03 & +/-0.04 & +/-0.05 &
Calibrate or repair \\
\end{longtable}
}

\textbf{Trend Analysis:}

Plot key parameters vs.~operating hours. Linear degradation is normal
(wear). Sudden changes indicate failure or misalignment requiring
immediate investigation.

\begin{center}\rule{0.5\linewidth}{0.5pt}\end{center}

\subsubsection{6.5 Safety Considerations}\label{safety-considerations-1}

\textbf{Critical Safety Features:}

\begin{enumerate}
\def\labelenumi{\arabic{enumi}.}
\tightlist
\item
  \textbf{Motor Brake:} Must engage on power loss to prevent carriage
  free-fall

  \begin{itemize}
  \tightlist
  \item
    Test monthly: Disable power while carriage elevated
  \item
    Brake should hold position with \textless{} 1 mm drop
  \end{itemize}
\item
  \textbf{Counterbalance:} Prevents rapid descent if brake fails

  \begin{itemize}
  \tightlist
  \item
    Should slow (not stop) carriage fall to \textless{} 50 mm/s
  \end{itemize}
\item
  \textbf{Software Limits:} Prevent over-travel beyond mechanical range

  \begin{itemize}
  \tightlist
  \item
    Set soft limits 5-10 mm inside hard limits
  \item
    Test: Jog toward limit, verify deceleration profile
  \end{itemize}
\item
  \textbf{Emergency Stop:} Immediately stops all motion

  \begin{itemize}
  \tightlist
  \item
    Test weekly: Press E-stop during motion
  \item
    Verify motion halts within 10 mm
  \end{itemize}
\item
  \textbf{Guarding:} Prevents hand/finger pinch points

  \begin{itemize}
  \tightlist
  \item
    Install bellows or covers over moving carriage
  \item
    Interlock access doors to disable motion when open
  \end{itemize}
\end{enumerate}

\textbf{Lockout/Tagout Procedure:}

Before performing maintenance: 1. Press emergency stop button 2.
Disconnect motor power at amplifier or main disconnect 3. Manually lower
carriage to bottom position (safe height) 4. Install mechanical blocks
or pins to prevent carriage motion 5. Attach lockout tag with technician
name and date

\begin{center}\rule{0.5\linewidth}{0.5pt}\end{center}

\subsection{6.6 Conclusion: Verification as Ongoing
Process}\label{conclusion-verification-as-ongoing-process}

Vertical axis verification is not a one-time commissioning activity--it
is an ongoing process of measurement, analysis, and adjustment. Regular
verification testing:

\begin{itemize}
\tightlist
\item
  Detects performance degradation before it causes part defects
\item
  Identifies wear patterns guiding proactive component replacement
\item
  Documents machine health over its lifecycle
\item
  Ensures continued compliance with specification
\item
  Maximizes uptime and productivity
\end{itemize}

\textbf{Best Practice:} Treat verification data as a predictive
maintenance tool. Trending analysis reveals when components are
approaching end-of-life, allowing scheduled replacement during planned
downtime rather than reactive failure responses.

The comprehensive checklist and procedures provided in this section form
the foundation of a world-class precision machine maintenance program.

\begin{center}\rule{0.5\linewidth}{0.5pt}\end{center}

\begin{center}\rule{0.5\linewidth}{0.5pt}\end{center}

\subsection{References}\label{references-4}

\begin{enumerate}
\def\labelenumi{\arabic{enumi}.}
\tightlist
\item
  \textbf{ISO 230-2:2014} - Test code for machine tools - Determination
  of accuracy and repeatability
\item
  \textbf{ISO 13849-1:2015} - Safety of machinery - Safety-related parts
  of control systems
\item
  \textbf{Renishaw Ballbar QC20-W User Guide} - Circular interpolation
  testing
\item
  \textbf{Heidenhain Linear Encoders Catalog} - Measurement system
  specifications
\item
  \textbf{ASME B89.3.4-2010} - Axes of Rotation: Methods for Specifying
  and Testing
\item
  \textbf{Machine Tool Maintenance Best Practices} - SME Technical Paper
  Series
\end{enumerate}

\newpage

\section{Module 2 - Vertical Axis \&
Z-Stage}\label{module-2---vertical-axis-z-stage-4}

\subsection{Overview}\label{overview-3}

Thermal expansion in vertical columns directly affects Z-axis
positioning accuracy. Unlike horizontal axes where thermal growth
affects positioning but not part dimensions, vertical thermal expansion
changes tool height relative to the workpiece. This section addresses
thermal analysis, design strategies, and compensation techniques for
thermally stable vertical axes.

\subsection{Thermal Expansion
Fundamentals}\label{thermal-expansion-fundamentals}

\subsubsection{Linear Expansion}\label{linear-expansion}

{[}\Delta L = \alpha \times L \times \Delta T{]}

Where: - alpha = coefficient of thermal expansion (/°C) - L = original
length (mm) - DeltaT = temperature change (°C)

\textbf{Material comparison:}

{\def\LTcaptype{} % do not increment counter
\begin{longtable}[]{@{}lll@{}}
\toprule\noalign{}
Material & alpha (/°C) & 1m length, 10°C rise \\
\midrule\noalign{}
\endhead
\bottomrule\noalign{}
\endlastfoot
\textbf{Steel} & 12 x 10\textsuperscript{-}6 & 120 µm \\
\textbf{Aluminum} & 23 x 10\textsuperscript{-}6 & 230 µm \\
\textbf{Cast iron} & 11 x 10\textsuperscript{-}6 & 110 µm \\
\textbf{Invar} & 1.2 x 10\textsuperscript{-}6 & 12 µm \\
\end{longtable}
}

\textbf{Example:}

Steel column: 800mm height, 5°C temperature rise {[}\Delta L = 12
\times 10\^{}\{-6\} \times 800 \times 5 = 48 \text{ µm}{]}

\textbf{Impact:} 48 µm error in Z-position (significant for precision
work).

\subsection{Heat Sources}\label{heat-sources}

\subsubsection{Internal Heat Sources}\label{internal-heat-sources}

\textbf{Spindle motor:} - Largest heat source (50-90\% of total) - Heat
rises by convection - Warms upper portion of column

\textbf{Ball screw friction:} - Continuous operation generates heat -
Distributed along screw length - Typically 5-15\% of total heat

\textbf{Linear guide friction:} - Minor contribution (\textless5\%) -
Distributed vertically

\textbf{Electronics:} - Drives and power supplies - Usually mounted away
from column - Ambient temperature contribution

\subsubsection{External Heat Sources}\label{external-heat-sources}

\textbf{Shop environment:} - Day/night temperature cycles (+/-5-10°C
typical) - Seasonal variations - HVAC system effects

\textbf{Solar radiation:} - Windows near machine - Can create
temperature gradients

\textbf{Coolant:} - Cutting fluid temperature variations - Coolant flow
affects local temperatures

\subsection{Thermal Design Strategies}\label{thermal-design-strategies}

\subsubsection{Minimize Heat Generation}\label{minimize-heat-generation}

\textbf{Efficient spindle:} - High-quality bearings (low friction) -
Adequate cooling - Proper preload (not over-tightened)

\textbf{Low-friction guides:} - Proper lubrication - Correct preload
selection - Regular maintenance

\textbf{Counterbalancing:} - Reduces motor RMS torque - Lower motor
heating - Less thermal drift

\subsubsection{Thermal Isolation}\label{thermal-isolation}

\textbf{Isolate spindle heat:} - Air gap between spindle and mount -
Thermal break materials (G10, ceramic spacers) - Active cooling of
spindle mount

\textbf{Shield from environment:} - Column covers/enclosures -
Insulation on exterior surfaces - Minimize exposed surface area

\subsubsection{Thermal Mass}\label{thermal-mass}

\textbf{Large thermal mass benefits:} - Resists rapid temperature
changes - Averages out short-term variations - Longer time constant for
stabilization

\textbf{Design approach:} - Solid column (not minimum-weight) - Cast
iron better than steel or aluminum for thermal stability - Internal mass
(fill cavities with sand or epoxy)

\subsubsection{Symmetric Design}\label{symmetric-design}

\textbf{Symmetric thermal expansion:} - Heat sources centered on column
- Symmetric cross-section - Equal expansion both sides minimizes
deflection

\textbf{Example:}

\begin{verbatim}
Poor:                    Better:
Spindle offset           Spindle centered
    |                        |
  /   \                   /     \
 |     |                 |   |   |
 |motor|                 | motor |
\end{verbatim}

Offset spindle causes column to bow; centered spindle expands
symmetrically.

\subsection{Active Thermal Management}\label{active-thermal-management}

\subsubsection{Spindle Cooling}\label{spindle-cooling}

\textbf{Water-cooled spindles:} - Circulating coolant removes heat -
Maintains constant spindle temperature - Prevents heat transfer to
column

\textbf{Coolant system requirements:} - Temperature-controlled chiller
(+/-0.5°C) - Adequate flow rate (2-4 L/min typical) - Low coolant
temperature (15-20°C)

\textbf{Air-cooled spindles:} - Less effective heat removal - Fan
directs hot air away from column - Acceptable for light-duty or
low-power applications

\subsubsection{Column Temperature
Control}\label{column-temperature-control}

\textbf{Coolant circulation through column:} - Tubes cast or welded
inside column - Circulating coolant maintains uniform temperature - Used
in high-precision machines

\textbf{Implementation:} - Copper tubing (8-12mm diameter) in serpentine
pattern - Flow rate: 1-2 L/min - Same chiller as spindle (temperature
stability critical)

\textbf{Cost vs.~benefit:} - High cost (complex fabrication) - Excellent
thermal stability (\textless5 µm drift) - Justified for precision
machining

\subsubsection{Environmental Control}\label{environmental-control}

\textbf{Air-conditioned machine room:} - +/-1-2°C temperature stability
- Eliminates day/night cycles - Expensive but most effective

\textbf{Machine enclosure:} - Isolates from shop environment - Internal
temperature control - Collects coolant and chips

\subsection{Thermal Compensation}\label{thermal-compensation}

\subsubsection{Software Compensation}\label{software-compensation}

\textbf{Temperature sensors:} - RTD or thermocouple on column - Multiple
sensors for gradient measurement - High accuracy (+/-0.1°C) required

\textbf{Compensation algorithm:}

{[}Z\_\{corrected\} = Z\_\{programmed\} + \alpha \times L
\times (T\_\{current\} - T\_\{reference\}){]}

\textbf{Implementation:} - LinuxCNC: HAL component reads sensor, applies
offset - Industrial controls: Built-in compensation features - Manual:
G-code offset based on measured temperature

\textbf{Accuracy:} - Simple linear compensation: +/-10-20 µm typical -
Multi-point compensation: +/-5 µm achievable - Requires calibration and
validation

\subsubsection{Mechanical Compensation}\label{mechanical-compensation}

\textbf{Bi-material compensator:} - Two materials with different alpha
expand at different rates - Designed to cancel column expansion -
Passive, no sensors or software required

\textbf{Invar-steel combination:} - Invar rod inside steel column -
Opposite expansion rates - Can achieve near-zero net expansion

\textbf{Complexity:} - Difficult to design and fabricate - Limited
adjustability - Rare in modern machines (software compensation
preferred)

\subsection{Thermal Stabilization}\label{thermal-stabilization}

\subsubsection{Warm-Up Procedure}\label{warm-up-procedure}

\textbf{New machines require thermal stabilization:}

\begin{enumerate}
\def\labelenumi{\arabic{enumi}.}
\tightlist
\item
  \textbf{Power on all systems} (spindle, motors, electronics)
\item
  \textbf{Run warm-up routine} (automated cycle)

  \begin{itemize}
  \tightlist
  \item
    Spindle at 50\% rated speed for 20-30 minutes
  \item
    Z-axis jogging over full travel
  \item
    All axes exercised
  \end{itemize}
\item
  \textbf{Monitor temperature} until stable (typically 30-60 minutes)
\item
  \textbf{Set tool offsets} after warm-up
\item
  \textbf{Begin production}
\end{enumerate}

\textbf{Daily procedure:} - Allow 15-30 minute warm-up before precision
work - Longer for machines in non-temperature-controlled environments

\subsubsection{Thermal Stability
Testing}\label{thermal-stability-testing}

\textbf{Procedure:} 1. Set up dial indicator on table, reading against
spindle face 2. Zero indicator when machine at stable temperature 3. Run
spindle at typical operating conditions 4. Monitor indicator reading
over time (1-2 hours) 5. Record thermal drift (µm/hour)

\textbf{Acceptable performance:} - Hobby/light-duty: \textless50 µm/hour
- General purpose: \textless20 µm/hour - Precision: \textless10 µm/hour
- Ultra-precision: \textless5 µm/hour

\subsection{Practical Design Examples}\label{practical-design-examples}

\subsubsection{Small Mill (Minimal Thermal
Control)}\label{small-mill-minimal-thermal-control}

\textbf{Design features:} - Aluminum column (accepts thermal expansion)
- Air-cooled spindle (500W) - No active thermal management - Shop
environment (+/-5°C)

\textbf{Expected drift:} +/-60 µm over temperature range

\textbf{Mitigation:} - Warm-up procedure before precision work - Part
dimension tolerance \textgreater=100 µm - Software compensation if
needed

\subsubsection{Medium Mill (Moderate Thermal
Control)}\label{medium-mill-moderate-thermal-control}

\textbf{Design features:} - Steel column - Water-cooled spindle (2.2kW)
with chiller - Temperature sensor for software compensation - Enclosed
machine - Shop environment (+/-3°C)

\textbf{Expected drift:} +/-15 µm with compensation

\textbf{Cost:} \textasciitilde(500 added (chiller, sensor, enclosure)

\subsubsection{Precision Mill (Advanced Thermal
Control)}\label{precision-mill-advanced-thermal-control}

\textbf{Design features:} - Cast iron column (thermal mass + low alpha)
- Water-cooled spindle with precision chiller (+/-0.5°C) - Coolant
circulation through column - Temperature-controlled enclosure (+/-1°C) -
Multi-sensor compensation

\textbf{Expected drift:} \textless5 µm

\textbf{Cost:} )5,000+ added (thermal management systems)

\subsection{Key Takeaways}\label{key-takeaways-3}

\begin{enumerate}
\def\labelenumi{\arabic{enumi}.}
\tightlist
\item
  \textbf{Vertical thermal expansion} directly affects Z-axis accuracy
  (unlike horizontal axes)
\item
  \textbf{Steel column, 1m height, 5°C rise} = 60 µm expansion
  (significant)
\item
  \textbf{Spindle heat} is largest contributor (50-90\% of thermal load)
\item
  \textbf{Water-cooled spindles} essential for thermal stability in
  precision applications
\item
  \textbf{Thermal mass} (cast iron, thick sections) resists rapid
  temperature changes
\item
  \textbf{Software compensation} cost-effective solution for moderate
  precision
\item
  \textbf{Active cooling} (column coolant circulation) for ultimate
  stability
\item
  \textbf{Warm-up procedures} mandatory before precision machining
\item
  \textbf{Environmental control} (air-conditioned room) most effective
  but expensive
\item
  \textbf{Design for application:} Light-duty accepts drift; precision
  requires active management
\end{enumerate}

\begin{center}\rule{0.5\linewidth}{0.5pt}\end{center}

\textbf{Next}: \href{section-2.8-spindle-mounting.md}{Section 2.8 -
Spindle Mounting}

\textbf{Previous}: \href{section-2.6-motor-drive-sizing.md}{Section 2.6
- Motor and Drive Sizing}

\newpage

\section{Module 2 - Vertical Axis \&
Z-Stage}\label{module-2---vertical-axis-z-stage-5}

\subsection{Overview}\label{overview-4}

The spindle mounting interface is critical for runout control, thermal
management, and structural rigidity. This section covers spindle
selection criteria for vertical mounting, interface design, cooling
integration, and toolholder standards.

\subsection{Spindle Selection for
Z-Axis}\label{spindle-selection-for-z-axis}

\subsubsection{Weight Considerations}\label{weight-considerations}

Spindle weight directly impacts: - Moving mass (affects motor sizing,
acceleration) - Column deflection under gravity - Counterbalance
requirements - Linear guide and ball screw loading

\textbf{Weight ranges by spindle type:} - ER-collet air-cooled
(400-800W): 2-4 kg - ER-collet water-cooled (1.5-2.2kW): 4-7 kg -
BT/CAT40 ATC spindle: 15-25 kg - HSK-A63 ATC spindle: 12-18 kg

\textbf{Design guideline:} Minimize spindle weight while meeting power
requirements.

\subsubsection{Cooling Requirements}\label{cooling-requirements}

\textbf{Air-cooled spindles:} - Suitable for light-duty (\textless1kW) -
Fan exhausts hot air (direct away from column) - Less thermal stability
- Lower cost, simpler installation

\textbf{Water-cooled spindles:} - Required for \textgreater1kW
continuous operation - Stable operating temperature (+/-1-2°C) - Chiller
required (adds (300-)1500 cost) - Better thermal performance

\textbf{Chiller specifications:} - Cooling capacity: 1.5x spindle power
minimum - Temperature control: +/-0.5°C for precision work - Flow rate:
2-4 L/min typical

\subsection{Mounting Interface Design}\label{mounting-interface-design}

\subsubsection{Spindle Mount Plate}\label{spindle-mount-plate}

\textbf{Requirements:} - Flatness: 10 µm (perpendicular to spindle axis)
- Rigidity: Minimal deflection under cutting forces - Thermal isolation:
Prevent heat transfer to carriage - Tool clearance: Access for tool
changes

\textbf{Material selection:} - Aluminum 7075-T6 (high
strength-to-weight) - Machined from solid (no welding for precision) -
Ribbed structure for stiffness

\textbf{Design example:}

Mount plate for 2.2kW water-cooled spindle: - Material: Al 7075-T6 -
Thickness: 20mm (ribbed to 15mm in non-critical areas) - Spindle
mounting: 4x M6 bolts, 80mm bolt circle - Carriage mounting: 6x M8 bolts
- Weight: 0.65 kg (optimized)

\subsubsection{Runout Control}\label{runout-control}

\textbf{Spindle runout tolerance:} - General milling: \textless10 µm TIR
at tool nose - Precision work: \textless5 µm TIR - Ultra-precision:
\textless2 µm TIR

\textbf{Mounting concentricity:} - Spindle bore to mounting face:
\textless5 µm (specified by manufacturer) - Mount plate to carriage:
\textless10 µm (installation tolerance) - Total stack-up: \textless15 µm

\textbf{Alignment procedure:} 1. Mount spindle to plate with indicator
in spindle taper 2. Indicate spindle nose runout while rotating by hand
3. Adjust shims if necessary (some spindles have adjustable mounting) 4.
Torque mounting bolts progressively to specification 5. Re-check runout
after torquing

\subsubsection{Thermal Isolation}\label{thermal-isolation-1}

\textbf{Spindle heat transfer to carriage:}

Heat flux: {[}Q = k \times A \times \frac{\Delta T}{d}{]}

Where: - k = thermal conductivity - A = contact area - DeltaT =
temperature difference - d = conduction path length

\textbf{Isolation strategies:} 1. Minimal contact area (raised bosses)
2. Thermal break material (G10 spacers, 0.5-1mm thick) 3. Air gap
between spindle and mount (5-10mm where possible) 4. Direct cooling of
mount plate (coolant passages)

\subsection{Toolholder Interfaces}\label{toolholder-interfaces}

\subsubsection{ER Collet System}\label{er-collet-system}

\textbf{Advantages:} - Lightweight (2-4 kg including spindle) - Wide
tool diameter range (ER20: 1-13mm, ER25: 1-16mm) - Low cost - Good for
small machines

\textbf{Disadvantages:} - Manual tool changes only - Moderate runout
(5-10 µm typical) - Lower rigidity than taper systems

\textbf{Applications:} Small DIY mills, hobby machines, router
conversions

\subsubsection{BT/CAT Taper (ISO 7388/ANSI
B5.50)}\label{btcat-taper-iso-7388ansi-b5.50}

\textbf{Common sizes:} - BT30/CAT30: Light-duty, smaller machines -
BT40/CAT40: Standard industrial size - BT50/CAT50: Heavy-duty
applications

\textbf{Advantages:} - Automatic tool change (ATC) compatible - High
rigidity (dual-contact: taper + face) - Widely available toolholders -
Industrial standard

\textbf{Disadvantages:} - Heavier spindles (15-25 kg for BT40) - More
expensive - Larger envelope

\textbf{Applications:} Production mills, machining centers

\subsubsection{HSK (Hollow Shank Taper)}\label{hsk-hollow-shank-taper}

\textbf{Modern high-speed interface:} - HSK-A63: General purpose
(comparable to BT40) - HSK-A100: Heavy-duty

\textbf{Advantages:} - Better high-speed performance (balanced,
symmetric) - Higher stiffness than BT/CAT - Shorter tool change time -
Face contact only (no taper friction at high speed)

\textbf{Disadvantages:} - Most expensive - Less toolholder availability
than BT/CAT - Heavier than ER systems

\textbf{Applications:} High-speed machining, modern industrial machines

\subsection{Cooling Integration}\label{cooling-integration}

\subsubsection{Spindle Coolant
Connections}\label{spindle-coolant-connections}

\textbf{Water-cooled spindle plumbing:} - Quick-disconnect fittings
(allow spindle removal) - Flexible hose (accommodates Z-axis motion) -
Routing through cable carrier - Drip loop at lowest point (prevents
water siphoning on shutdown)

\textbf{Flow direction:} - Inlet at bottom, outlet at top (natural
convection assists) - Or follow manufacturer specification

\subsubsection{Through-Spindle Coolant
(TSC)}\label{through-spindle-coolant-tsc}

\textbf{High-pressure coolant delivery:} - Coolant fed through spindle
center, out tool - Pressure: 300-1000 PSI (20-70 bar) - Requires rotary
union on spindle - Special toolholders with coolant passage

\textbf{Benefits:} - Direct cooling at cutting edge - Excellent chip
evacuation - Enables deep hole drilling

\textbf{Implementation challenges:} - Adds weight (rotary union: 0.5-1
kg) - Requires high-pressure pump - More complex plumbing - Higher cost
(\$500-2000 additional)

\subsection{Vibration and Damping}\label{vibration-and-damping}

\subsubsection{Spindle-Induced
Vibration}\label{spindle-induced-vibration}

\textbf{Sources:} - Unbalanced tool or tool holder - Bearing wear -
Chatter during cutting

\textbf{Effects on vertical axis:} - Accelerates linear guide wear -
Reduces surface finish - Can excite column resonances

\textbf{Mitigation:} - Balance all toolholders (on-machine or external
balancer) - Quality spindle with low vibration (\textless0.5 mm/s\^{}2
typical) - Rigid mounting (no compliance in spindle mount) - Adequate
column stiffness

\subsubsection{Damping Strategies}\label{damping-strategies}

\textbf{Passive damping:} - Viscoelastic damping layers in mount
structure - Tuned mass dampers (for specific resonances) - Cast iron
components (inherent damping)

\textbf{Active damping:} - Rarely used except in ultra-precision
machines - Piezo-actuated damping systems - Requires sensors and control
system

\subsection{Key Takeaways}\label{key-takeaways-4}

\begin{enumerate}
\def\labelenumi{\arabic{enumi}.}
\tightlist
\item
  \textbf{Spindle weight} is critical design parameter for
  Z-axis--minimize while meeting power needs
\item
  \textbf{Water-cooling} essential for \textgreater1kW spindles and
  thermal stability
\item
  \textbf{Runout control} requires precise mounting: target \textless10
  µm TIR at tool nose
\item
  \textbf{Thermal isolation} prevents heat transfer from spindle to
  column
\item
  \textbf{ER collets} for lightweight applications, BT/CAT for
  industrial use, HSK for high-speed
\item
  \textbf{Mount plate stiffness} important for rigidity under cutting
  forces
\item
  \textbf{Tool change} considerations: Manual (ER) vs.~automatic
  (BT/CAT/HSK)
\item
  \textbf{Through-spindle coolant} beneficial for production but adds
  complexity and weight
\item
  \textbf{Vibration control} through balancing and rigid mounting
\item
  \textbf{Integration planning}: Coolant routing, electrical
  connections, tool clearance
\end{enumerate}

\begin{center}\rule{0.5\linewidth}{0.5pt}\end{center}

\textbf{Next}: \href{section-2.9-cable-management.md}{Section 2.9 -
Cable Management}

\textbf{Previous}: \href{section-2.7-thermal-management.md}{Section 2.7
- Thermal Management}

\newpage

\section{Module 2 - Vertical Axis \&
Z-Stage}\label{module-2---vertical-axis-z-stage-6}

\subsection{Overview}\label{overview-5}

Vertical axis cable management presents unique challenges: cables must
travel with the moving carriage while avoiding snagging, minimizing
strain, and resisting gravity's tendency to pull cables downward. This
section covers cable carrier selection, routing strategies, and strain
relief techniques specific to Z-axis applications.

\subsection{Cable Types and
Requirements}\label{cable-types-and-requirements}

\subsubsection{Power Cables}\label{power-cables}

\textbf{Motor power:} - 3-phase for servo motors (4-wire plus ground) -
AWG 18-14 typical for NEMA 17-34 - Continuous flex rated (millions of
cycles)

\textbf{Spindle power:} - AWG 12-8 for 1.5-7.5kW spindles - VFD output
(high dV/dt, requires shielding) - Separate from signal cables

\subsubsection{Signal Cables}\label{signal-cables}

\textbf{Encoder feedback:} - Differential signals (RS-422 typical) -
Twisted-pair, shielded - Low capacitance (\textless100 pF/m) - Keep
\textless3m total length

\textbf{Limit switches / Home sensors:} - 22-24 AWG - Shielded to
prevent EMI

\textbf{Temperature sensors:} - Thermocouple or RTD extension wire -
Shielded, keep away from power cables

\subsubsection{Fluid Lines}\label{fluid-lines}

\textbf{Spindle coolant:} - 6-10mm ID flexible hose - Rated for coolant
temperature and pressure - Quick-disconnects for serviceability

\textbf{Lubrication lines:} - 4-6mm tubing for automatic lubrication
systems - Low-pressure (\textless10 bar)

\subsection{Cable Carrier Systems}\label{cable-carrier-systems}

\subsubsection{Carrier Types}\label{carrier-types}

\textbf{Plastic drag chain:} - Most common for small-medium machines -
Low cost ((2-5/ft) - Quiet operation - Adequate for most applications

\textbf{Steel cable carrier:} - Heavy-duty applications - Higher cost
()10-20/ft) - Better for long travels (\textgreater2m) - Oil/coolant
resistant

\textbf{Spiral wrap / Conduit:} - Lightweight, flexible - Not suitable
for Z-axis (gravity pulls loops) - Use only for stationary routing

\subsubsection{Carrier Sizing}\label{carrier-sizing}

\textbf{Internal dimensions:}

{[}A\_\{required\} = \frac{\sum A_{cables}}{0.4}{]}

Fill ratio: 40\% maximum (allows cable movement, prevents binding)

\textbf{Example:} - 4x AWG 18 power: 4x 8 mm\^{}2 = 32 mm\^{}2 - 2x
shielded twisted pair: 2x 12 mm\^{}2 = 24 mm\^{}2 - 1x coolant hose: 80
mm\^{}2 (10mm OD) - Total: 136 mm\^{}2 - Required carrier area: 136 /
0.4 = 340 mm\^{}2

Select 25x25mm carrier (internal area = 625 mm\^{}2, adequate margin)

\subsubsection{Bend Radius}\label{bend-radius}

\textbf{Minimum bend radius:}

{[}R\_\{min\} = 10 \times D\_\{largest\textbackslash\_cable\}{]}

Conservative guideline ensures long cable life.

\textbf{Example:} - Largest cable: 10mm OD coolant hose - Minimum
radius: 100mm - Selected carrier: 150mm radius (safety margin)

\subsection{Mounting Configurations}\label{mounting-configurations}

\subsubsection{Top-Mounted Carrier
(Recommended)}\label{top-mounted-carrier-recommended}

\begin{verbatim}
    Column
        |
    [Fixed end]
        |
      Cable
      Carrier <- Moves with carriage
        |
    [Moving end]
        |
    Carriage
\end{verbatim}

\textbf{Advantages:} - Gravity assists cable return (reduces wear) -
Clean appearance - Protected from chips/coolant

\textbf{Disadvantages:} - Requires space above column - More complex
mounting at top

\subsubsection{Side-Mounted Carrier}\label{side-mounted-carrier}

\textbf{Carrier mounted to side of column:} - Works when top mounting
not feasible - Requires rigid carrier support bracket - Gravity pulls
carrier outward (support every 300-500mm)

\subsubsection{Internal Carrier
(Advanced)}\label{internal-carrier-advanced}

\textbf{Carrier routed inside hollow column:} - Clean external
appearance - Maximum protection - Difficult to service - Requires larger
column cross-section

\subsection{Installation Best
Practices}\label{installation-best-practices}

\subsubsection{Fixed End Mounting}\label{fixed-end-mounting}

\textbf{Anchor point (top of column or machine frame):} - Rigid mounting
(no deflection) - Strain relief for all cables exiting carrier - Cable
ties every 100mm - Drip loop for coolant lines (lowest point before
entry)

\subsubsection{Moving End Mounting}\label{moving-end-mounting}

\textbf{Carriage attachment:} - Secure mounting to carriage plate -
Flexible boot or conduit from carrier to components - Allow 50-100mm
extra cable length (service loop) - Strain relief clamps at carrier exit

\subsubsection{Cable Routing Inside
Carrier}\label{cable-routing-inside-carrier}

\textbf{Layout strategy:} 1. Power cables on one side 2. Signal cables
on opposite side (EMI separation) 3. Coolant hoses in center or separate
chamber 4. Secure cables every 200-300mm inside carrier (prevent
tangling)

\textbf{Avoid:} - Mixing power and signal in same bundle - Excessive
tension on cables - Sharp bends at carrier entry/exit - Cables crossing
over each other

\subsection{Strain Relief}\label{strain-relief}

\subsubsection{At Carrier Exits}\label{at-carrier-exits}

\textbf{Fixed end:} - 90° conduit fitting or cable gland - Cable ties to
support bracket - Maintain bend radius - Allow thermal expansion (don't
over-tighten)

\textbf{Moving end:} - Flexible conduit (last 150-200mm) - Spiral wrap
for small cables - Service loop (coiled excess cable) - Secure to
carriage but allow slight movement

\subsubsection{At Component Connections}\label{at-component-connections}

\textbf{Motor/spindle:} - Connector with strain relief collar - Cable
clamped 100mm before connector - No tension on connector pins - Avoid
sharp bends near connector

\textbf{Drives/electronics:} - Cable entry through grommet or gland -
Secured to enclosure wall - Drip loop before entry (if moisture
possible)

\subsection{Vertical-Specific
Considerations}\label{vertical-specific-considerations}

\subsubsection{Gravity Effects}\label{gravity-effects}

\textbf{Cable weight pulls downward:} - Use lightweight cables when
possible - Support carrier at intervals (side-mount) - Top-mount
preferred (gravity assists)

\textbf{Coolant draining:} - Route coolant hoses with continuous
downward slope - Avoid traps where fluid can pool - Drip loop at lowest
point

\subsubsection{Travel Length
Calculation}\label{travel-length-calculation}

\textbf{Carrier length:}

{[}L\_\{carrier\} = L\_\{travel\} + 2 \times R\_\{bend\} + L\_\{fixed\}
+ L\_\{moving\}{]}

Where: - (L\_\{travel\}) = Z-axis travel - (R\_\{bend\}) = carrier bend
radius - (L\_\{fixed\}) = fixed end length (typically 150-300mm) -
(L\_\{moving\}) = moving end length (typically 200-400mm)

\textbf{Example:} - Z-travel: 400mm - Bend radius: 150mm - Fixed end:
200mm - Moving end: 300mm - Total carrier length: 400 + 2(150) + 200 +
300 = 1200mm

Order 1500mm carrier (allows adjustment)

\subsubsection{Dynamic Loading}\label{dynamic-loading}

\textbf{Acceleration forces:} - Cables experience inertial loading
during acceleration - Secure cables inside carrier to prevent shifting -
Use cable separators if high-acceleration application

\subsection{Troubleshooting Common
Issues}\label{troubleshooting-common-issues-1}

\subsubsection{Cable Wear}\label{cable-wear}

\textbf{Symptoms:} Frayed insulation, broken conductors

\textbf{Causes:} - Bend radius too tight - Cables rubbing against
carrier edges - Over-filled carrier (cables binding)

\textbf{Solutions:} - Increase bend radius - Add cable separators -
Route cables to avoid edges - Reduce fill ratio

\subsubsection{Signal Interference}\label{signal-interference}

\textbf{Symptoms:} Encoder errors, erratic motion, noise on sensors

\textbf{Causes:} - Power and signal cables not separated - Inadequate
shielding - Ground loops

\textbf{Solutions:} - Separate power and signal (opposite sides of
carrier) - Use shielded cables, ground shields at one end only - Twisted
pair for differential signals - Ferrite beads on signal cables near
drives

\subsubsection{Carrier Binding}\label{carrier-binding}

\textbf{Symptoms:} Jerky motion, increased motor current, audible
clicking

\textbf{Causes:} - Carrier over-filled - Misaligned carrier mounting -
Insufficient lubrication

\textbf{Solutions:} - Reduce cable count or increase carrier size -
Realign carrier mounting (parallel to motion axis) - Lubricate carrier
links (plastic-compatible grease)

\subsubsection{Coolant Leaks}\label{coolant-leaks}

\textbf{Symptoms:} Dripping coolant, pooling in carrier

\textbf{Causes:} - Inadequate hose strain relief - Hose kinked or
damaged - Quick-disconnect leaking

\textbf{Solutions:} - Add hose support near connections - Replace
damaged hose - Use higher-quality quick-disconnects with seals

\subsection{Key Takeaways}\label{key-takeaways-5}

\begin{enumerate}
\def\labelenumi{\arabic{enumi}.}
\tightlist
\item
  \textbf{Top-mounted carriers} preferred for Z-axis (gravity assists
  cable return)
\item
  \textbf{Fill ratio} must be \textless40\% to prevent binding and allow
  cable movement
\item
  \textbf{Bend radius} minimum 10x largest cable diameter for long life
\item
  \textbf{Separate power and signal} cables to prevent EMI (opposite
  sides of carrier)
\item
  \textbf{Strain relief} critical at both fixed and moving ends
\item
  \textbf{Service loops} provide extra cable length for maintenance and
  adjustment
\item
  \textbf{Drip loops} required for coolant hoses before entering
  components
\item
  \textbf{Carrier length} calculation includes travel + bends +
  fixed/moving ends
\item
  \textbf{Cable ties} every 100-200mm prevent tangling inside carrier
\item
  \textbf{Continuous downward slope} for fluid lines prevents trapping
  and leaks
\end{enumerate}

\begin{center}\rule{0.5\linewidth}{0.5pt}\end{center}

\textbf{Next}: \href{section-2.10-assembly-alignment.md}{Section 2.10 -
Assembly and Alignment}

\textbf{Previous}: \href{section-2.8-spindle-mounting.md}{Section 2.8 -
Spindle Mounting}

\end{document}
